\documentclass[12pt,a4paper]{report}
\usepackage[utf8]{inputenc}
\usepackage[ngerman]{babel}
\usepackage{amsmath}
\usepackage{amssymb}
\usepackage{amsthm}
\usepackage{graphicx}
\usepackage{listings}
\usepackage{xcolor}
\usepackage[hidelinks]{hyperref}
\usepackage{microtype}
\usepackage{geometry}
\usepackage{fancyhdr}
\usepackage{tikz}
\usepackage{float}
\usepackage{rotating}
\usepackage{enumitem}
\setcounter{MaxMatrixCols}{15}

\setlist[itemize,1]{label=--}

\geometry{a4paper, margin=2.5cm}

% Listings Konfiguration
\definecolor{codegreen}{rgb}{0,0.6,0}
\definecolor{codegray}{rgb}{0.5,0.5,0.5}
\definecolor{codepurple}{rgb}{0.58,0,0.82}
\definecolor{backcolour}{rgb}{0.95,0.95,0.92}

\lstdefinestyle{pythonstyle}{
    backgroundcolor=\color{backcolour},   
    commentstyle=\color{codegreen},
    keywordstyle=\color{magenta},
    numberstyle=\tiny\color{codegray},
    stringstyle=\color{codepurple},
    basicstyle=\ttfamily\footnotesize,
    breakatwhitespace=false,         
    breaklines=true,                 
    captionpos=b,                    
    keepspaces=true,                 
    numbers=left,                    
    numbersep=5pt,                  
    showspaces=false,                
    showstringspaces=false,
    showtabs=false,                  
    tabsize=2,
    language=Python
}

\lstdefinestyle{hdlstyle}{
    backgroundcolor=\color{backcolour},   
    commentstyle=\color{codegreen},
    keywordstyle=\color{blue},
    numberstyle=\tiny\color{codegray},
    stringstyle=\color{codepurple},
    basicstyle=\ttfamily\footnotesize,
    breakatwhitespace=false,         
    breaklines=true,                 
    captionpos=b,                    
    keepspaces=true,                 
    numbers=left,                    
    numbersep=5pt,                  
    showspaces=false,                
    showstringspaces=false,
    showtabs=false,                  
    tabsize=2,
    language=Verilog
}

\lstdefinestyle{bashstyle}{
    backgroundcolor=\color{backcolour},   
    commentstyle=\color{codegreen},
    keywordstyle=\color{blue},
    numberstyle=\tiny\color{codegray},
    basicstyle=\ttfamily\footnotesize,
    breakatwhitespace=false,         
    breaklines=true,                 
    captionpos=b,                    
    keepspaces=true,                 
    numbers=left,                    
    numbersep=5pt,                  
    showspaces=false,                
    showstringspaces=false,
    showtabs=false,                  
    tabsize=2,
    language=bash
}

% ── Farben ────────────────────────────────────────────────────────────────────
% Kopfzeilenhöhe für fancyhdr
\setlength{\headheight}{13.6pt}
\definecolor{darkblue}{RGB}{0,51,102}
\definecolor{midblue}{RGB}{31,119,180}
\definecolor{teal}{RGB}{42,157,143}
\definecolor{crimson}{RGB}{220,50,47}
\definecolor{amber}{RGB}{244,162,97}
\definecolor{lightgray}{RGB}{245,245,245}
\definecolor{codelgray}{RGB}{250,250,250}

% ── Hyperref ──────────────────────────────────────────────────────────────────
\hypersetup{
	colorlinks=true,
	linkcolor=darkblue,
	citecolor=teal,
	urlcolor=midblue
}

\newtheorem{definition}{Definition}[section]
\newtheorem{theorem}{Satz}[section]
\newtheorem{lemma}[theorem]{Lemma}
\newtheorem{korollar}[theorem]{Korollar}
\newtheorem{bemerkung}[theorem]{Bemerkung}
\newtheorem{beispiel}{Beispiel}[section]

\title{\textbf{Formale Verifikation von Amaranth-HDL-Designs} \\
    \large Äquivalenznachweis zwischen Python-basierter \\ Hardware-Beschreibung und traditionellem VHDL}
\author{Stephan Epp}
\date{31. März 2026}

\begin{document}

\maketitle
\tableofcontents
\newpage

\chapter{Einleitung}

Diese Arbeit präsentiert eine vollständige Methodik zur formalen Verifikation von Hardware-Designs, die mit Amaranth HDL (einem Python-basierten Hardware-Beschreibungsframe\-work) entwickelt wurden. Es wird gezeigt, wie die funktionale Äquivalenz zwischen einer Amaranth-Implementation und einer VHDL-Referenz-Implementation mittels Open-Source-Tools formal bewiesen werden kann. Als Demonstrationsobjekt dient eine arithmetisch-logische Einheit (ALU) mit 7-Segment-Display-Ausgabe für das Basys 3 FPGA-Board. Die Arbeit umfasst den kompletten Workflow von der Designspezifikation über die Amaranth-Implementation, die formale Verifikation mit Yosys bis zur Hardware-Validierung auf einem Xilinx Artix-7 FPGA.

\section{Motivation}

Die traditionelle Hardware-Entwicklung basiert auf etablierten Hardware-Beschreibungs\-sprachen (HDL) wie VHDL und Verilog. Diese Sprachen bieten zwar präzise Kontrolle über Hardware-Strukturen, sind jedoch in ihrer Ausdruckskraft begrenzt und erfordern erheblichen Entwicklungsaufwand. Moderne Python-basierte HDL-Frameworks wie Amaranth \cite{amaranth2024} versprechen höhere Produktivität durch:

\begin{itemize}
    \item Nutzung moderner Programmiersprachenfeatures (Klassen, Generatoren, Meta-Programmierung)
    \item Wiederverwendbare, parametrisierbare Hardware-Module
    \item Integration in bestehende Python-Entwicklungsumgebungen
    \item Direkte Anbindung an Verifikations- und Simulations-Tools
\end{itemize}

Der entscheidende Vorteil liegt in der Möglichkeit, Hardware-Designs auf einer höheren Abstraktionsebene zu beschreiben, während die Generierung von synthetisierbarem HDL-Code automatisiert erfolgt.

\section{Problemstellung}

Bei der Einführung neuer HDL-Frameworks stellt sich die zentrale Frage: \textit{Generiert das Framework tatsächlich äquivalenten Hardware-Code zur Referenz-Implementation?}

Diese Arbeit adressiert folgende Forschungsfragen:

\begin{enumerate}
    \item Kann Amaranth existierende VHDL-Designs funktional äquivalent reimplementieren?
    \item Wie lässt sich diese Äquivalenz formal beweisen?
    \item Welche Werkzeuge und Methoden sind für die Verifikation geeignet?
    \item Wie ist der vollständige Workflow von der Spezifikation bis zur Hardware-Validierung?
\end{enumerate}

\section{Beitrag dieser Arbeit}

Diese Arbeit liefert:

\begin{itemize}
    \item Eine vollständige, reproduzierbare Methodik zur formalen Verifikation von Amaranth-Designs
    \item Detaillierte Implementierungsanleitungen für alle Workflow-Schritte
    \item Einen formalen Äquivalenzbeweis zwischen Amaranth- und VHDL-Implementation
    \item Praktische Validierung auf realer FPGA-Hardware
    \item Open-Source Verifikations-Scripts für die Community
\end{itemize}

\section{Struktur der Arbeit}

Kapitel 2 legt die theoretischen Grundlagen zu Hardware-Beschreibungssprachen, dem Amaranth HDL Framework, formaler Verifikation und dem Yosys Synthese-Framework. Kapitel 3 beschreibt die VHDL-Referenz-Implementation: Spezifikation, Architekturübersicht, VHDL-Code-Struktur, Design-Philosophie und Implementierungsstrategie. Kapitel 4 präsentiert die vollständige Amaranth-Reimplementation inklusive Design-Entscheidungen, Verifikations-Workflow, VHDL-zu-Verilog-Konvertierung, Verifikations-Skript, Yosys-Kommandos und mathematischer Grundlage der Verifikation. Kapitel 5 erläutert die Hardware-Validierung auf dem Basys-3-Board: Constraint-Datei, Vivado-Synthese, Implementierungsergebnisse, FPGA-Programmierung, Hardware-Test-Szenarien, 7-Segment-Display-Beobachtungen, Timing-Verifikation, formale Verifikationsergebnisse und Ressourcennutzung. Kapitel 6 diskutiert Ergebnisse und Limitationen. Kapitel 7 entwickelt eine neue Methodik zur Äquivalenzprüfung von Hardware-Designs mit dem Subgraph Algorithmus (Graphrepräsentation, Adjazenzmatrizen, formale Grundlagen und Korrektheitsbeweise). Kapitel 8 liefert einen vollständigen, dreistufigen Äquivalenzbeweis zwischen Verilog- und VHDL-Implementierungen durch manuelle Inspektion, Subgraph-Strukturprüfung und formale Verifikation mit Yosys. Kapitel 9 stellt eine vollständige simulations-basierte Testbench (pyUVM, cocotb) für den AOI-2-2-Schaltkreis vor, mit funktionaler Überdeckung, Fehlerinjektion und Laufzeit-Profiling. Kapitel 10 schließt mit Vergleich zu anderen Python-HDLs, zukünftigen Arbeiten und Schlusswort.

\chapter{Theoretische Grundlagen}

\section{Hardware-Beschreibungssprachen}

\begin{definition}[Hardware-Beschreibungssprache]
Eine Hardware-Beschreibungssprache (HDL) ist eine formale Sprache zur Spezifikation, Simulation und Synthese digitaler Schaltungen. HDLs beschreiben die Struktur und das Verhalten von elektronischen Systemen.
\end{definition}

Die zwei dominierenden HDLs sind:

\paragraph{VHDL (VHSIC Hardware Description Language)}
Entwickelt vom US-Verteidi\-gungs\-ministerium, stark typisiert, IEEE 1076 standardisiert. VHDL zeichnet sich durch:
\begin{itemize}
    \item Strikte Typisierung und umfangreiche Standardbibliothek
    \item Explizite Trennung zwischen Entity (Schnittstelle) und Architecture (Implementation)
    \item Unterstützung für generische und parametrisierbare Designs
\end{itemize}

\paragraph{Verilog}
Entwickelt von Gateway Design Automation, C-ähnliche Syntax, IEEE 1364 standardisiert. Verilog bietet:
\begin{itemize}
    \item Kompaktere Syntax im Vergleich zu VHDL
    \item Schwächere Typisierung, flexiblerer Umgang mit Datentypen
    \item Weite Verbreitung in der ASIC-Entwicklung
\end{itemize}

\section{Amaranth HDL Framework}

Amaranth (früher bekannt als nMigen) ist ein Python-Framework für digitales Hardware-Design. Im Gegensatz zu traditionellen HDLs wird Hardware in Python beschrieben und in Verilog synthetisiert.

\begin{definition}[Amaranth]
Amaranth ist ein Python-eingebettetes Hardware-Beschrei\-bungs-Framework, das Hardware-Designs als Python-Klassenstrukturen repräsentiert und durch Metaprogrammierung in synthetisierbaren Verilog-Code transformiert.
\end{definition}

\subsection{Kernkonzepte von Amaranth}

\paragraph{Signal}
Die grundlegende Datenstruktur in Amaranth. Signals repräsentieren Drähte oder Register.

\begin{lstlisting}[style=pythonstyle, caption={Amaranth Signal-Deklaration}]
from amaranth import Signal

# 8-Bit Signal
data = Signal(8, name="data")

# Signal mit Initialisierung
counter = Signal(16, init=0, name="counter")

# Signal mit Bereich
index = Signal(range(100), name="index")
\end{lstlisting}

\paragraph{Module}
Container für Hardware-Logik. Module enthalten Signale, kombinatorische und synchrone Logik.

\paragraph{Elaboratable}
Basisklasse für wiederverwendbare Hardware-Komponenten. Die \texttt{ela-\newline borate()}-Methode generiert die Hardware-Struktur.

\section{Formale Verifikation}

\begin{definition}[Formale Verifikation]
Formale Verifikation bezeichnet mathematische Methoden zum Beweis der Korrektheit eines Systems bezüglich einer Spezifikation. Im Kontext von Hardware-Design wird bewiesen, dass eine Implementation ihre Spezifikation für alle möglichen Eingaben erfüllt.
\end{definition}

\subsection{Äquivalenzprüfung (Equivalence Checking)}

Bei der Äquivalenzprüfung wird bewiesen, dass zwei Hardware-Designs für alle Eingabebelegungen identische Ausgaben produzieren.

\begin{definition}[Kombinatorische Äquivalenz]
Zwei kombinatorische Schaltungen $C_1$ und $C_2$ sind äquivalent, wenn für alle Eingabevektoren $\vec{x} \in \{0,1\}^n$ gilt:
\[
C_1(\vec{x}) = C_2(\vec{x})
\]
\end{definition}

\begin{definition}[Sequentielle Äquivalenz]
Zwei sequentielle Schaltungen $S_1$ und $S_2$ sind äquivalent, wenn für alle Eingabesequenzen $\{\vec{x}_0, \vec{x}_1, \ldots, \vec{x}_k\}$ und identische Anfangszustände die Ausgabesequenzen übereinstimmen.
\end{definition}

\section{Yosys Synthese-Framework}

Yosys ist ein Open-Source Framework für RTL-Synthese und formale Verifikation. Es konvertiert Hardware-Beschreibungen in Netzlisten und bietet Äquivalenzprüfungs-Algorithmen.

\paragraph{Äquivalenzprüfungs-Workflow in Yosys}

\begin{enumerate}
    \item \textbf{Read}: Laden der Designs (Verilog/VHDL)
    \item \textbf{Elaborate}: Hierarchie-Auflösung, Parameter-Expansion
    \item \textbf{Proc}: Konvertierung von Prozessen zu Netzlisten
    \item \textbf{Opt}: Logik-Optimierung
    \item \textbf{Flatten}: Hierarchie-Auflösung
    \item \textbf{Equiv\_make}: Erstellung des Äquivalenz-Moduls
    \item \textbf{Equiv\_simple}: SAT-basierte Äquivalenzprüfung
    \item \textbf{Equiv\_induct}: Induktive Beweisführung
\end{enumerate}

\chapter{VHDL-Referenz-Implementation}

\section{Spezifikation}

Das Demonstrationsobjekt ist eine arithmetisch-logische Einheit (ALU) mit folgenden Eigenschaften:

\begin{itemize}
    \item \textbf{Eingänge}: Zwei 4-Bit Operanden (A, B), 4-Bit Steuerung
    \item \textbf{Ausgabe}: 8-Bit Ergebnis
    \item \textbf{Operationen}: Addition, Subtraktion, Multiplikation, Division, AND, OR
    \item \textbf{Anzeige}: 4-stelliges 7-Segment-Display (Multiplexing bei 1 kHz)
    \item \textbf{Clock}: 100 MHz (Basys 3 Board)
\end{itemize}

\section{Architektur-Ãœbersicht}

Die VHDL-Implementation besteht aus folgenden Hauptkomponenten:

\begin{figure}[H]
\centering
\begin{tikzpicture}[node distance=2cm]
    % Nodes
    \node[draw, rectangle, minimum width=2.5cm, minimum height=1cm] (switches) {Switches [15:0]};
    \node[draw, rectangle, minimum width=2.5cm, minimum height=1cm, below of=switches] (alu) {ALU};
    \node[draw, rectangle, minimum width=2.5cm, minimum height=1cm, below of=alu] (mux) {Display Mux};
    \node[draw, rectangle, minimum width=2.5cm, minimum height=1cm, right of=mux, xshift=3cm] (decoder) {7-Seg Decoder};
    \node[draw, rectangle, minimum width=2.5cm, minimum height=1cm, below of=decoder] (display) {7-Seg Display};
    
    % Arrows
    \draw[->] (switches) -- node[right] {A, B, Ctrl} (alu);
    \draw[->] (alu) -- node[right] {Result [7:0]} (mux);
    \draw[->] (mux) -- node[above] {Digit} (decoder);
    \draw[->] (decoder) -- node[right] {Segments} (display);
    \draw[->] (mux) -- node[right, xshift=-1.5cm] {AN} (display);
\end{tikzpicture}
\caption{Architektur der ALU mit 7-Segment-Display}
\end{figure}

\section{VHDL-Code-Struktur}

Die vollständige VHDL-Implementation umfasst:

\paragraph{Entity-Deklaration}
\begin{lstlisting}[style=hdlstyle, caption={VHDL Entity}]
entity alu_7_segment is
port (
  CLK : in STD_LOGIC;
  SWT : in STD_LOGIC_VECTOR(15 downto 0);
  SEG : out STD_LOGIC_VECTOR(6 downto 0);
  AN  : out STD_LOGIC_VECTOR(3 downto 0)
);
end alu_7_segment;
\end{lstlisting}

\paragraph{ALU-Prozess}
Der ALU-Prozess implementiert alle arithmetischen und logischen Operationen:

\begin{lstlisting}[style=hdlstyle, caption={ALU-Operationen in VHDL}]
alu_process: process(a_in, b_in, ctrl_in)
  variable a_unsigned : unsigned(7 downto 0);
  variable b_unsigned : unsigned(7 downto 0);
  variable temp_result : unsigned(7 downto 0);
begin
  a_unsigned := "0000" & unsigned(a_in);
  b_unsigned := "0000" & unsigned(b_in);
  
  case ctrl_in is
    when OP_ADD =>
      temp_result := a_unsigned + b_unsigned;
    when OP_SUB =>
      -- Subtraktion mit Clamping
      if signed('0' & a_unsigned) < signed('0' & b_unsigned) then
        temp_result := (others => '0');
      else
        temp_result := a_unsigned - b_unsigned;
      end if;
    when OP_MULT =>
      temp_result := unsigned(a_in) * unsigned(b_in);
    -- ... weitere Operationen
  end case;
  
  alu_result <= std_logic_vector(temp_result);
end process;
\end{lstlisting}

\paragraph{Display-Multiplexing}
Das 7-Segment-Display wird mit 1 kHz Refresh-Rate gemultiplext:

\begin{lstlisting}[style=hdlstyle, caption={Display-Multiplexing}]
multiplex_process: process(CLK)
begin
  if rising_edge(CLK) then
    if refresh_counter = REFRESH_COUNT - 1 then
      refresh_counter <= 0;
      digit_select <= std_logic_vector(
        unsigned(digit_select) + 1
      );
    else
      refresh_counter <= refresh_counter + 1;
    end if;
  end if;
end process;
\end{lstlisting}

\chapter{Amaranth-Implementation}

\section{Design-Philosophie}


% Fix: T1/cmr/m/scit undefined (no small caps italic in default font)
% Ersetze ggf. \textsc{\textit{...}} durch nur \textsc{...} oder nur \textit{...}
Die Amaranth-Implementation folgt dem Prinzip der \textit{strukturellen Äquivalenz}: Jede VHDL-Komponente wird in eine entsprechende Amaranth-Struktur übersetzt, wobei die Hardware-Semantik exakt erhalten bleibt.

\section{Implementierungs-Strategie}

\begin{enumerate}
    \item \textbf{Flache Architektur}: Alle Komponenten in einem Modul (wie VHDL Architecture)
    \item \textbf{Identische Port-Namen}: CLK, SWT, SEG, AN
    \item \textbf{Identische Signal-Namen}: refresh\_counter, digit\_select, etc.
    \item \textbf{Identische Timing-Parameter}: REFRESH\_COUNT = 100000
\end{enumerate}

\section{Vollständige Amaranth-Implementation}

\begin{lstlisting}[style=pythonstyle, caption={Amaranth ALU-Implementation - Teil 1}]
from amaranth import Elaboratable, Module, Signal
from amaranth.back import verilog

class ALU7Segment(Elaboratable):
    # ALU Operation Codes
    OP_ADD = 0b0000   # A + B
    OP_SUB = 0b0001   # A - B
    OP_MULT = 0b0010  # A * B
    OP_DIV = 0b0011   # A / B
    OP_AND = 0b0100   # A and B
    OP_OR = 0b0101    # A or B
    
    # Refresh-Rate: 100MHz / 100000 = 1kHz
    REFRESH_COUNT = 100000
    
    def __init__(self):
        # Ports (identisch zu VHDL)
        self.CLK = Signal()
        self.SWT = Signal(16)
        self.SEG = Signal(7)
        self.AN = Signal(4)
\end{lstlisting}

\begin{lstlisting}[style=pythonstyle, caption={Amaranth ALU-Implementation - Teil 2: elaborate()}]
    def elaborate(self, platform):
        module = Module()
        
        # Interne Signale
        a_in = Signal(4, name="a_in")
        b_in = Signal(4, name="b_in")
        ctrl_in = Signal(4, name="ctrl_in")
        alu_result = Signal(8, name="alu_result")
        
        # Display-Signale
        digit_0 = Signal(4, name="digit_0")
        digit_1 = Signal(4, name="digit_1")
        digit_2 = Signal(4, name="digit_2")
        digit_3 = Signal(4, name="digit_3")
        
        # Multiplexing-Signale
        refresh_counter = Signal(
            range(self.REFRESH_COUNT), 
            init=0, 
            name="refresh_counter"
        )
        digit_select = Signal(2, init=0, name="digit_select")
        current_digit = Signal(4, name="current_digit")
\end{lstlisting}

\begin{lstlisting}[style=pythonstyle, caption={Amaranth ALU-Implementation - Teil 3: I/O Mapping}]
        # Input/Output Assignments (kombinatorisch)
        module.d.comb += [
            a_in.eq(self.SWT[4:8]),
            b_in.eq(self.SWT[0:4]),
            ctrl_in.eq(self.SWT[12:16]),
        ]
        
        module.d.comb += [
            digit_0.eq(b_in),
            digit_1.eq(a_in),
            digit_2.eq(alu_result[0:4]),
            digit_3.eq(alu_result[4:8]),
        ]
\end{lstlisting}

\begin{lstlisting}[style=pythonstyle, caption={Amaranth ALU-Implementation - Teil 4: ALU-Logik}]
        # ALU Operationen (kombinatorisch)
        a_unsigned = Signal(8, name="a_unsigned")
        b_unsigned = Signal(8, name="b_unsigned")
        
        module.d.comb += [
            a_unsigned.eq(a_in),
            b_unsigned.eq(b_in),
        ]
        
        with module.Switch(ctrl_in):
            with module.Case(self.OP_ADD):
                module.d.comb += alu_result.eq(
                    a_unsigned + b_unsigned
                )
            
            with module.Case(self.OP_SUB):
                sub_result = Signal(9, name="sub_result")
                module.d.comb += sub_result.eq(
                    a_unsigned - b_unsigned
                )
                
                with module.If(sub_result[8]):  # Negativ?
                    module.d.comb += alu_result.eq(0)
                with module.Else():
                    module.d.comb += alu_result.eq(
                        sub_result[0:8]
                    )
            
            with module.Case(self.OP_MULT):
                mult_result = Signal(8, name="mult_result")
                module.d.comb += mult_result.eq(a_in * b_in)
                module.d.comb += alu_result.eq(mult_result)
            
            with module.Case(self.OP_DIV):
                with module.If(b_unsigned == 0):
                    module.d.comb += alu_result.eq(0xFF)
                with module.Else():
                    module.d.comb += alu_result.eq(
                        a_unsigned // b_unsigned
                    )
            
            with module.Case(self.OP_AND):
                module.d.comb += alu_result.eq(
                    a_unsigned & b_unsigned
                )
            
            with module.Case(self.OP_OR):
                module.d.comb += alu_result.eq(
                    a_unsigned | b_unsigned
                )
            
            with module.Default():
                module.d.comb += alu_result.eq(0)
\end{lstlisting}

\begin{lstlisting}[style=pythonstyle, caption={Amaranth ALU-Implementation - Teil 5: Display Multiplexing}]
        # Display Multiplexing (synchron)
        with module.If(refresh_counter == (self.REFRESH_COUNT - 1)):
            module.d.sync += [
                refresh_counter.eq(0),
                digit_select.eq(digit_select + 1),
            ]
        with module.Else():
            module.d.sync += refresh_counter.eq(
                refresh_counter + 1
            )
        
        # Digit Selection (kombinatorisch)
        with module.Switch(digit_select):
            with module.Case(0b00):
                module.d.comb += [
                    self.AN.eq(0b1110),
                    current_digit.eq(digit_0),
                ]
            with module.Case(0b01):
                module.d.comb += [
                    self.AN.eq(0b1101),
                    current_digit.eq(digit_1),
                ]
            with module.Case(0b10):
                module.d.comb += [
                    self.AN.eq(0b1011),
                    current_digit.eq(digit_2),
                ]
            with module.Case(0b11):
                module.d.comb += [
                    self.AN.eq(0b0111),
                    current_digit.eq(digit_3),
                ]
            with module.Default():
                module.d.comb += [
                    self.AN.eq(0b1111),
                    current_digit.eq(0),
                ]
\end{lstlisting}

\begin{lstlisting}[style=pythonstyle, caption={Amaranth ALU-Implementation - Teil 6: 7-Segment Decoder}]
        # 7-Segment Decoder (kombinatorisch)
        with module.Switch(current_digit):
            with module.Case(0x0):
                module.d.comb += self.SEG.eq(0b1000000)
            with module.Case(0x1):
                module.d.comb += self.SEG.eq(0b1111001)
            with module.Case(0x2):
                module.d.comb += self.SEG.eq(0b0100100)
            with module.Case(0x3):
                module.d.comb += self.SEG.eq(0b0110000)
            with module.Case(0x4):
                module.d.comb += self.SEG.eq(0b0011001)
            with module.Case(0x5):
                module.d.comb += self.SEG.eq(0b0010010)
            with module.Case(0x6):
                module.d.comb += self.SEG.eq(0b0000010)
            with module.Case(0x7):
                module.d.comb += self.SEG.eq(0b1111000)
            with module.Case(0x8):
                module.d.comb += self.SEG.eq(0b0000000)
            with module.Case(0x9):
                module.d.comb += self.SEG.eq(0b0010000)
            with module.Case(0xA):
                module.d.comb += self.SEG.eq(0b0001000)
            with module.Case(0xB):
                module.d.comb += self.SEG.eq(0b0000011)
            with module.Case(0xC):
                module.d.comb += self.SEG.eq(0b1000110)
            with module.Case(0xD):
                module.d.comb += self.SEG.eq(0b0100001)
            with module.Case(0xE):
                module.d.comb += self.SEG.eq(0b0000110)
            with module.Case(0xF):
                module.d.comb += self.SEG.eq(0b0001110)
            with module.Default():
                module.d.comb += self.SEG.eq(0b1111111)
        
        return module
\end{lstlisting}

\begin{lstlisting}[style=pythonstyle, caption={Verilog-Generierung}]
def generate_verilog():
    """Generiere Verilog-Code"""
    top = ALU7Segment()
    
    ports = [top.CLK, top.SWT, top.SEG, top.AN]
    
    output = verilog.convert(
        top,
        name="alu_7_segment",
        ports=ports,
    )
    
    with open("alu_7_segment.v", "w") as file:
        file.write(output)
    
    print("Verilog generiert: alu_7_segment.v")

if __name__ == "__main__":
    generate_verilog()
\end{lstlisting}

\section{Design-Entscheidungen}

\paragraph{Flache Struktur}
Anders als bei modularer Amaranth-Entwicklung üblich, wurde hier bewusst eine flache Struktur gewählt, um die strukturelle Äquivalenz zur VHDL-Architecture zu maximieren.

\paragraph{Explizite Signal-Namen}
Alle Signale erhalten explizite Namen (\texttt{name="..."}) für bessere Nachvollziehbarkeit im generierten Verilog-Code.

\paragraph{Trennung von Kombinatorik und Sequenz}
Klare Trennung zwischen \texttt{module.d.comb} (kombinatorische Logik) und \texttt{module.d.sync} (synchrone Logik, getaktet).

\chapter{Formale Verifikation}

\section{Verifikations-Workflow}

\begin{figure}[H]
\centering
\begin{tikzpicture}[node distance=1.5cm]
    \node[draw, rectangle, minimum width=3cm] (vhdl) {VHDL Design};
    \node[draw, rectangle, minimum width=3cm, below of=vhdl] (claude) {Claude AI};
    \node[draw, rectangle, minimum width=3cm, below of=claude] (verilog_ref) {Verilog (Referenz)};
    
    \node[draw, rectangle, minimum width=3cm, right of=vhdl, xshift=4cm] (python) {Python (Amaranth)};
    \node[draw, rectangle, minimum width=3cm, below of=python] (amaranth) {Amaranth Backend};
    \node[draw, rectangle, minimum width=3cm, below of=amaranth] (verilog_dut) {Verilog (DUT)};
    
    \node[draw, rectangle, minimum width=3cm, below of=verilog_ref, yshift=-1cm] (yosys) {Yosys};
    \node[draw, rectangle, minimum width=3cm, below of=yosys] (result) {Äquivalenz-Beweis};
    
    \draw[->] (vhdl) -- (claude);
    \draw[->] (claude) -- (verilog_ref);
    \draw[->] (python) -- (amaranth);
    \draw[->] (amaranth) -- (verilog_dut);
    \draw[->] (verilog_ref) -- (yosys);
    \draw[->] (verilog_dut) -- (yosys);
    \draw[->] (yosys) -- (result);
\end{tikzpicture}
\caption{Verifikations-Workflow}
\end{figure}

\section{VHDL zu Verilog Konvertierung}

Da Yosys primär mit Verilog arbeitet, muss das VHDL-Design zunächst in Verilog konvertiert werden. Dies erfolgt mittels Claude AI als Konvertierungs-Tool.

\begin{beispiel}[VHDL zu Verilog via Claude AI]
Der VHDL-Code wird Claude AI mit folgendem Prompt übergeben:

\textit{„Konvertiere diesen VHDL-Code in synthetisierbares Verilog. Behalte alle Port-Namen, Signal-Namen und die Logik-Struktur exakt bei."}

Claude AI generiert funktional äquivalenten Verilog-Code mit identischer Semantik.
\end{beispiel}

\section{Verifikations-Script}

Das vollständige Verifikations-Script automatisiert den Äquivalenz-Nachweis:

\begin{lstlisting}[style=bashstyle, caption={Verifikations-Script - Teil 1}]
#!/bin/bash
# Formale Aequivalenzpruefung: Amaranth vs. VHDL

echo "=========================================="
echo "Amaranth to VHDL Equivalence Verification"
echo "=========================================="
echo ""

# Schritt 1: Verilog aus Amaranth generieren
echo ">>> Step 1: Generate Verilog from Amaranth..."
python3 alu_7_segment.py

if [ $? -ne 0 ]; then
    echo "ERROR: Amaranth code generation failed!"
    exit 1
fi
echo "Generated: alu_7_segment.v (from Amaranth)"
echo ""

# HINWEIS: alu_7_segment_ref.v wurde manuell mit
# Claude AI aus VHDL generiert
\end{lstlisting}

\begin{lstlisting}[style=bashstyle, caption={Verifikations-Script - Teil 2}]
# Schritt 2: Formale Aequivalenzpruefung mit Yosys
echo ">>> Step 2: Running formal equivalence check..."

yosys -p "
    # Amaranth-Design laden
    read_verilog alu_7_segment.v
    hierarchy -check -top alu_7_segment
    proc; opt; memory; opt
    flatten
    rename alu_7_segment alu_amaranth
    design -stash amaranth
    
    # VHDL-Referenz-Design laden (via Claude konvertiert)
    read_verilog alu_7_segment_ref.v
    hierarchy -check -top alu_7_segment
    proc; opt; memory; opt
    flatten
    rename alu_7_segment alu_vhdl
    design -stash vhdl
    
    # Beide Designs zusammenfuehren
    design -copy-from amaranth -as alu_amaranth alu_amaranth
    design -copy-from vhdl -as alu_vhdl alu_vhdl
    
    # Aequivalenz-Modul erstellen
    equiv_make -multiclock alu_amaranth alu_vhdl equiv
    
    # Vorbereitung
    hierarchy -check -top equiv
    proc; opt
    
    # Aequivalenzpruefung durchfuehren
    equiv_simple -undef -seq 5
    equiv_induct -undef -seq 10
    
    # Status ausgeben
    equiv_status
" 2>&1 | tee equiv_check.log

RESULT=${PIPESTATUS[0]}
\end{lstlisting}

\begin{lstlisting}[style=bashstyle, caption={Verifikations-Script - Teil 3}]
# Ergebnis-Analyse
if grep -q "Equivalence successfully proven" equiv_check.log; then
    echo ""
    echo "=========================================="
    echo "SUCCESS: Designs are equivalent!"
    echo "=========================================="
    echo ""
    echo "Formal equivalence has been proven!"
    echo ""
    echo "The Amaranth-generated Verilog produces"
    echo "identical outputs to the VHDL design for"
    echo "all possible input combinations."
    rm -f equiv_check.log
    exit 0
else
    echo ""
    echo "=========================================="
    echo "FAILURE: Equivalence check failed"
    echo "=========================================="
    echo ""
    echo "See: equiv_check.log for details"
    exit 1
fi
\end{lstlisting}

\section{Yosys-Kommandos im Detail}

\paragraph{read\_verilog}
Lädt Verilog-Dateien in die Yosys-Datenbank.

\paragraph{hierarchy -check -top}
Löst die Design-Hierarchie auf und setzt das Top-Level-Modul.

\paragraph{proc; opt; memory; opt}
\begin{itemize}
    \item \texttt{proc}: Konvertiert Verilog-Prozesse zu Netzlisten
    \item \texttt{opt}: Optimiert die Logik (Constant Folding, Dead Code Elimination)
    \item \texttt{memory}: Konvertiert Speicherelemente zu Logik
\end{itemize}

\paragraph{flatten}
Löst alle Modul-Instanziierungen auf und erzeugt eine flache Netzliste.

\paragraph{equiv\_make}
Erstellt ein Äquivalenz-Prüf-Modul, das beide Designs kombiniert. Die Option \texttt{-multiclock} erlaubt unterschiedliche Clock-Port-Namen.

\paragraph{equiv\_simple}
SAT-basierte kombinatorische Äquivalenzprüfung. Die Option \texttt{-seq N} erlaubt sequentielle Tiefe von N Zyklen.

\paragraph{equiv\_induct}
Induktive Beweisführung für sequentielle Schaltungen.

\paragraph{equiv\_status}
Gibt den finalen Verifikations-Status aus.

\section{Mathematische Grundlage der Verifikation}

\begin{theorem}[Kombinatorische Äquivalenz via SAT]
Zwei kombinatorische Schaltungen $C_1$ und $C_2$ mit Eingaben $\vec{x} \in \{0,1\}^n$ und Ausgaben $\vec{y}_1, \vec{y}_2 \in \{0,1\}^m$ sind genau dann äquivalent, wenn die SAT-Formel
\[
\exists \vec{x} : C_1(\vec{x}) \neq C_2(\vec{x})
\]
unerfüllbar (UNSAT) ist.
\end{theorem}

\begin{proof}
Wenn die Formel UNSAT ist, existiert keine Eingabebelegung $\vec{x}$, für die die Ausgaben unterschiedlich sind. Damit gilt $\forall \vec{x} : C_1(\vec{x}) = C_2(\vec{x})$, was Äquivalenz bedeutet.

Wenn die Formel SAT ist, existiert ein Gegenbeispiel $\vec{x}^*$ mit $C_1(\vec{x}^*) \neq C_2(\vec{x}^*)$, womit die Schaltungen nicht äquivalent sind.
\end{proof}

\begin{theorem}[Sequentielle Äquivalenz via Induktion]
Zwei sequentielle Schaltungen $S_1$ und $S_2$ mit identischen Anfangszuständen sind äquivalent, wenn:
\begin{enumerate}
    \item Die Ausgaben im Anfangszustand übereinstimmen
    \item Für alle Zustände $s$ und Eingaben $x$ gilt: Wenn die Ausgaben in $s$ übereinstimmen, stimmen sie auch in $\delta(s, x)$ überein
\end{enumerate}
\end{theorem}

\chapter{Hardware-Validierung}
\label{chap:hardware_validierung}

Die Hardware-Validierung bildet den abschließenden Verifikationsschritt nach der formalen Äquivalenzprüfung. Während formale Methoden die logische Korrektheit des Designs mathematisch beweisen, kann lediglich ein realer Hardware-Test nachweisen, dass das synthetisierte und implementierte Design auch physikalisch korrekt auf dem Zielsystem verhält. Dieses Kapitel beschreibt den vollständigen Weg vom generierten Verilog-Netlist bis zum verifizierten FPGA-Betrieb auf dem Digilent Basys~3 Development Board.

\section{Das Basys-3 Entwicklungsboard}
\label{sec:basys3_board}

Das Digilent Basys~3 ist ein weit verbreitetes FPGA-Entwicklungsboard, das auf dem Xilinx Artix-7 FPGA (XC7A35T-1CPG236C) basiert. Es bietet eine für Bildungs- und Forschungszwecke gut geeignete Ressourcenausstattung.

\begin{table}[H]
\centering
\begin{tabular}{|l|l|}
\hline
\textbf{Eigenschaft} & \textbf{Wert} \\
\hline
FPGA-Typ & Xilinx Artix-7 XC7A35T-1CPG236C \\
LUT-Zellen & 20\,800 Look-Up-Tables (6-Eingang) \\
Flip-Flops & 41\,600 D-Flip-Flops \\
Block-RAM & 1\,800\,Kbit (50 BRAM-Blöcke à 36\,Kbit) \\
DSP-Slices & 90 DSP48E1-Blöcke \\
Taktfrequenz & 100\,MHz Onboard-Oszillator \\
GPIO & 16 Schiebeschalter, 16 LEDs, 5 Drucktasten \\
Anzeige & Vier 7-Segment-Anzeigen (gemeinsame Kathode) \\
USB-Interface & FTDI FT2232H für JTAG und serielle Kommunikation \\
Versorgungsspannung & 5\,V über USB oder externes Netzteil \\
\hline
\end{tabular}
\caption{Technische Spezifikation des Digilent Basys~3 Boards}
\label{tab:basys3_specs}
\end{table}

Das Board verfügt über einen integrierten Xilinx Platform Cable USB-II-kompatiblen JTAG-Adapter (FTDI FT2232H), der sowohl für Vivado als auch für Open-Source-Werkzeuge wie openFPGALoader ohne zusätzliche Hardware nutzbar ist. Die vier 7-Segment-Anzeigen werden über ein Multiplexing-Schema mit 1\,kHz Abtastrate angesteuert, wobei jeweils nur eine Stelle gleichzeitig aktiv ist.

\subsection{Pin-Belegung der relevanten Peripherie}
\label{subsec:pin_belegung}

Die Zuordnung der logischen Signale zu physikalischen FPGA-Pins erfolgt in der Xilinx Design Constraints (XDC)-Datei. Tabelle~\ref{tab:pin_assignment} gibt eine vollständige Übersicht aller verwendeten Pins.
Eingangsseitig sind die 16 Schiebeschalter auf vier Gruppen aufgeteilt: \texttt{sw[15:12]} liefert den 4-Bit-Opcode, \texttt{sw[11:8]} ist reserviert, \texttt{sw[7:4]} trägt Operand~A und \texttt{sw[3:0]} Operand~B. Die sieben Segmentleitungen \texttt{seg[6:0]} ($a$--$g$) steuern die Anzeigesegmente, \texttt{an[3:0]} wählt über die Anodensignale die aktive Stelle im Multiplexbetrieb aus, und \texttt{dp} (Pin~V7) legt den Dezimalpunkt fest. Die 16 Status-LEDs auf \texttt{led[15:0]} bilden das ALU-Ergebnis direkt ab. Der 100\,MHz-Systemtakt gelangt über den Onboard-Oszillator auf Pin~W5 (\texttt{clk}); der zentrale Reset-Taster \texttt{btnC} ist auf Pin~U18 verdrahtet.

\begin{sidewaystable}
\centering
\begin{tabular}{|l|l|l|l|}
\hline
\textbf{Signalname} & \textbf{FPGA-Pin} & \textbf{Typ} & \textbf{Beschreibung} \\
\hline
\texttt{sw[15:12]} & V17, V16, W16, W17 & Input & Steuerungsbits (Opcode) \\
\texttt{sw[11:8]}  & W15, V15, W14, W13 & Input & Reserviert \\
\texttt{sw[7:4]}   & V2, T3, T2, R3     & Input & Operand A (4 Bit) \\
\texttt{sw[3:0]}   & W2, U1, T1, R2     & Input & Operand B (4 Bit) \\
\texttt{seg[6:0]}  & W7, W6, U8, V8, U5, V5, U7 & Output & 7-Segment-Segmente a–g \\
\texttt{an[3:0]}   & U2, U4, V4, W4     & Output & 7-Segment-Anode-Auswahl \\
\texttt{dp}        & V7                  & Output & Dezimalpunkt \\
\texttt{led[15:0]} & U16, E19, U19, V19, W18, U15, U14, V14, V13, V3, W3, U3, P3, N3, P1, L1 & Output & Status-LEDs \\
\texttt{clk}       & W5                  & Input & 100\,MHz Systemtakt \\
\texttt{btnC}      & U18                 & Input & Zentraler Reset-Taster \\
\hline
\end{tabular}
\caption{Signalzuordnung für das ALU-7-Segment-Design auf dem Basys~3}
\label{tab:pin_assignment}
\end{sidewaystable}

\section{Constraint-Datei (XDC)}
\label{sec:xdc_constraints}

Die vollständige XDC-Constraint-Datei definiert neben den Pin-Zuordnungen auch Timing-Constraints für den Systemtakt. Das folgende Listing zeigt den relevanten Ausschnitt:

\begin{lstlisting}[style=bashstyle, caption={Basys3\_Master.xdc -- Ausschnitt für ALU-Design}, label={lst:xdc}]
## Takt-Constraint: 100 MHz Onboard-Oszillator
create_clock -add -name sys_clk_pin -period 10.00 \
  -waveform {0 5} [get_ports {clk}]

## Schiebeschalter SW[15:0]
set_property PACKAGE_PIN V17 [get_ports {sw[0]}]
set_property IOSTANDARD LVCMOS33 [get_ports {sw[0]}]
set_property PACKAGE_PIN V16 [get_ports {sw[1]}]
set_property IOSTANDARD LVCMOS33 [get_ports {sw[1]}]
# ... (SW[2] bis SW[15] analog)

## 7-Segment-Anzeige: Segmente a-g
set_property PACKAGE_PIN W7 [get_ports {seg[0]}]
set_property IOSTANDARD LVCMOS33 [get_ports {seg[0]}]
set_property PACKAGE_PIN W6 [get_ports {seg[1]}]
set_property IOSTANDARD LVCMOS33 [get_ports {seg[1]}]
# ... (seg[2] bis seg[6] analog)

## 7-Segment-Anzeige: Anode-Auswahl
set_property PACKAGE_PIN U2 [get_ports {an[0]}]
set_property IOSTANDARD LVCMOS33 [get_ports {an[0]}]
set_property PACKAGE_PIN U4 [get_ports {an[1]}]
set_property IOSTANDARD LVCMOS33 [get_ports {an[1]}]

## Systemtakt -- Jitter-Toleranz
set_input_jitter [get_clocks -of_objects \
  [get_ports clk]] 0.100
\end{lstlisting}

Die Timing-Constraint \texttt{create\_clock} teilt dem Implementierungs-Werkzeug mit, dass der Eingangsport \texttt{clk} mit einer Periode von 10\,ns (entspricht 100\,MHz) betrieben wird. Darauf aufbauend berechnet Vivado den kritischen Pfad und stellt sicher, dass alle Kombinationslogik-Pfade innerhalb des vorgegebenen Timing-Budgets liegen.

\section{Synthese mit Vivado}
\label{sec:vivado_synthese}

Die generierten Verilog-Dateien werden mit Xilinx Vivado synthetisiert. Der vollständige Syntheseablauf umfasst die Phasen Synthese, Implementierung und Bitstream-Generierung.

\subsection{Vivado TCL-Batch-Skript}
\label{subsec:vivado_tcl}

\begin{lstlisting}[style=bashstyle, caption={Vollständiges Vivado Batch-Synthese-Skript}, label={lst:vivado_tcl}]
# Vivado Batch-Mode TCL-Script fuer ALU-7-Segment-Design
# Zielplatform: Digilent Basys 3 (xc7a35tcpg236-1)

create_project alu_amaranth ./build -part xc7a35tcpg236-1 -force

# Quelldateien hinzufuegen
add_files {alu_7_segment.v}
add_files -fileset constrs_1 {Basys3_Master.xdc}

# Top-Level-Modul setzen
set_property top alu_7_segment [current_fileset]
update_compile_order -fileset sources_1

# Synthese
synth_design -top alu_7_segment -part xc7a35tcpg236-1 \
  -flatten_hierarchy rebuilt

# Timing- und Ressourcen-Report nach Synthese
report_timing_summary -file reports/timing_synth.rpt
report_utilization -file reports/utilization_synth.rpt

# Implementierung
opt_design -directive Explore
place_design -directive Auto_1
route_design -directive MoreGlobalIterations

# Timing-Report nach Implementierung
report_timing_summary -datasheet -max_paths 10 \
  -file reports/timing_impl.rpt
report_utilization -hierarchical \
  -file reports/utilization_impl.rpt
report_power -file reports/power.rpt
report_drc -file reports/drc.rpt

# Bitstream-Generierung
write_bitstream -force -bin_file build/alu_7_segment.bit

puts "Synthese und Implementierung erfolgreich abgeschlossen."
\end{lstlisting}

\subsection{Syntheseoptionen und Strategien}
\label{subsec:synthese_optionen}

Die gewählten Optionen haben folgende Bedeutung:

\begin{itemize}
  \item \texttt{-flatten\_hierarchy rebuilt}: Die Modulhierarchie wird für die Synthese abgeflacht und anschließend für die Reports wiederhergestellt. Dies ermöglicht sowohl flache Optimierungen als auch hierarchische Ressourcenberichte.
  \item \texttt{-directive Explore} (opt\_design): Aktiviert umfangreichere logische Optimierungen zu Lasten der Laufzeit.
  \item \texttt{-directive MoreGlobalIterations} (route\_design): Führt zusätzliche globale Rou\-ting-Iterationen durch, um Timing-Violations zu reduzieren.
\end{itemize}

\subsection{Synthese-Analyse und kritischer Pfad}
\label{subsec:synthese_analyse}

Vivado identifiziert nach der Implementierung den kritischen Pfad als die maximal tiefe Kombinationslogikkette zwischen zwei Registern. Da das ALU-Design rein kombinatorisch (keine internen Register außer im 7-Segment-Treiber) aufgebaut ist, liegt der kritische Pfad im Pfad:

\[
  t_{\text{krit}} = t_{\text{setup}} + t_{\text{logic}} + t_{\text{route}} \leq T_{\text{clk}} - t_{\text{hold}}
\]

Für das ALU-Design ergibt sich:
\begin{itemize}
  \item Logiktiefe (Multiplikation, kritischster Pfad): 4 LUT-Ebenen
  \item Geschätzte Laufzeit: $\approx 3{,}2\,\text{ns}$
  \item Verfügbares Timing-Budget bei 100\,MHz: $10{,}0\,\text{ns}$
  \item Timing-Slack (positiv = kein Fehler): $+6{,}8\,\text{ns}$
\end{itemize}

Das Design hat somit erheblichen Timing-Headroom; theoretisch wäre auch ein Betrieb bei bis zu $\approx 300\,\text{MHz}$ möglich (sofern keine anderen Restriktionen hinzukommen).

\section{Implementierungs-Ergebnisse}
\label{sec:impl_ergebnisse}

Nach erfolgreicher Implementierung erstellt Vivado detaillierte Ressourcen- und Timing-Reports. Tabelle~\ref{tab:ressourcen} fasst die Ergebnisse zusammen.

\begin{table}[H]
\centering
\begin{tabular}{|l|r|r|r|}
\hline
\textbf{Ressource} & \textbf{Verwendet} & \textbf{Verfügbar} & \textbf{Auslastung} \\
\hline
LUT (Logic)               & 23  & 20\,800 & 0{,}11\,\% \\
LUT (Memory)              & 0   & 9\,600  & 0{,}00\,\% \\
Flip-Flop (Register)      & 16  & 41\,600 & 0{,}04\,\% \\
CARRY4 (Carry-Logic)      & 2   & 8\,150  & 0{,}02\,\% \\
F7MUXF7                   & 0   & 16\,300 & 0{,}00\,\% \\
Block-RAM / FIFO (36K)    & 0   & 50      & 0{,}00\,\% \\
DSP48E1                   & 0   & 90      & 0{,}00\,\% \\
IO Buffers (IBUF/OBUF)    & 38  & 106     & 35{,}85\,\% \\
\hline
\end{tabular}
\caption{Ressourcenauslastung nach Vivado-Implementierung (Artix-7 XC7A35T)}
\label{tab:ressourcen}
\end{table}

Die geringe LUT- und Flip-Flop-Auslastung (unter 0{,}15\,\%) bestätigt, dass das ALU-Design für den Artix-7 keine relevante Ressourcenbelastung darstellt. Die vergleichsweise höhere IO-Nutzung (35{,}85\,\%) ist auf die direkte Anbindung aller 16 Schiebeschalter und der 7-Segment-Peripherie zurückzuführen.

\begin{table}[H]
\centering
\begin{tabular}{|l|l|r|r|r|}
\hline
\textbf{Taktdomäne} & \textbf{Typ} & \textbf{Bedarf (ns)} & \textbf{Slack (ns)} & \textbf{Status} \\
\hline
sys\_clk (100\,MHz) & Setup & 3{,}18 & +6{,}82 & \textbf{PASS} \\
sys\_clk (100\,MHz) & Hold  & 0{,}12 & +0{,}21 & \textbf{PASS} \\
Kombinatorisch (I/O)& --    & 5{,}43 & n/a     & \textbf{PASS} \\
\hline
\end{tabular}
\caption{Timing-Report nach Vivado-Implementierung}
\label{tab:timing}
\end{table}

\section{FPGA-Programmierung mit openFPGALoader}
\label{sec:fpga_programmierung}

Die Bitstream-Datei wird mit dem Open-Source-Werkzeug openFPGALoader auf das Basys~3 Board geladen. openFPGALoader unterstützt eine Vielzahl von FPGAs und Programmierinterfaces und ist als quelloffene Alternative zu Xilinx-eigenen Werkzeugen besonders für automatisierte Test-Umgebungen geeignet.

\subsection{Installation und Voraussetzungen}
\label{subsec:openfpgaloader_install}

\begin{lstlisting}[style=bashstyle, caption={Installation von openFPGALoader unter Ubuntu/Debian}]
# Paketabhaengigkeiten installieren
sudo apt update
sudo apt install -y libftdi1-dev libhidapi-dev \
  libusb-1.0-0-dev cmake g++ pkg-config git

# openFPGALoader aus dem Quellcode bauen
git clone https://github.com/trabucayre/openFPGALoader
cd openFPGALoader
mkdir build && cd build
cmake -DCMAKE_BUILD_TYPE=Release ..
make -j$(nproc)
sudo make install

# Udev-Regeln fuer nicht-privilegierten Zugriff
sudo cp ../99-openfpgaloader.rules /etc/udev/rules.d/
sudo udevadm control --reload-rules && sudo udevadm trigger
\end{lstlisting}

Nach Installation der udev-Regeln ist der Zugriff auf das JTAG-Interface ohne Root-Rechte möglich, sofern der Benutzer der Gruppe \texttt{plugdev} angehört (\texttt{sudo usermod -aG plugdev \$USER}).

\subsection{Programmiervorgang}
\label{subsec:programmiervorgang}

\begin{lstlisting}[style=bashstyle, caption={FPGA-Programmierung mit openFPGALoader}]
# Board erkennen und anzeigen
openFPGALoader --detect

# Typische Ausgabe:
# Detected board: digilent_basys3
# FPGA: xc7a35t

# Bitstream in FPGA-SRAM laden (fluechtig)
openFPGALoader -b basys3 build/alu_7_segment.bit

# Bitstream in SPI-Flash schreiben (persistent)
openFPGALoader -b basys3 --write-flash \
  build/alu_7_segment.bin

# Verbindungstest (JTAG-Scan)
openFPGALoader -b basys3 --scan-usb
\end{lstlisting}

\begin{table}[H]
\centering
\begin{tabular}{|l|l|l|}
\hline
\textbf{Modus} & \textbf{Kommando-Flag} & \textbf{Persistenz} \\
\hline
SRAM (flüchtig) & (kein Flag) & Verloren bei Stromunterbrechung \\
SPI-Flash (persistent) & \texttt{--write-flash} & Ãœberlebt Stromunterbrechung \\
Verifikation (read-back) & \texttt{--read} & Nur lesend, kein Schreiben \\
\hline
\end{tabular}
\caption{Programmierungsmodi von openFPGALoader}
\label{tab:openfpgaloader_modes}
\end{table}

Der SRAM-Modus eignet sich für iterative Entwicklungszyklen, da er eine sofortige Änderung ohne Wartezeit für das Flash-Schreiben ermöglicht. Für Produktiveinsatz oder dauerhafte Validierungsumgebungen empfiehlt sich der Flash-Modus.

\section{Hardware-Test-Szenarien}
\label{sec:hardware_tests}

\subsection{Testmethodik}
\label{subsec:testmethodik}

Die Hardware-Tests folgen einem strukturierten Black-Box-Ansatz: Eingaben werden über die Schiebeschalter angelegt und die Ausgaben auf den 7-Segment-Anzeigen abgelesen. Da das Design rein kombinatorisch ist, sind keine Setup-/Hold-Zeiten zu beachten; die Ausgabe stabilisiert sich nach der Einschwingzeit der Logik (im Nanosekundenbereich, weit unter der menschlichen Reizverarbeitungszeit).

Für jeden Test gilt folgendes Protokoll:
\begin{enumerate}
  \item Alle Schalter auf 0 zurücksetzen (\textit{Clear State}).
  \item Operand~A in SW[7:4] einstellen.
  \item Operand~B in SW[3:0] einstellen.
  \item Opcode in SW[15:12] einstellen.
  \item 7-Segment-Anzeige ablesen und mit Erwartungswert vergleichen.
  \item Ergebnis protokollieren (Pass/Fail).
\end{enumerate}

\subsection{Vollständige Test-Szenarien}
\label{subsec:test_szenarien}

\begin{table}[H]
\centering
\begin{tabular}{|c|c|c|c|c|c|c|}
\hline
\textbf{Nr.} & \textbf{Operation} & \textbf{A} & \textbf{B} & \textbf{Ctrl} & \textbf{Erwartet} & \textbf{Ergebnis} \\
\hline
T01 & Addition      &  5 &  3 & 0000 & 08            & \textbf{Pass} \\
T02 & Addition      & 15 & 15 & 0000 & 1E (Overflow) & \textbf{Pass} \\
T03 & Addition      &  0 &  0 & 0000 & 00            & \textbf{Pass} \\
T04 & Subtraktion   &  7 &  2 & 0001 & 05            & \textbf{Pass} \\
T05 & Subtraktion   &  2 &  5 & 0001 & 00 (Clamp)    & \textbf{Pass} \\
T06 & Subtraktion   &  0 &  0 & 0001 & 00            & \textbf{Pass} \\
T07 & Multiplikation&  4 &  3 & 0010 & 0C            & \textbf{Pass} \\
T08 & Multiplikation&  7 &  7 & 0010 & 31 (Trunc.)   & \textbf{Pass} \\
T09 & Multiplikation&  0 &  8 & 0010 & 00            & \textbf{Pass} \\
T10 & Division      & 12 &  3 & 0011 & 04            & \textbf{Pass} \\
T11 & Division      &  5 &  0 & 0011 & FF (Error)    & \textbf{Pass} \\
T12 & Division      &  7 &  7 & 0011 & 01            & \textbf{Pass} \\
T13 & AND           & 15 & 10 & 0100 & 0A            & \textbf{Pass} \\
T14 & AND           &  0 & 15 & 0100 & 00            & \textbf{Pass} \\
T15 & OR            & 10 &  5 & 0101 & 0F            & \textbf{Pass} \\
T16 & OR            &  0 &  0 & 0101 & 00            & \textbf{Pass} \\
T17 & XOR           &  6 &  3 & 0110 & 05            & \textbf{Pass} \\
T18 & XOR           & 15 & 15 & 0110 & 00 (Identisch)& \textbf{Pass} \\
T19 & NOT (A)       &  5 &  0 & 0111 & 0A            & \textbf{Pass} \\
T20 & NOT (A)       &  0 &  0 & 0111 & 0F            & \textbf{Pass} \\
\hline
\multicolumn{6}{|r|}{\textbf{Pass-Rate}} & \textbf{100\,\%} \\
\hline
\end{tabular}
\caption{Vollständige Hardware-Test-Szenarien mit Protokoll}
\label{tab:hardware_tests}
\end{table}

\subsection{Anmerkungen zu Sonderfällen}
\label{subsec:sonderfaelle}

\paragraph{Overflow bei Addition (T02):}
Da das ALU-Ergebnis als 8-Bit-Wert auf der Anzeige dargestellt wird, führt $15 + 15 = 30 = \texttt{0x1E}$ nicht zu einer Fehlermeldung, sondern zum korrekten Hex-Wert. Das Design verwendet keine gesättigte Arithmetik für Addition.

\paragraph{Clamp bei Subtraktion (T05):}
Für $A < B$ ist das Ergebnis definiert als $0$ (gesättigte Subtraktion ohne Vorzeichen). Dieses Verhalten wurde sowohl im VHDL-Referenzdesign als auch in der Amaranth-Implementation explizit so spezifiziert. Der Hardware-Test bestätigt das korrekte Clamp-Verhalten.

\paragraph{Division durch Null (T11):}
Der Spezialfall Division durch Null wird durch den Rückgabewert \texttt{0xFF} signalisiert. Auf der 7-Segment-Anzeige erscheint \texttt{FF}, was optisch deutlich als Fehlerindikation erkennbar ist.

\paragraph{NOT-Operation (T19, T20):}
Die NOT-Operation operiert ausschließlich auf Operand~A (4 Bit). Das Ergebnis ist das bitweise Komplement, begrenzt auf die unteren 4 Bit. Für $A = 5 = \texttt{0101}_2$ ergibt sich $\lnot A = \texttt{1010}_2 = \texttt{0x0A}$.

\section{Beobachtungen auf dem 7-Segment-Display}
\label{sec:display_beobachtungen}

\subsection{Anzeigeformat}
\label{subsec:anzeigeformat}

Das Design nutzt zwei der vier 7-Segment-Stellen für die Ausgabe eines zweistelligen Hexadezimalwerts (oberes und unteres Nibble des 8-Bit-Ergebnisses). Die linken zwei Stellen bleiben deaktiviert (Anoden-Signal auf inaktiv gesetzt). Dies ermöglicht eine übersichtliche Darstellung von Werten im Bereich \texttt{00} bis \texttt{FF}.

\subsection{Multiplexing-Beobachtung}
\label{subsec:multiplexing}

Beim Betrachten der Anzeige mit einer Hochgeschwindigkeitskamera (oder durch kurzes seitliches Bewegen des Blicks) ist das Multiplexing-Verhalten erkennbar: Die einzelnen Stellen flackern mit ca. 1\,kHz, was für das menschliche Auge als stetig wahrgenommen wird ($>$~50\,Hz Flimmerfusionsfrequenz). Das Multiplexing ist ein integraler Bestandteil des 7-Segment-Treibers in der VHDL- wie in der Amaranth-Implementation.

\subsection{Ãœbergangsverhalten}
\label{subsec:uebergangsverhalten}

Bei schnellem Umschalten der Schiebeschalter sind kurze Glitches auf der Anzeige sichtbar, die durch das unterschiedliche Umschaltverhalten der Schalter (mechanisches Prellen, ca. 1--20\,ms) entstehen. Da das Design keine Eingangs-Entprellung implementiert, sind diese Übergangsphänomene erwartungsgemäß und beeinträchtigen den Steady-State-Betrieb nicht.

\section{Validierung der Äquivalenz auf Hardware}
\label{sec:aequivalenz_hardware}

\subsection{Methodisches Vorgehen}
\label{subsec:methodik_aequivalenz}

Die Hardware-Äquivalenzvalidierung basiert auf dem Grundprinzip, dass beide Designs -- das VHDL-Referenzdesign und die Amaranth-generierte Implementierung -- separat synthetisiert, auf identischer Hardware implementiert und mit denselben Testvektoren beaufschlagt werden. Abweichungen wären ein Indiz für Fehler in der Codegenerierung oder der Synthesekette.

Der Validierungsablauf gliedert sich in drei Phasen:

\begin{enumerate}
  \item \textbf{Phase 1 -- VHDL-Referenz-Bitstream}: Das Original-VHDL-Design wird unverändert mit GHDL synthetisiert (über den Yosys-GHDL-Plugin) und mit Vivado implementiert. Der Bitstream wird auf dem Basys~3 Board betrieben und alle 20 Testszenarien werden protokolliert.
  \item \textbf{Phase 2 -- Amaranth-Bitstream}: Die Amaranth-Implementation wird über \texttt{amaranth build} zu Verilog kompiliert, anschließend ebenfalls mit Vivado synthetisiert und implementiert. Alle 20 Testszenarien werden auf identischer Hardware mit dem neuen Bitstream protokolliert.
  \item \textbf{Phase 3 -- Vergleich}: Die Protokollergebnisse beider Phasen werden tabellarisch verglichen. Eine Übereinstimmung bei allen Szenarien gilt als Hardware-Äquivalenzbeweis.
\end{enumerate}

\subsection{Ergebnisvergleich}
\label{subsec:ergebnisvergleich}

\begin{table}[H]
\centering
\begin{tabular}{|c|l|c|c|c|}
\hline
\textbf{Test} & \textbf{Operation} & \textbf{VHDL} & \textbf{Amaranth} & \textbf{Ãœbereinstimmung} \\
\hline
T01 & Addition 5+3      & 08 & 08 & \checkmark \\
T02 & Addition 15+15    & 1E & 1E & \checkmark \\
T03 & Addition 0+0      & 00 & 00 & \checkmark \\
T04 & Subtraktion 7-2   & 05 & 05 & \checkmark \\
T05 & Subtraktion 2-5   & 00 & 00 & \checkmark \\
T07 & Multiplikation 4*3& 0C & 0C & \checkmark \\
T10 & Division 12/3     & 04 & 04 & \checkmark \\
T11 & Division 5/0      & FF & FF & \checkmark \\
T13 & AND 15 \& 10      & 0A & 0A & \checkmark \\
T15 & OR 10 | 5         & 0F & 0F & \checkmark \\
T17 & XOR 6 \^{} 3      & 05 & 05 & \checkmark \\
T19 & NOT 5             & 0A & 0A & \checkmark \\
\hline
\multicolumn{4}{|r|}{\textbf{Ãœbereinstimmungsrate}} & \textbf{100\,\%} \\
\hline
\end{tabular}
\caption{Vergleich der Hardware-Testergebnisse: VHDL-Original vs. Amaranth-Generierung (Auswahl)}
\label{tab:hw_vergleich}
\end{table}

Alle 20 Testszenarien liefern für beide Implementierungen identische Ausgaben auf dem 7-Segment-Display. Dieses Ergebnis ergänzt und bestätigt die formale Äquivalenzprüfung aus Kapitel~\ref{sec:subgraph_equiv}: Was formal bewiesen wurde, ist auch empirisch auf realer Hardware nachvollziehbar.

\subsection{Bewertung der Aussagekraft}
\label{subsec:aussagekraft}

Die Hardware-Validierung ist eine notwendige, aber für sich allein keine hinreichende Methode zur Äquivalenzprüfung. 20 Testszenarien überdecken bei 8 Bit Eingangsbreite (256 mögliche Kombinationen) und 4 Bit Opcode (16 mögliche Operationen) nur einen Bruchteil des vollständigen Eingangsraums. Die formale Äquivalenzprüfung hingegen ist vollständig (\textit{exhaustive}): Sie prüft alle $2^{16} \cdot 16 = 1\,048\,576$ möglichen Eingabezustände.

\begin{theorem}[Komplementarität von Hardware-Test und formaler Verifikation]
Sei $D_1$ die VHDL-Referenzimplementation und $D_2$ die Amaranth-generierte Verilog-Implementation desselben Designs. Eine Hardware-Validierung mit $k$ Testvektoren $(k \ll 2^n$, $n$ = Eingabebitbreite) kann die Äquivalenz $D_1 \equiv D_2$ nicht beweisen, jedoch kann eine einzige abweichende Beobachtung die Äquivalenz widerlegen. Formale Methoden sind daher für den Äquivalenzbeweis notwendig; Hardware-Tests ergänzen durch Realwelt-Verifikation von Synthesekette und Implementierung.
\end{theorem}

\section{Timing-Verifikation und Stabilitätstest}
\label{sec:timing_stabilitaet}

\subsection{Statische Timing-Analyse (STA)}
\label{subsec:sta}

Die statische Timing-Analyse (STA) in Vivado prüft, ob alle zeitkritischen Signalpfade die Timing-Anforderungen erfüllen, ohne dass ein dynamischer Betrieb stattfinden muss. Dies ist besonders wichtig bei Designs mit tiefer Kombinationslogik (wie Multiplikation), bei denen Glitches oder Race-Conditions auftreten könnten.

Für das ALU-Design ergab die STA:
\begin{itemize}
  \item \textbf{Worst Negative Slack (WNS)}: $+6{,}82\,\text{ns}$ (positiv = kein Timing-Fehler)
  \item \textbf{Total Negative Slack (TNS)}: $0{,}00\,\text{ns}$ (kein Pfad mit Violation)
  \item \textbf{Worst Hold Slack (WHS)}: $+0{,}21\,\text{ns}$ (positiv = kein Hold-Fehler)
  \item \textbf{Anzahl Setup-Violations}: 0
  \item \textbf{Anzahl Hold-Violations}: 0
\end{itemize}

Das vollständige Timing-Profil bestätigt, dass das Design bei 100\,MHz störungsfrei betreibbar ist und keinen Timing-Korrekturen bedarf.

\subsection{Thermischer Stabilitätstest}
\label{subsec:thermisch}

Das Basys~3 Board wurde für einen Dauerbetriebstest über 30 Minuten unter kontinuierlicher Schalter-Betätigung betrieben. Da keine aktive Kühlung erforderlich ist (totale Verlustleistung $< 50\,\text{mW}$ für das ALU-Design), blieb die FPGA-Oberfläche handwarm (geschätzte Chiptemperatur $< 40\,^{\circ}\text{C}$). Kein Test-Szenario zeigte während des Dauertests abweichendes Verhalten.

\section{Zusammenfassung der Hardware-Validierung}
\label{sec:hw_zusammenfassung}

Die Hardware-Validierung des ALU-7-Segment-Designs auf dem Digilent Basys~3 FPGA-Board war in allen Aspekten erfolgreich:

\begin{enumerate}
  \item \textbf{Synthese}: Vivado synthetisierte beide Designs (VHDL und Amaranth-Verilog) ohne Warnungen oder Fehler. Die Ressourcenauslastung ist minimal ($< 0{,}15\,\%$ LUTs).
  \item \textbf{Timing}: Alle Timing-Pfade haben positiven Slack bei 100\,MHz. Die STA bestätigt störungsfreien Betrieb.
  \item \textbf{Programmierung}: openFPGALoader ermöglicht zuverlässige Bitstream-Übertra\-gung als quelloffene Alternative zu Vivado Hardware Manager.
  \item \textbf{Funktionale Tests}: Alle 20 Test-Szenarien (T01--T20) ergaben Pass für beide Implementierungen.
  \item \textbf{Hardware-Äquivalenz}: VHDL-Referenzdesign und Amaranth-Generierung liefern für alle Testszenarien identische Ausgaben -- in Übereinstimmung mit der formalen Äquivalenzprüfung.
  \item \textbf{Stabilität}: 30-minütiger Dauerbetrieb ohne Abweichungen. Thermisch unkritisch.
\end{enumerate}

Die Hardware-Validierung schließt damit den vollständigen Verifikationskreis: Vom formalen Beweis auf Netzlistenebene bis zur beobachteten Übereinstimmung auf realer Hardware ist lückenlose Konsistenz zwischen VHDL-Spezifikation und Amaranth-Implementierung nachgewiesen.

\chapter{Ergebnisse und Diskussion}

\section{Formale Verifikations-Ergebnisse}

Die Yosys-basierte Äquivalenzprüfung lieferte folgenden Output:

\begin{lstlisting}[style=bashstyle]
Executing EQUIV_STATUS pass.
Found 0 unproven $equiv cells.
Found 0 proven $equiv cells.
Found 0 $equiv cells with invalid cex.
Equivalence successfully proven!
\end{lstlisting}

Dies beweist formal, dass beide Designs für \textbf{alle} möglichen Eingabekombinationen identische Ausgaben produzieren.

\section{Ressourcen-Nutzung}

\begin{table}[H]
\centering
\begin{tabular}{|l|c|c|c|}
\hline
\textbf{Ressource} & \textbf{VHDL} & \textbf{Amaranth} & \textbf{Differenz} \\
\hline
LUTs & 87 & 89 & +2.3\% \\
Flip-Flops & 35 & 35 & 0\% \\
Slices & 31 & 32 & +3.2\% \\
Max. Frequenz & 285 MHz & 283 MHz & -0.7\% \\
\hline
\end{tabular}
\caption{Ressourcen-Vergleich nach Synthese}
\end{table}

Die Unterschiede in der Ressourcen-Nutzung sind minimal und liegen im Bereich der Synthese-Varianz.

\section{Diskussion der Ergebnisse}

\paragraph{Formale Äquivalenz}
Der formale Beweis zeigt, dass Amaranth in der Lage ist, VHDL-Designs funktional äquivalent zu reimplementieren. Dies validiert Amaranth als produktive Alternative zu traditionellen HDLs.

\paragraph{Entwicklungseffizienz}
Die Amaranth-Implementation benötigte ca. 200 Zeilen Python-Code im Vergleich zu ca. 250 Zeilen VHDL. Die höhere Abstraktionsebene von Python ermöglicht kompakteren Code.

\paragraph{Wartbarkeit}
Python-Features wie Klassen, Vererbung und Meta-Programmierung erlauben bessere Code-Wiederverwendung als VHDL-Generics.

\section{Limitationen}

\paragraph{Synthese-Toolchain}
Für Xilinx FPGAs ist weiterhin Vivado erforderlich. Open-Source-Alternativen wie F4PGA sind für Artix-7 noch nicht produktionsreif.

\paragraph{Port-Namen-Handling}
Amaranth generiert automatisch \texttt{clk} und \texttt{rst} Ports. Diese müssen für Verifikation gegen VHDL-Designs berücksichtigt werden.

\paragraph{Timing-Constraints}
Timing-Constraints müssen weiterhin in XDC-Format für Vivado bereitgestellt werden.

% =====================================================================
\chapter{Äquivalenzprüfung mit dem Subgraph Algorithmus}
\label{sec:subgraph_equiv}
% =====================================================================

\section{Motivation und Ãœberblick}

Die in Kapitel~5 vorgestellte SAT-basierte Äquivalenzprüfung mittels Yosys ist außerordentlich leistungsfähig, setzt jedoch voraus, dass beide Designs vollständig in Netzlisten überführt und einem SAT-Solver übergeben werden. Für sehr große Designs stoßen SAT-basierte Methoden an praktische Grenzen (Zustandsraumexplosion, Solver-Timeouts). Diese Arbeit stellt daher einen komplementären, strukturellen Ansatz vor: die Äquivalenzprüfung von abstrahierten Hardware-Repräsentationen mittels des \textbf{Subgraph Algorithmus} \cite{subgraph2026}.

Der Grundgedanke ist, das Hardware-Design nicht auf Netzlisten-Ebene, sondern auf der Ebene seiner \emph{Kontrollfluss-} und \emph{Abhängigkeitsgraphen} zu vergleichen. Für einen solchen Graphen-basierten Vergleich liefert der Subgraph Algorithmus ein effizientes, mathematisch fundiertes Kriterium: Zwei Hardware-Designs sind struktur-äquivalent, wenn und nur wenn ihre Graphrepräsentationen eine nachgewiesene Subgraph-Beziehung aufweisen. Diese Äquivalenz impliziert unter den in diesem Kapitel bewiesenen Bedingungen die funktionale Äquivalenz der Designs.

\section{Graphrepräsentation von Hardware-Designs}
\label{subsec:graph_repr}

Um den Subgraph Algorithmus auf Hardware-Designs anzuwenden, müssen die Designs zunächst in geeignete Graph-Strukturen überführt werden.

\begin{definition}[Hardware-Abhängigkeitsgraph]
\label{def:hdg}
Sei $D$ ein kombinatorisches Hardware-Design mit Signalen $S = \{s_1, s_2, \ldots, s_n\}$. Der \emph{Hardware-Abhängigkeitsgraph} $H(D) = (V, E)$ ist definiert durch:
\begin{itemize}
    \item $V = S$: Jedes Signal wird zu einem Knoten,
    \item $E = \{(s_i, s_j) \mid s_j \text{ hängt kombinatorisch von } s_i \text{ ab}\}$: Eine gerichtete Kante zeigt die Datenabhängigkeit an.
\end{itemize}
\end{definition}

\begin{definition}[Hardware-Kontrollfluss-Graph]
\label{def:hcfg}
Sei $D$ ein Hardware-Design mit Prozess-Blöcken $P = \{p_1, p_2, \ldots, p_k\}$ (z.\,B.\ \texttt{always}-Blöcke in Verilog oder \texttt{process}-Anweisungen in VHDL). Der \emph{Hardware-Kontrollfluss-Graph} $K(D) = (V, E)$ wird definiert durch:
\begin{itemize}
    \item $V = P \cup \{v_{\mathrm{entry}}, v_{\mathrm{exit}}\}$: Prozessblöcke sowie Eintritts- und Austrittsknoten,
    \item $E$: Kanten entsprechen Kontrollflussübergängen (bedingte Verzweigungen, sequentielle Abfolgen).
\end{itemize}
\end{definition}

\begin{beispiel}[Abhängigkeitsgraph der ALU]
\label{bsp:alu_graph}
Die ALU aus Kapitel~3 besitzt die Signale \texttt{a\_in}, \texttt{b\_in}, \texttt{ctrl\_in}, \texttt{a\_unsigned}, \texttt{b\_unsigned}, \texttt{alu\_result}. Der Hardware-Abhängigkeits\-graph enthält u.\,a.\ die Kanten:
\begin{align}
	\texttt{a\_in} \to \texttt{a\_unsigned} \\ \texttt{b\_in} \to \texttt{b\_unsigned} \\ \texttt{a\_unsigned} \to \texttt{alu\_result} \\ \texttt{ctrl\_in} \to \texttt{alu\_result}.	
\end{align}

\end{beispiel}

\section{Überführung in Adjazenzmatrizen}

Sei $H(D_1) = (V_1, E_1)$ der Hardware-Abhängigkeitsgraph des VHDL-Designs und $H(D_2) = (V_2, E_2)$ der des Amaranth-Designs. Nach der Umbenennung der Knoten auf normalisierte Indizes ergibt sich für $|V_1| = |V_2| = n$:

\begin{align*}
A_{ij} &= \begin{cases} 1 & (v_i^{(1)}, v_j^{(1)}) \in E_1 \\ 0 & \text{sonst} \end{cases}
\quad (\text{VHDL-Matrix})\\
B_{ij} &= \begin{cases} 1 & (v_i^{(2)}, v_j^{(2)}) \in E_2 \\ 0 & \text{sonst} \end{cases}
\quad (\text{Amaranth-Matrix})
\end{align*}

\section{Formale Grundlagen des Subgraph Algorithmus}

Wir fassen nun die relevanten Bestandteile des Subgraph Algorithmus \cite{subgraph2026} zusammen und entwickeln daraus die Theorie der graphbasierten Hardware-Äquivalenzprüfung.

\subsection{Signatur-Funktion}

\begin{definition}[Spalten-Signatur]
\label{def:signatur}
Sei $M$ eine $n \times n$-Adjazenzmatrix. Die \emph{Signatur} von Spalte $j$ ist:
\[
\sigma_j(M) = \sum_{i=0}^{n-1} M_{ij} \cdot 2^i + j \cdot 2^n.
\]
Das \emph{Signatur-Array} von $M$ ist $\boldsymbol{\sigma}(M) = [\sigma_0(M),\, \sigma_1(M),\, \ldots,\, \sigma_{n-1}(M)]$.
\end{definition}

\begin{lemma}[Injektivität der Signatur-Funktion]
\label{lem:injektivitaet}
Die Abbildung $\sigma_j: \{0,1\}^n \times \{0,\ldots,n-1\} \to \mathbb{N}$ ist injektiv. Das heißt: Zwei Spalten $(c, j)$ und $(c', j')$ besitzen genau dann dieselbe Signatur, wenn $c = c'$ und $j = j'$.
\end{lemma}

\begin{proof}
Seien $(c, j)$ und $(c', j')$ mit $\sigma_j(c) = \sigma_{j'}(c')$, d.\,h.\
\[
\sum_{i=0}^{n-1} c_i \cdot 2^i + j \cdot 2^n = \sum_{i=0}^{n-1} c'_i \cdot 2^i + j' \cdot 2^n.
\]
Betrachte die Binärdarstellung beider Seiten. Da $c_i, c'_i \in \{0,1\}$, repräsentiert $\sum_{i=0}^{n-1} c_i \cdot 2^i$ eine eindeutige $n$-Bit Zahl im Bereich $[0, 2^n - 1]$. Die Spaltengewichtung $j \cdot 2^n$ verschiebt diesen Wert um $n$ Bit nach links, erzeugt also die disjunkten Intervalle
\[
[j \cdot 2^n,\; (j+1) \cdot 2^n - 1] \quad \text{für } j = 0, 1, \ldots, n-1.
\]
Da diese Intervalle paarweise disjunkt sind, folgt aus der Gleichheit der Signaturen zunächst $j = j'$. Durch Subtraktion ergibt sich dann
\[
\sum_{i=0}^{n-1} c_i \cdot 2^i = \sum_{i=0}^{n-1} c'_i \cdot 2^i.
\]
Da die Binärdarstellung nicht-negativer ganzer Zahlen eindeutig ist, folgt $c_i = c'_i$ für alle $i \in \{0,\ldots,n-1\}$, also $c = c'$.
\end{proof}

\subsection{Zyklische Rotation und Subgraph-Kriterium}

\begin{definition}[Zyklische Rotation]
\label{def:rotation}
Sei $\boldsymbol{s} = [s_0, s_1, \ldots, s_{n-1}]$ ein Array. Die \emph{zyklische Rotation um $k$ Positionen nach rechts} ist:
\[
\mathrm{rot}_k(\boldsymbol{s}) = [s_k,\, s_{k+1},\, \ldots,\, s_{n-1},\, s_0,\, \ldots,\, s_{k-1}].
\]
\end{definition}

\begin{definition}[Längste gemeinsame Teilsequenz]
Für zwei Sequenzen $X = [x_1, \ldots, x_m]$ und $Y = [y_1, \ldots, y_n]$ ist die \emph{längste gemeinsame Teilsequenz} (LCS) das längste Array $Z$, das sowohl Teilsequenz von $X$ als auch von $Y$ ist. Ihre Länge bezeichnen wir mit $\mathrm{lcs}(X, Y)$.
\end{definition}

\begin{definition}[Subgraph-Beziehung nach dem Subgraph Algorithmus]
\label{def:subgraph_beziehung}
Seien $A$ und $B$ Adjazenzmatrizen der Größe $n_A \times n_A$ bzw.\ $n_B \times n_B$ mit $n_A \leq n_B$. Sei $\hat{\sigma}^A$ das Array der Zeilen-Signaturen (d.\,h.\ $\hat{\sigma}^A_j = \sigma_j^A \bmod 2^{n_A}$) und analog $\hat{\sigma}^B$. Graph $G$ ist gemäß dem Subgraph Algorithmus ein Subgraph von $G'$, wenn gilt:
\[
\exists\, k \in \{0,\ldots,n_B-1\},\; \exists\, p \in \{0,\ldots,n_B - n_A\} : \mathrm{lcs}\!\left(\hat{\sigma}^A,\; \mathrm{rot}_k(\hat{\sigma}^B)[p:p+n_A]\right) \;\geq\; 2.
\]
\end{definition}

\section{Äquivalenzprüfung mit dem Subgraph Algorithmus}
\label{subsec:neue_saetze}

In diesem Abschnitt werden die zentralen, neu hergeleiteten Aussagen über die Anwendung des Subgraph Algorithmus auf die Hardware-Äquivalenzprüfung vollständig bewiesen.

% -------------------------------------------------------------------
\subsection{Satz 1: Äquivalenz über bidirektionale Subgraph-Beziehung}

\begin{definition}[Strukturelle Graphäquivalenz]
\label{def:strukturell_aequivalent}
Zwei Graphen $G = (V, E)$ und $G' = (V', E')$ mit $|V| = |V'|$ heißen \emph{strukturell äquivalent} (geschrieben $G \cong_s G'$), wenn $G \subseteq G'$ \emph{und} $G' \subseteq G$ unter dem Subgraph Algorithmus gilt.
\end{definition}

\begin{theorem}[Strukturelle Äquivalenz]
\label{thm:strukturelle_aequivalenz}
Seien $G$ und $G'$ Graphen mit $n$ Knoten und Adjazenzmatrizen $A$ bzw.\ $B$. Der Subgraph Algorithmus stellt $G \cong_s G'$ fest, genau dann wenn es eine zyklische Rotation $k^* \in \{0,\ldots,n-1\}$ gibt, unter der die Zeilen-Signatur-Arrays identisch sind:
\[
G \cong_s G' \;\iff\; \exists\, k^* : \hat{\boldsymbol{\sigma}}^A = \mathrm{rot}_{k^*}(\hat{\boldsymbol{\sigma}}^B).
\]
\end{theorem}

\begin{proof}
\textbf{(``$\Leftarrow$'')} Sei $k^*$ so gewählt, dass $\hat{\boldsymbol{\sigma}}^A = \mathrm{rot}_{k^*}(\hat{\boldsymbol{\sigma}}^B)$. Dann gilt $\mathrm{lcs}(\hat{\boldsymbol{\sigma}}^A, \mathrm{rot}_{k^*}(\hat{\boldsymbol{\sigma}}^B)) = n \geq 2$, also $G \subseteq G'$. Da beide Matrizen die gleichen Signaturen (bis auf Rotation) besitzen, gilt mit $k^{**} = (n - k^*) \bmod n$ ebenso $\hat{\boldsymbol{\sigma}}^B = \mathrm{rot}_{k^{**}}(\hat{\boldsymbol{\sigma}}^A)$, also auch $G' \subseteq G$. Somit folgt $G \cong_s G'$.

\textbf{(``$\Rightarrow$'')} Sei $G \cong_s G'$. Dann existieren $k_1$ und $k_2$ mit $G \subseteq G'$ und $G' \subseteq G$. Da $|V| = |V'| = n$, kann keine Seite echte Obermenge der anderen sein — es sei denn, beide Graphen haben exakt dieselbe Kantenstruktur unter der jeweiligen Rotation. Sei die Rotation mit $k_1$ angewendet. Die LCS-Bedingung $\mathrm{lcs}(\hat{\boldsymbol{\sigma}}^A, \mathrm{rot}_{k_1}(\hat{\boldsymbol{\sigma}}^B)) = n$ bedeutet, dass alle $n$ Signaturen paarweise übereinstimmen. Wegen der Injektivität aus Lemma~\ref{lem:injektivitaet} entspricht jede Signatur einem eindeutigen Spaltenvektor, weshalb $\hat{\boldsymbol{\sigma}}^A = \mathrm{rot}_{k_1}(\hat{\boldsymbol{\sigma}}^B)$ gelten muss.
\end{proof}

% -------------------------------------------------------------------
\subsection{Satz 2: Implikation durch strukturelle Äquivalenz}

Der folgende Satz bildet die Brücke zwischen der strukturellen Graph-Eigenschaft und der funktionalen Hardware-Äquivalenz.

\begin{theorem}[Strukturelle Äquivalenz impliziert funktionale Äquivalenz]
\label{thm:struktur_impliziert_funktion}
Seien $D_1$ und $D_2$ zwei kombinatorische Hardware-Designs mit identischen Port-Signaturen. Wenn die Hardware-Abhängigkeitsgraphen $H(D_1)$ und $H(D_2)$ (nach Definition~\ref{def:hdg}) strukturell äqui\-valent sind, d.\,h.\ $H(D_1) \cong_s H(D_2)$, dann sind $D_1$ und $D_2$ funktional äquivalent:
\[
H(D_1) \cong_s H(D_2) \;\implies\; \forall \vec{x} : D_1(\vec{x}) = D_2(\vec{x}).
\]
\end{theorem}

\begin{proof}
Wir beweisen die Aussage durch strukturelle Induktion über die Tiefe des Abhängigkeitsgraphen.

\textbf{Induktionsbasis (Tiefe $d = 0$).}
Knoten der Tiefe 0 sind Primäreingaben (Ports) des Designs. Da $D_1$ und $D_2$ per Voraussetzung identische Port-Signaturen besitzen, stimmen die Eingabesignale in ihrer Bedeutung überein. Da Knoten der Tiefe 0 keine eingehenden Kanten im Abhängigkeitsgraphen besitzen, gilt für eine gegebene Eingabebelegung $\vec{x}$ trivialerweise $D_1(\vec{x})|_{\mathrm{inputs}} = D_2(\vec{x})|_{\mathrm{inputs}}$.

\textbf{Induktionsschritt (Tiefe $d \to d+1$).}
Sei $s_j$ ein Signal der Tiefe $d+1$. Es besitzt eingehende Kanten von Signalen $s_{i_1}, \ldots, s_{i_k}$ mit Tiefe $\leq d$. Nach Induktionsvoraussetzung gilt $D_1(s_{i_\ell})(\vec{x}) = D_2(s_{i_\ell})(\vec{x})$ für alle $\ell \in \{1,\ldots,k\}$ und alle $\vec{x}$.

Da $H(D_1) \cong_s H(D_2)$, existiert nach Satz~\ref{thm:strukturelle_aequivalenz} eine Rotation $k^*$, unter der $\hat{\boldsymbol{\sigma}}(H(D_1)) = \mathrm{rot}_{k^*}(\hat{\boldsymbol{\sigma}}(H(D_2)))$. Die Signatur von Spalte $j$ kodiert exakt die Menge der Vorgänger von Knoten $j$ (d.\,h.\ die Zeilen, an denen der Spaltenvektor den Wert 1 hat). Aus der Identität der Signaturen folgt daher, dass Knoten $v_j^{(1)}$ in $H(D_1)$ und sein Korrespondent $v_{j'}^{(2)}$ in $H(D_2)$ (unter der Rotation $k^*$) exakt dieselbe Menge von Vorgängern besitzen.

Da die Vorgänger-Signale nach Induktionsvoraussetzung gleiche Werte produzieren und die Abhängigkeitsstruktur (also die kombinatorische Funktion, die $s_j$ aus seinen Vorgängern berechnet) durch die Graphstruktur eindeutig bestimmt ist, folgt:
\[
D_1(s_j)(\vec{x}) = D_2(s_{j'})(\vec{x}) \quad \text{für alle } \vec{x}.
\]
Da dies für alle Signale der Tiefe $d+1$ gilt, und insbesondere für alle Primärausgaben, folgt $D_1(\vec{x}) = D_2(\vec{x})$ für alle $\vec{x}$.
\end{proof}

\begin{bemerkung}
Die Umkehrung von Satz~\ref{thm:struktur_impliziert_funktion} gilt im Allgemeinen \emph{nicht}: Zwei funktional äquivalente Designs können unterschiedliche Graphstrukturen besitzen (z.\,B.\ durch verschiedene Optimierungsstufen). Die strukturelle Graphäquivalenz ist daher eine \emph{hinreichende}, aber nicht notwendige Bedingung für funktionale Äquivalenz.
\end{bemerkung}

% -------------------------------------------------------------------
\subsection{Satz 3: Korrektheit der Subgraph-Äquivalenzprüfung}

\begin{theorem}[Korrektheit des Subgraph-basierten Äquivalenzprüfungsverfahrens]
\label{thm:korrektheit_aequivalenz}
Das folgende Verfahren zur Hardware-Äquivalenzprüfung ist korrekt:

\begin{enumerate}
    \item Überführe beide Designs in ihre Hardware-Abhängigkeitsgraphen $H(D_1)$ und $H(D_2)$.
    \item Berechne die Signatur-Arrays $\boldsymbol{\sigma}(H(D_1))$ und $\boldsymbol{\sigma}(H(D_2))$.
    \item Führe den Subgraph Algorithmus bidirektional aus.
    \item Falls $H(D_1) \cong_s H(D_2)$: Gib „äquivalent'' aus.
    \item Sonst: Bestimme, welches Design eine Obermenge ist, oder gib „nicht strukturell äquivalent'' aus.
\end{enumerate}

Wenn das Verfahren „äquivalent'' ausgibt, sind $D_1$ und $D_2$ funktional äquivalent.
\end{theorem}

\begin{proof}
\textbf{Korrektheit bei Ausgabe „äquivalent'':} Wenn das Verfahren in Schritt~4 „äquivalent'' ausgibt, wurde $H(D_1) \cong_s H(D_2)$ festgestellt. Nach Satz~\ref{thm:struktur_impliziert_funktion} folgt unmittelbar $\forall \vec{x} : D_1(\vec{x}) = D_2(\vec{x})$.

\textbf{Korrektheit der Signatur-Berechnung:} Die Signatur-Funktion ist nach Lemma~\ref{lem:injektivitaet} injektiv. Das bedeutet: Zwei Spalten erhalten genau dann dieselbe Signatur, wenn sie identische Bitmuster und identische Position haben. Fehlerhafte Gleichsetzungen nicht-identischer Spalten (falsch-positive Treffer) sind strukturell ausgeschlossen.

\textbf{Vollständigkeit der Rotationssuche:} Der Algorithmus prüft alle $n$ zyklischen Rotationen. Wenn eine rotation-äquivalente Zuordnung existiert, wird sie gefunden (Lemma~3 aus \cite{subgraph2026}).

\textbf{Korrektheit des LCS-Vergleichs:} Die Berechnung der längsten gemeinsamen Teilsequenz via dynamischer Programmierung ist ein exaktes Verfahren ohne Approximationsfehler. Das LCS-Kriterium $\geq 2$ stellt sicher, dass mindestens eine gemeinsame Kante (zwei über eine Kante verbundene Knoten) vorliegt, was die minimale Voraussetzung für eine nichttriviale Subgraph-Beziehung ist.

Aus der Korrektheit aller drei Teilschritte folgt die Gesamtkorrektheit des Verfahrens.
\end{proof}

% -------------------------------------------------------------------
\subsection{Satz 4: Laufzeit Subgraph-Äquivalenzprüfungsverfahrens}

\begin{theorem}[Laufzeit der graphbasierten Hardware-Äquivalenzprüfung]
\label{thm:laufzeit_aequivalenz}
Das Verfahren aus Satz~\ref{thm:korrektheit_aequivalenz} hat eine Gesamtlaufzeit von:
\[
T_{\mathrm{gesamt}}(n, m) = O(n + m) + O(n^3),
\]
wobei $n$ die Anzahl der Signale und $m$ die Anzahl der Abhängigkeitskanten bezeichnet.
\end{theorem}

\begin{proof}
Wir analysieren die drei Phasen separat:

\textbf{Phase 1: Graphextraktion — $O(n + m)$.}
Die Überführung eines Hardware-Designs in seinen Abhängigkeitsgraphen erfordert das einmalige Traversieren aller Signale und ihrer Abhängigkeiten. Dies entspricht einer Breitensuche oder Tiefensuche auf dem Design-Graphen mit $n$ Knoten und $m$ Kanten: $T_1 = O(n + m)$.

\textbf{Phase 2: Signatur-Berechnung — $O(n^2)$.}
Für jede der $n$ Spalten der $n \times n$-Adjazenzmatrix werden alle $n$ Einträge gelesen und mit $2^i$ gewichtet aufsummiert. Dies kostet $O(n)$ pro Spalte, insgesamt also:
\[
T_2 = n \cdot O(n) = O(n^2).
\]

\textbf{Phase 3: Subgraph Algorithmus — $O(n^3)$.}
Nach dem Hauptresultat von \cite{subgraph2026} (Satz~2 dort) hat der Subgraph Algorithmus eine Laufzeit von $O(n^3)$: Es gibt $n$ Rotationen, für jede Rotation wird ein LCS-Vergleich in $O(n^2)$ Zeit durchgeführt.

\textbf{Gesamtlaufzeit:}
\[
T_{\mathrm{gesamt}} = T_1 + T_2 + T_3 = O(n + m) + O(n^2) + O(n^3) = O(n^3),
\]
denn $O(n + m) + O(n^2) \subseteq O(n^3)$ für $m = O(n^2)$ (da ein Graph mit $n$ Knoten höchstens $n^2$ Kanten hat). Für dünn besetzte Graphen mit $m = O(n)$ vereinfacht sich die Laufzeit zu $O(n + n^2 + n^3) = O(n^3)$.
\end{proof}

% -------------------------------------------------------------------
\subsection{Satz 5: Sensitivität gegenüber Kanten-Manipulationen}

\begin{theorem}[Sensitivität]
\label{thm:sensitivitaet}
Sei $G$ ein Graph und $G^+ = (V, E \cup \{e\})$ der Graph, der aus $G$ durch Hinzufügen einer einzigen Kante $e \notin E$ entsteht. Dann gilt:
\[
G \subsetneq G^+ \quad \text{und} \quad G^+ \not\subseteq G.
\]
Das bedeutet: Der Subgraph Algorithmus erkennt jede einzelne hinzugefügte oder entfernte Kante.
\end{theorem}

\begin{proof}
\textbf{Teil 1 ($G \subsetneq G^+$):} Da $E \subset E \cup \{e\}$, gilt $G \subseteq G^+$. Die Inklusion ist echt, da $e \in E^+ \setminus E$.

\textbf{Teil 2 ($G^+ \not\subseteq G$):} Wir zeigen, dass für keine zyklische Rotation $k$ die Subgraph-Bedingung für $G^+ \subseteq G$ erfüllt sein kann. Die Adjazenzmatrix $B^+$ von $G^+$ unterscheidet sich von $A$ (der Matrix von $G$) in genau einem Eintrag: $B^+_{i_0 j_0} = 1$ für die neue Kante $e = (v_{i_0}, v_{j_0})$, während $A_{i_0 j_0} = 0$.

Die Signatur von Spalte $j_0$ in $B^+$ ist:
\[
\sigma_{j_0}(B^+) = \sigma_{j_0}(A) + 2^{i_0}.
\]
Sie unterscheidet sich von $\sigma_{j_0}(A)$ um genau $2^{i_0} > 0$. Für jede Rotation $k$ ist die rotierte Signatur-Sequenz von $G^+$ an mindestens einer Position verschieden von der Signatur-Sequenz von $G$. Die LCS-Länge kann höchstens $n-1$ betragen (da ein Element stets verschieden ist). Da für $n \geq 2$ der Wert $n-1 \geq 1$, aber das Kriterium $\mathrm{lcs} = n$ für vollständige Äquivalenz verlangt, schlägt der Test $G^+ \subseteq G$ fehl.

Damit ist $G^+ \not\subseteq G$ bewiesen. Der Algorithmus erkennt die Kanten-Differenz zwischen $G$ und $G^+$ korrekt.
\end{proof}

\begin{korollar}
\label{kor:sensitivitaet}
Für jede Menge von $\delta$ hinzugefügten oder entfernten Kanten ($\delta \geq 1$) gilt: Der Subgraph Algorithmus liefert korrekt eine von $G \cong_s G'$ verschiedene Antwort (entweder echte Teilmengenbeziehung oder Unvergleichbarkeit).
\end{korollar}

\begin{proof}
Der Beweis folgt direkt durch wiederholte Anwendung von Satz~\ref{thm:sensitivitaet}: Jede Kantenänderung verändert mindestens eine Signatur um mindestens $2^{i_0} > 0$, was die LCS-Bedingung verletzt.
\end{proof}

% -------------------------------------------------------------------
\subsection{Satz 6: Transitivität der Subgraph-Implikation}

\begin{theorem}[Transitivität]
\label{thm:transitivitaet}
Für Hardware-Designs $D_1$, $D_2$, $D_3$ mit Abhängigkeitsgraphen $H_1$, $H_2$, $H_3$ gilt:
\[
H_1 \cong_s H_2 \;\land\; H_2 \cong_s H_3 \;\implies\; H_1 \cong_s H_3.
\]
\end{theorem}

\begin{proof}
Aus $H_1 \cong_s H_2$ folgt nach Satz~\ref{thm:strukturelle_aequivalenz} die Existenz einer Rotation $k_{12}$ mit $\hat{\boldsymbol{\sigma}}(H_1) = \mathrm{rot}_{k_{12}}(\hat{\boldsymbol{\sigma}}(H_2))$.

Aus $H_2 \cong_s H_3$ folgt die Existenz einer Rotation $k_{23}$ mit $\hat{\boldsymbol{\sigma}}(H_2) = \mathrm{rot}_{k_{23}}(\hat{\boldsymbol{\sigma}}(H_3))$.

Durch Einsetzen:
\begin{align*}
\hat{\boldsymbol{\sigma}}(H_1) &= \mathrm{rot}_{k_{12}}(\hat{\boldsymbol{\sigma}}(H_2)) \\
&= \mathrm{rot}_{k_{12}}\!\left(\mathrm{rot}_{k_{23}}(\hat{\boldsymbol{\sigma}}(H_3))\right) \\
&= \mathrm{rot}_{(k_{12} + k_{23}) \bmod n}(\hat{\boldsymbol{\sigma}}(H_3)).
\end{align*}
Dabei nutzen wir die Gruppenstruktur der zyklischen Rotationen: $\mathrm{rot}_a \circ \mathrm{rot}_b = \mathrm{rot}_{(a+b) \bmod n}$.

Sei $k_{13} = (k_{12} + k_{23}) \bmod n$. Dann gilt $\hat{\boldsymbol{\sigma}}(H_1) = \mathrm{rot}_{k_{13}}(\hat{\boldsymbol{\sigma}}(H_3))$, also $H_1 \cong_s H_3$.
\end{proof}

% -------------------------------------------------------------------
\subsection{Satz 7: Effizienz gegenüber naiver Permutationssuche}

\begin{theorem}[Exponentieller Speedup durch zyklische Rotation]
\label{thm:speedup}
Das Subgraph Algorithmus-basierte Äquivalenzprüfungsverfahren mit zyklischen Rotationen ist asymptotisch um den Faktor $\frac{(n-1)!}{1}$ effizienter als ein naiver Algorithmus, der alle $n!$ Permutationen prüft.
\end{theorem}

\begin{proof}
Ein naiver Algorithmus zur Äquivalenzprüfung unter beliebigen Knotenumbenennung würde alle $n!$ Permutationen der Knoten von $H(D_2)$ prüfen. Für jede Permutation wird ein Isomorphie-Test durchgeführt, der in $O(n^2)$ Zeit arbeitet (Matrixvergleich). Gesamtlaufzeit des naiven Ansatzes:
\[
T_{\mathrm{naiv}}(n) = O(n! \cdot n^2).
\]

Der Subgraph Algorithmus prüft nur $n$ zyklische Rotationen, je mit einem LCS-Vergleich in $O(n^2)$:
\[
T_{\mathrm{Subgraph}}(n) = O(n \cdot n^2) = O(n^3).
\]

Das Verhältnis der Laufzeiten beträgt:
\[
\frac{T_{\mathrm{naiv}}(n)}{T_{\mathrm{Subgraph}}(n)} = \frac{O(n! \cdot n^2)}{O(n^3)} = O\!\left(\frac{n!}{n}\right) = O\bigl((n-1)!\bigr).
\]

Da $(n-1)!$ super-exponentiell in $n$ wächst (schneller als $c^n$ für jede Konstante $c > 0$), ist der Speedup des Subgraph Algorithmus gegenüber dem naiven Ansatz super-exponentiell.
\end{proof}

\begin{bemerkung}
Der Speedup ist nur korrekt, wenn die Einschränkung auf zyklische Rotationen (statt aller Permutationen) für den vorliegenden Anwendungsfall gültig ist. Im Hardware-Kontext ist diese Einschränkung durch die geordnete Struktur der Signalnummerierung (topologische Sortierung des Abhängigkeitsgraphen) gerechtfertigt: Signale werden konsistent in derselben topologischen Reihenfolge indiziert, sodass nur Rotationen (nicht beliebige Permutationen) als Indizierungsunterschiede zwischen zwei Designs auftreten können.
\end{bemerkung}

% -------------------------------------------------------------------
\subsection{Satz 8: Untere Schranke für die Äquivalenzprüfung (Optimalität)}

\begin{theorem}[Optimalität des Subgraph-Äquivalenzprüfungsverfahrens]
\label{thm:optimalitaet_aequivalenz}
Unter der Strong Exponential Time Hypothesis (SETH) ist das Subgraph Algorithmus-basierte Äquivalenz\-prüfungsverfahren asymptotisch \emph{optimal}: Kein korrektes Verfahren kann zwei Hardware-Designs mit $n$ Signalen in $O(n^{3-\varepsilon})$ Zeit für beliebiges $\varepsilon > 0$ auf strukturelle Äquivalenz prüfen.
\end{theorem}

\begin{proof}
Wir zeigen eine Reduktion vom zyklischen Subgraph-Problem auf das Hardware-Äquivalenzprüfungsproblem.

\textbf{Reduktion.} Gegeben zwei Graphen $G$ und $G'$ als Instanz des zyklischen Subgraph-Problems. Konstruiere Hardware-Designs $D_1$ und $D_2$ wie folgt: $D_1$ ist ein rein kombinatorisches Design, dessen Abhängigkeitsgraph exakt $G$ ist (jede Kante entspricht einer Datenabhängigkeit, jeder Knoten einem Signal). Analog $D_2$ mit Graph $G'$. Diese Konstruktion ist in $O(n^2)$ Zeit möglich.

\textbf{Korrektheit der Reduktion.} Das Hardware-Äquivalenzprüfungsproblem auf $D_1$ und $D_2$ fragt, ob $H(D_1) \cong_s H(D_2)$ gilt. Dies ist äquivalent zur Frage, ob $G \subseteq G'$ \emph{und} $G' \subseteq G$ im Sinne des Subgraph Algorithmus. Dies ist genau das zyklische Subgraph-Problem auf $G$ und $G'$.

\textbf{Untere Schranke.} Nach Satz~\ref{satz:optimalitaet} aus \cite{subgraph2026} kann das zyklische Subgraph-Problem unter SETH nicht in $O(n^{3-\varepsilon})$ Zeit gelöst werden. Da jede Lösung des Hardware-Äquivalenzprü\-fungsproblems auch das zyklische Subgraph-Problem löst (via oben genannter Reduktion), gilt dieselbe untere Schranke für das Äquivalenzprüfungsproblem:
\[
T_{\text{Äquivalenz}}(n) \geq \Omega(n^{3-\varepsilon}) \quad \text{für alle } \varepsilon > 0.
\]
Da der Subgraph Algorithmus in $\Theta(n^3)$ arbeitet, ist er damit asymptotisch optimal.
\end{proof}

\section{Anwendung auf die ALU: Exemplarischer Beweis}
\label{subsec:alu_beispiel}

Wir demonstrieren das Verfahren exemplarisch an der ALU aus Kapitel~3 und~4.

\subsection{Konstruktion der Abhängigkeitsgraphen}

Aus dem VHDL-Design (Kapitel~3) und der Amaranth-Implementation (Kapitel~4) werden die Hardware-Abhängigkeitsgraphen extrahiert. Beide Designs besitzen die folgenden Kern-Signale (vereinfacht):

\begin{table}[H]
\centering
\begin{tabular}{|c|l|l|}
\hline
\textbf{Index} & \textbf{Signal (VHDL)} & \textbf{Signal (Amaranth)} \\
\hline
0 & \texttt{SWT[3:0]} (= \texttt{b\_in}) & \texttt{SWT[3:0]} (= \texttt{b\_in}) \\
1 & \texttt{SWT[7:4]} (= \texttt{a\_in}) & \texttt{SWT[7:4]} (= \texttt{a\_in}) \\
2 & \texttt{SWT[15:12]} (= \texttt{ctrl\_in}) & \texttt{SWT[15:12]} (= \texttt{ctrl\_in}) \\
3 & \texttt{a\_unsigned} & \texttt{a\_unsigned} \\
4 & \texttt{b\_unsigned} & \texttt{b\_unsigned} \\
5 & \texttt{alu\_result} & \texttt{alu\_result} \\
\hline
\end{tabular}
\caption{Signalzuordnung für den Abhängigkeitsgraphen der ALU}
\end{table}

Die Kanten des Abhängigkeitsgraphen für das VHDL-Design (und analog für Amaranth) lauten:
\[
E_1 = \{(1,3), (0,4), (3,5), (4,5), (2,5)\},
\]
was den Abhängigkeiten $\texttt{a\_in} \to \texttt{a\_unsigned}$, $\texttt{b\_in} \to \texttt{b\_unsigned}$, $\texttt{a\_unsigned} \to \texttt{alu\_result}$, $\texttt{b\_unsigned} \to \texttt{alu\_result}$, $\texttt{ctrl\_in} \to \texttt{alu\_result}$ entspricht.

\subsection{Adjazenzmatrizen}

Für $n = 6$ Knoten (Indizes $0,\ldots,5$) ergibt sich:
\[
A = B = \begin{pmatrix}
0 & 0 & 0 & 0 & 0 & 0 \\
0 & 0 & 0 & 1 & 0 & 0 \\
0 & 0 & 0 & 0 & 0 & 1 \\
0 & 0 & 0 & 0 & 0 & 1 \\
0 & 0 & 0 & 0 & 0 & 1 \\
0 & 0 & 0 & 0 & 0 & 0 \\
\end{pmatrix}
\quad \text{(VHDL = Amaranth)}
\]

\subsection{Berechnung der Signaturen}

Mit $n = 6$ gilt für die Zeilen-Signaturen (Modulo $2^6 = 64$):
\begin{align*}
\hat{\sigma}_0 &= 0 \quad \text{(Spalte 0: Nullvektor)}\\
\hat{\sigma}_1 &= 0 \quad \text{(Spalte 1: Nullvektor)}\\
\hat{\sigma}_2 &= 0 \quad \text{(Spalte 2: Nullvektor)}\\
\hat{\sigma}_3 &= 2^1 = 2 \quad \text{(Kante von Knoten 1: } 2^1\text{)}\\
\hat{\sigma}_4 &= 2^0 = 1 \quad \text{(Kante von Knoten 0: } 2^0\text{)}\\
\hat{\sigma}_5 &= 2^2 + 2^3 + 2^4 = 4 + 8 + 16 = 28 \quad \text{(Kanten von Knoten 2, 3, 4)}
\end{align*}

Da $A = B$, gilt $\hat{\boldsymbol{\sigma}}^A = \hat{\boldsymbol{\sigma}}^B = [0, 0, 0, 2, 1, 28]$.

\subsection{Ergebnis der Äquivalenzprüfung}

Die Rotation $k^* = 0$ liefert:
\[
\hat{\boldsymbol{\sigma}}^A = \mathrm{rot}_0(\hat{\boldsymbol{\sigma}}^B) = [0, 0, 0, 2, 1, 28].
\]
$\mathrm{lcs}(\hat{\boldsymbol{\sigma}}^A, \hat{\boldsymbol{\sigma}}^B) = 6 \geq 2$, also $H(D_1) \cong_s H(D_2)$.

\textbf{Schlussfolgerung:} Nach Satz~\ref{thm:struktur_impliziert_funktion} sind die VHDL-Implementation und die Amaranth-Implementation der ALU funktional äquivalent.

\section{SAT-basierte vs.\ Subgraph-basierte Äquivalenz}
\label{subsec:vergleich}

\begin{table}[H]
\centering
\begin{tabular}{|l|c|c|}
\hline
\textbf{Eigenschaft} & \textbf{SAT-basiert (Yosys)} & \textbf{Subgraph Algorithmus} \\
\hline
Laufzeit & $O(2^n)$ (worst case) & $O(n^3)$ \\
Vollständigkeit & vollständig & strukturell vollständig \\
Abstraktionsebene & Netzlisten & Abhängigkeitsgraph \\
Optimierungstoleranz & ja (Logikoptimiz.) & nein (strukturell exact) \\
Open-Source & ja (Yosys) & ja \\
Optimalität & — & SETH-optimal \\
\hline
\end{tabular}
\caption{Vergleich der Äquivalenzprüfungsverfahren}
\end{table}

Die beiden Ansätze sind komplementär:
\begin{itemize}
    \item Der SAT-basierte Ansatz prüft funktionale Äquivalenz auch unter Logikoptimierungen.
    \item Der Subgraph Algorithmus prüft strukturelle Äquivalenz mit polynomieller Garantie und ist besonders geeignet für Designs, bei denen die Strukturerhaltung explizit gefordert ist.
\end{itemize}

\section{Zusammenfassung}

Dieses Kapitel hat gezeigt, dass der Subgraph Algorithmus \cite{subgraph2026} eine theoretisch fundierte, polynomial-zeitliche Methode zur strukturellen Äquivalenzprüfung von Hardware-Designs liefert. Die wesentlichen Ergebnisse sind:

\begin{enumerate}
    \item \textbf{Satz~\ref{thm:strukturelle_aequivalenz}} charakterisiert strukturelle Äquivalenz über bidirektionale Subgraph-Beziehungen.
    \item \textbf{Satz~\ref{thm:struktur_impliziert_funktion}} beweist, dass strukturelle Äquivalenz funktionale Äquivalenz impliziert.
    \item \textbf{Satz~\ref{thm:korrektheit_aequivalenz}} garantiert die Korrektheit des Verfahrens.
    \item \textbf{Satz~\ref{thm:laufzeit_aequivalenz}} zeigt die polynomielle Laufzeit $O(n^3)$.
    \item \textbf{Satz~\ref{thm:sensitivitaet}} beweist die Empfindlichkeit gegenüber einzelnen Kantenänderungen.
    \item \textbf{Satz~\ref{thm:transitivitaet}} zeigt die Transitivität der Äquivalenzrelation.
    \item \textbf{Satz~\ref{thm:speedup}} quantifiziert den super-exponentiellen Speedup gegenüber naiver Permutationssuche.
    \item \textbf{Satz~\ref{thm:optimalitaet_aequivalenz}} beweist die SETH-bedingte Optimalität.
\end{enumerate}

\chapter{Äquivalenz: Verilog gegen VHDL}
\label{sec:verilog_vhdl_direkt}

\section{Motivation und Einordnung}

Die bisherigen Kapitel dieser Arbeit betrachteten die Äquivalenz zwischen einer
Amaranth-generierten Verilog-Netzliste und einem VHDL-Referenzdesign.  Es gibt
jedoch einen grundlegenderen Fall, der in der Praxis häufig auftritt: zwei
\emph{manuell geschriebene} Implementierungen \emph{desselben Entwurfs} in zwei
verschiedenen HDLs.  Dieser Fall tritt beispielsweise auf, wenn

\begin{itemize}
	\item ein existierendes VHDL-Design für eine Simulationsumgebung nach Verilog
	portiert wird,
	\item zwei Entwickler unabhängig voneinander dieselbe Spezifikation in
	verschiedenen Sprachen codieren (klassisches Vier-Augen-Prinzip auf
	Sprach\-ebene), oder
	\item ein Legacy-Design für neue Synthesewerkzeuge in der jeweils präferierten
	HDL neu geschrieben wird.
\end{itemize}

Als konkretes Demonstrationsobjekt dient die \texttt{alu\_7\_segment}-Schaltung
des CITEC / Bielefeld University FPGA-Labors (Kurs-Praktikum, Revision~1.0,
Martin Kaiser und Sarah Pilz, 25.\,September\,2017).  Sie liegt in zwei
gleichwertigen Fassungen vor:

\begin{itemize}
	\item \textbf{Verilog-Fassung}: \texttt{alu\_7\_segment.v} –
	kompakte, C-ähnliche Syntax, \texttt{always\,@(*)}-Blöcke,
	implizite Typumwandlung.
	\item \textbf{VHDL-Fassung}: \texttt{alu\_7\_segment.vhd} –
	streng typisiertes IEEE-1076-Modell mit expliziten
	Variablendeklarationen und mehrspaltiger \texttt{process}-Semantik.
\end{itemize}

Das Ziel dieses Kapitels ist ein \emph{vollständiger, schrittweiser
	Äquivalenzbeweis} dieser beiden Fassungen.  Wir verwenden dazu drei
komplementäre Methoden:

\begin{enumerate}
	\item \textbf{Manuelle Inspektion} der Schnittstellensignaturen (Abschnitt~\ref{subsec:schnittstellenvergleich}).
	\item \textbf{Strukturelle Äquivalenzprüfung} mit dem Subgraph Algorithmus
	\cite{subgraph2026} (Abschnitt \ref{subsec:strukturell_vv}).
	\item \textbf{Formale Verifikation} mit Yosys\,/\,SAT-Solver
	(Abschnitt~\ref{subsec:yosys_direkt}).
\end{enumerate}

\section{Spezifikation und Schnittstellenvergleich}
\label{subsec:schnittstellenvergleich}

\subsection{Gemeinsame Schnittstellenspezifikation}

Beide Implementierungen realisieren dasselbe Top-Level-Interface.
Tabelle~\ref{tab:schnittstelle} stellt die vier Ports beider Fassungen gegenüber.
Das Modul besitzt zwei Eingänge: den einbitigen Systemtakt \texttt{CLK} sowie den
16\,Bit breiten Schalterbus \texttt{SWT}, über den Operanden und Opcode zugeführt
werden.  Auf der Ausgangsseite liefert \texttt{SEG} die sieben Segmentsignale
($a$--$g$) und \texttt{AN} die vier Anodensignale für die Multiplexansteuerung der
7-Segment-Anzeigen.  In Verilog werden die Ausgänge als \texttt{output~reg}
deklariert, was die getaktete Registrierung im selben Modul ausdrückt; das VHDL-
Äquivalent verwendet \texttt{out~STD\_LOGIC\_VECTOR} mit einer internen
\texttt{signal}-Variablen.  Trotz dieser syntaktischen Unterschiede stimmen Name,
Richtung und Bitbreite jedes Ports in beiden Sprachen exakt überein, sodass die
Schnittstellensignatur als äquivalent nachgewiesen ist.

\begin{definition}[Schnittstellensignatur $\mathcal{I}$]
	\label{def:schnittstellensignatur}
	Die Schnittstellensignatur eines HDL-Moduls ist das Tupel
	\[
	\mathcal{I} = (P_\mathrm{in},\, P_\mathrm{out},\, W)
	\]
	wobei $P_\mathrm{in}$ die Menge der Eingangs-Ports, $P_\mathrm{out}$ die Menge
	der Ausgangs-Ports und $W : P_\mathrm{in} \cup P_\mathrm{out} \to \mathbb{N}$
	die Bitbreiten-Abbildung bezeichnet.
\end{definition}

\begin{sidewaystable}
	\centering
	\begin{tabular}{|l|c|c|c|c|}
		\hline
		\textbf{Port} & \textbf{Rich.} & \textbf{Breite} & \textbf{Verilog} & \textbf{VHDL} \\
		\hline
		\texttt{CLK}  & in  & 1  & \texttt{input wire}               & \texttt{in STD\_LOGIC} \\
		\texttt{SWT}  & in  & 16 & \texttt{input wire [15:0]}        & \texttt{in STD\_LOGIC\_VECTOR(15 downto 0)} \\
		\texttt{SEG}  & out & 7  & \texttt{output reg [6:0]}         & \texttt{out STD\_LOGIC\_VECTOR(6 downto 0)} \\
		\texttt{AN}   & out & 4  & \texttt{output reg [3:0]}         & \texttt{out STD\_LOGIC\_VECTOR(3 downto 0)} \\
		\hline
	\end{tabular}
	\caption{Schnittstellenvergleich: Verilog vs.\ VHDL}
	\label{tab:schnittstelle}
\end{sidewaystable}

\begin{lemma}[Schnittstellenäquivalenz]
	\label{lem:schnittstellen_equiv}
	Die Verilog-Fassung und die VHDL-Fassung besitzen identische Schnittstellensignaturen:
	\[
	\mathcal{I}_\mathrm{Verilog} = \mathcal{I}_\mathrm{VHDL}.
	\]
\end{lemma}

\begin{proof}
	Direkter Vergleich der Port-Deklarationen (Tabelle~\ref{tab:schnittstelle}).
	Für jeden Port stimmen Name, Richtung und Bitbreite überein.  Die syntaktischen
	Unterschiede (\texttt{input wire} vs.\ \texttt{in STD\_LOGIC} bzw.\
	\texttt{[15:0]} vs.\ \texttt{(15 downto 0)}) sind rein sprachliche Konventionen
	ohne semantischen Unterschied für die Netzlistensemantik.
\end{proof}

\subsection{Interne Parametrierung}

Beide Fassungen deklarieren denselben Teilungsparameter für die
Display-Taktung.  Tabelle~\ref{tab:paras} listet sämtliche Konstanten
und Opcodes beider Designs nebeneinander auf.
Die wichtigste Konstante ist \texttt{REFRESH\_COUNT~=~100\,000}: Bei einem
100\,MHz-Systemtakt ergibt sich daraus eine Display-Refresh-Rate von 1\,kHz,
die für ein flimmerfreies Multiplexbild der vier 7-Segment-Stellen benötigt
wird.  In Verilog wird sie als \texttt{localparam} definiert, in VHDL als
\texttt{constant} mit explizitem Typ \texttt{integer}; der numerische Wert
ist identisch.  Darüber hinaus sind sechs ALU-Opcodes festgelegt:
\texttt{OP\_ADD} (0), \texttt{OP\_SUB} (1), \texttt{OP\_MULT} (2),
\texttt{OP\_DIV} (3), \texttt{OP\_AND} (4) und \texttt{OP\_OR} (5).
Verilog codiert sie als 4-Bit-Binärliterale (\texttt{4'b0000} usw.),
VHDL als Bit-String-Literale (\texttt{"0000"} usw.)~-- semantisch
bezeichnen beide Notationen dieselben Ganzzahlwerte, wie die
gegenüberstellende Tabelle unmittelbar zeigt.

\begin{sidewaystable}
	\centering
	\begin{tabular}{|l|l|l|}
		\hline
		\textbf{Parameter} & \textbf{Verilog} & \textbf{VHDL} \\
		\hline
		Refresh-Konstante & \texttt{localparam REFRESH\_COUNT = 100000} & \texttt{constant REFRESH\_COUNT : integer := 100000} \\
		OP\_ADD  & \texttt{4'b0000} & \texttt{"0000"} \\
		OP\_SUB  & \texttt{4'b0001} & \texttt{"0001"} \\
		OP\_MULT & \texttt{4'b0010} & \texttt{"0010"} \\
		OP\_DIV  & \texttt{4'b0011} & \texttt{"0011"} \\
		OP\_AND  & \texttt{4'b0100} & \texttt{"0100"} \\
		OP\_OR   & \texttt{4'b0101} & \texttt{"0101"} \\
		\hline
	\end{tabular}
	\caption{Parametrierung}
	\label{tab:paras}
\end{sidewaystable}

Alle Operation Codes sowie die Taktteiler-Konstante sind numerisch identisch.

\section{Funktionale Dekomposition}
\label{subsec:dekomposition}

Beide Designs lassen sich in drei funktional unabhängige Teilblöcke zerlegen:

\begin{figure}[H]
	\centering
	\begin{tikzpicture}[
		block/.style={draw, rectangle, minimum width=7.0cm, text width=4.3cm,
			minimum height=1.1cm, text centered, font=\small},
		node distance=1.9cm
		]
		\node[block] (io)    {Input-Mapping
			{\small\texttt{a\_in, b\_in, ctrl\_in}}};
		\node[block, below of=io]  (alu)  {ALU-Logik
			{\small(kombinatorisch)}};
		\node[block, below of=alu] (mux)  {Display-Multiplexer
			{\small(sequentiell)}};
		\node[block, below of=mux] (dec)  {7-Segment-Decoder
			{\small(16 Eintr\"{a}ge)}};

		\node[block, below of=dec, node distance=1.9cm] (segout) {\texttt{SEG[6:0]}};

		\draw[->] (io)  -- node[right]{\small SWT-Bits} (alu);
		\draw[->] (alu) -- node[right]{\small \texttt{alu\_result[7:0]}} (mux);
		\draw[->] (mux) -- node[right]{\small \texttt{current\_digit, AN}} (dec);
		\draw[->] (dec) -- node[right]{\small Ausgang} (segout);
	\end{tikzpicture}
	\caption{Funktionale Dekomposition beider Designs in vier unabhängig vergleichbare Teilblöcke.}
	\label{fig:dekomposition}
\end{figure}

Die Äquivalenzprüfung wird im Folgenden blockweise geführt.  Gemäß dem
Kompositionalitätsprinzip gilt:

\begin{theorem}[Kompositionale Äquivalenz]
	\label{thm:kompositional}
	Seien $D_1$ und $D_2$ zwei Hardware-Designs, die sich als Komposition
	$D_i = B_i^{(1)} \parallel B_i^{(2)} \parallel \cdots \parallel B_i^{(k)}$
	disjunkter Teilblöcke darstellen lassen.  Wenn jedes Blockpaar
	$B_1^{(j)} \cong B_2^{(j)}$ für $j = 1,\ldots,k$ äquivalent ist und die
	Verbindungsstruktur zwischen den Blöcken in beiden Designs identisch ist, dann
	gilt $D_1 \cong D_2$.
\end{theorem}

\begin{proof}
	Sei $\vec{x}$ eine beliebige Eingabe.  Da die Verbindungsstrukturen identisch
	sind, empfängt $B_1^{(1)}$ und $B_2^{(1)}$ dieselbe Eingabe, und wegen
	$B_1^{(1)} \cong B_2^{(1)}$ produzieren sie dieselbe Ausgabe.  Diese Ausgabe
	wird als Eingabe an den nächsten Block weitergeleitet; per Induktion über die
	Tiefe der Komposition folgt, dass alle Blockpaare unter identischen Eingaben
	identische Ausgaben liefern.  Insbesondere stimmen die Top-Level-Ausgaben
	überein, also $D_1(\vec{x}) = D_2(\vec{x})$ für alle $\vec{x}$.
\end{proof}

\section{Schritt-für-Schritt-Äquivalenz der Teilblöcke}
\label{subsec:schritte}

\subsection{Schritt 1 – Input-Mapping}
\label{sssec:input_mapping}

\paragraph{Verilog:}
\begin{lstlisting}[style=hdlstyle, caption={Input-Mapping in Verilog}]
	assign a_in    = SWT[7:4];      // Bits 7..4
	assign b_in    = SWT[3:0];      // Bits 3..0
	assign ctrl_in = SWT[15:12];    // Bits 15..12
	
	assign digit_0 = b_in;
	assign digit_1 = a_in;
	assign digit_2 = alu_result[3:0];
	assign digit_3 = alu_result[7:4];
\end{lstlisting}

\paragraph{VHDL:}
\begin{lstlisting}[style=hdlstyle, caption={Input-Mapping in VHDL}]
	a_in    <= SWT(7 downto 4);
	b_in    <= SWT(3 downto 0);
	ctrl_in <= SWT(15 downto 12);
	
	digit_0 <= b_in;
	digit_1 <= a_in;
	digit_2 <= alu_result(3 downto 0);
	digit_3 <= alu_result(7 downto 4);
\end{lstlisting}

\paragraph{Vergleich:}
Beide Fassungen führen strukturell identische bit-selektive Zuweisungen durch.
Die Verilog-Notation \texttt{[7:4]} (absteigende Indizierung) und die
VHDL-Notation \texttt{(7~downto~4)} referenzieren denselben Bit-Bereich.
Formell:

\begin{lemma}[Äquivalenz des Input-Mappings]
	\label{lem:input_mapping}
	Für alle $\vec{s} \in \{0,1\}^{16}$ gelten die Gleichungen
	\begin{align*}
		a_\mathrm{in} &= \vec{s}[7{:}4], &
		b_\mathrm{in} &= \vec{s}[3{:}0], &
		\mathrm{ctrl\_in} &= \vec{s}[15{:}12]
	\end{align*}
	in beiden Implementierungen.  \hfill $\square$
\end{lemma}

\subsection{Schritt 2 – ALU-Logik}
\label{sssec:alu_logik}

Dies ist der komplexeste Teilblock.  Wir prüfen alle sechs Operationen einzeln.

\paragraph{Operationsauswahl durch \texttt{ctrl\_in}.}
Beide Designs verwenden eine \texttt{case}-Anweisung auf denselben 4-Bit
Steuerleitungen.  Die Default-Ausgabe ist in beiden Fällen \texttt{0x00}.

\paragraph{OP\_ADD ($ctrl = 0000_2$).}
\begin{itemize}
	\item Verilog: \quad\texttt{alu\_result = a\_in + b\_in;}\\
	Verilog-Addition von \texttt{reg [3:0]}-Signalen liefert aufgrund der
	impliziten Null-Erwei\-terung ein 8-Bit-Ergebnis; maximal $15+15=30$, kein
	Ãœberlauf in 8~Bit.
	\item VHDL: \quad\texttt{temp\_result := a\_unsigned + b\_unsigned;}\\
	Beide Operanden wurden zuvor auf 8~Bit null-erweitert:
	\texttt{a\_unsigned := "0000" \& unsigned(a\_in)}.  Die Addition
	\texttt{unsigned(8)} + \texttt{unsigned(8)} liefert dasselbe Ergebnis.
\end{itemize}

\begin{lemma}[Äquivalenz OP\_ADD]
	Für alle $a,b \in \{0,\ldots,15\}$ gilt
	\[
	\mathrm{ADD}_\mathrm{Verilog}(a,b) = \mathrm{ADD}_\mathrm{VHDL}(a,b) = a + b.
	\]
\end{lemma}
\begin{proof}
	Da $a, b \leq 15$, ist $a + b \leq 30 < 256$.  In beiden Implementierungen
	wird die mathematische Addition in einem 8-Bit-Ergebnis ohne Ãœberlauf
	dargestellt.
\end{proof}

\paragraph{OP\_SUB ($ctrl = 0001_2$) – Clamp-Semantik.}
Hier die Programmzeilen für \textbf{Verilog}:
\begin{lstlisting}[style=hdlstyle]
if (a_in >= b_in)
alu_result = a_in - b_in;
else
alu_result = 8'h00;
\end{lstlisting}
Hier die Programmzeilen für \textbf{VHDL}:
\begin{lstlisting}[style=hdlstyle]
sub_result := signed('0' & a_unsigned) - signed('0' & b_unsigned);
if sub_result < 0 then
temp_result := (others => '0');
else
temp_result := unsigned(sub_result(7 downto 0));
end if;
\end{lstlisting}

Auf den ersten Blick unterscheiden sich die Implementierungen: Verilog vergleicht
\emph{vor} der Subtraktion (\texttt{if a\_in >= b\_in}), VHDL \emph{nach} der
Subtraktion (Sign-Bit des 9-Bit-Ergebnisses).  Es gilt jedoch:

\begin{lemma}[Äquivalenz OP\_SUB]
	\label{lem:sub_equiv}
	Für alle $a, b \in \{0,\ldots,15\}$ gilt
	\[
	\mathrm{SUB}_\mathrm{Verilog}(a,b)
	= \mathrm{SUB}_\mathrm{VHDL}(a,b)
	= \max(a - b,\; 0).
	\]
\end{lemma}
\begin{proof}
	Da $a$ und $b$ nichtnegative ganze Zahlen sind, gilt $a < b
	\Leftrightarrow a - b < 0$.  In der Verilog-Fassung wird bei $a < b$ der Wert
	0 ausgegeben, sonst $a - b \geq 0$.  In der VHDL-Fassung wird ein 9-Bit vorzeichenbehafteter
	Wert $s = a_\mathrm{unsigned} - b_\mathrm{unsigned}$ berechnet.  Das Bit
	\texttt{sub\_result(8)} (Vorzeichenbit) ist genau dann gesetzt, wenn
	$s < 0$, d.\,h.\ wenn $a < b$.  In diesem Fall wird 0 ausgegeben, sonst die
	unteren 8~Bit von~$s$, was $a - b$ entspricht.  Beide Fallunterscheidungen
	sind daher semantisch äquivalent.
\end{proof}

\begin{figure}[H]
	\centering
	\begin{tikzpicture}
		% Number line
		\draw[->] (-0.5,0) -- (8,0) node[right]{\small $a - b$};
		\draw[thick, blue] (0,0) -- (7.5,0);
		\draw[thick, red, dashed] (-0.5,0) -- (0,0);
		\fill[blue] (0,0) circle (2pt) node[above]{\small $0$};
		% Labels
		\node[below] at (3.5,-0.3) {\small Ausgabe $= a - b$ (blau: $a \geq b$)};
		\node[below] at (-0.2,-0.3) {\small};
		%\draw[->] (-0.2,-0.7) node[left]{\small $a < b$: Ausgabe $= 0$} -- (0,-0.1);
	\end{tikzpicture}
	\caption{Clamp-Semantik der Subtraktion: für $a < b$ wird auf 0 gesättigt.}
\end{figure}

\paragraph{OP\_MULT ($ctrl = 0010_2$).}
\begin{itemize}
	\item Verilog: \texttt{alu\_result = a\_in * b\_in;} –
	Produkt zweier \texttt{[3:0]}-Werte, 8-Bit-Ergebnis, maximal $15 \times 15 = 225 < 256$.
	\item VHDL: \texttt{mult\_result := unsigned(a\_in) * unsigned(b\_in);} –
	explizite Typumwandlung, dann Multiplikation im \texttt{unsigned}-Typ;
	das Ergebnis wird in 8~Bit abgelegt.
\end{itemize}

\begin{lemma}[Äquivalenz OP\_MULT]
	Für alle $a, b \in \{0,\ldots,15\}$ gilt $a \cdot b \leq 225 < 256$, also
	kein Ãœberlauf.  Beide Implementierungen liefern $a \cdot b$. \hfill $\square$
\end{lemma}

\paragraph{OP\_DIV ($ctrl = 0011_2$) – Division-by-Zero-Behandlung.}
Hier die Programmzeilen für \textbf{Verilog}:
\begin{lstlisting}[style=hdlstyle]
if (b_in == 4'h0)
alu_result = 8'hFF;
else
alu_result = a_in / b_in;
\end{lstlisting}
Hier die Programmzeilen für \textbf{VHDL}:
\begin{lstlisting}[style=hdlstyle]
if b_unsigned = 0 then
temp_result := (others => '1');   -- ergibt 0xFF
else
div_result := a_unsigned / b_unsigned;
temp_result := div_result;
end if;
\end{lstlisting}

Beide Fassungen geben bei $b = 0$ den Wert \texttt{0xFF} (alle Bits gesetzt)
aus, was als Fehlerkode vereinbart ist.  Für $b \neq 0$ wird die ganzzahlige
Division berechnet.  In VHDL ist die Division auf \texttt{unsigned}-Typen
definiert und entspricht der ganzzahligen Division.  In Verilog entspricht
\texttt{/} für unsigned Operanden ebenfalls der ganzzahligen Division.

\begin{lemma}[Äquivalenz OP\_DIV]
	\label{lem:div_equiv}
	Für alle $a \in \{0,\ldots,15\}$, $b \in \{0,\ldots,15\}$ gilt:
	\[
	\mathrm{DIV}_\mathrm{Verilog}(a,b) = \mathrm{DIV}_\mathrm{VHDL}(a,b)
	= \begin{cases} \mathtt{0xFF} & \text{falls } b = 0, \\ \lfloor a/b \rfloor & \text{sonst.} \end{cases}
	\]
	\hfill $\square$
\end{lemma}

\paragraph{OP\_AND und OP\_OR ($ctrl = 0100_2$ bzw.\ $0101_2$).}

\textbf{Verilog}: \texttt{alu\_result = \{4'h0, a\_in \& b\_in\};}
sowie \texttt{alu\_result = \{4'h0, a\_in | b\_in\};}

\textbf{VHDL}: \texttt{temp\_result := a\_unsigned and b\_unsigned;}
sowie \texttt{temp\_result := a\_un\newline signed or b\_unsigned;}

In der Verilog-Fassung erfolgt die Null-Erweiterung explizit durch
Konkatenation (\texttt{\{4'h0,~...\}}).  In VHDL sind die Operanden bereits
auf 8~Bit null-erweitert (\texttt{a\_unsigned := "0000" \& unsigned(a\_in)}),
sodass AND und OR auf den oberen 4~Bit trivialerweise 0 liefern.  Das Ergebnis
ist in beiden Fällen die bitweise Operation auf den unteren 4~Bit, obere
4~Bit sind 0.

\begin{lemma}[Äquivalenz OP\_AND / OP\_OR]
	Für alle $a, b \in \{0,\ldots,15\}$:
	\begin{align*}
		\mathrm{AND}_\mathrm{Verilog}(a,b) &= \mathrm{AND}_\mathrm{VHDL}(a,b) = a \mathbin{\&} b, \\
		\mathrm{OR}_\mathrm{Verilog}(a,b)  &= \mathrm{OR}_\mathrm{VHDL}(a,b)  = a \mathbin{|} b.
	\end{align*}
	\hfill $\square$
\end{lemma}

\paragraph{Default-Fall.}
Verilog: \texttt{alu\_result = 8'h00;}; VHDL: \texttt{temp\_result := (others => \newline'0');}.
Beide ergeben $\texttt{0x00}$ für alle undefinierten Steuercodes. \hfill $\square$

\begin{theorem}[Funktionale Äquivalenz der ALU]
	\label{thm:alu_equiv}
	Für alle Eingaben $a \in \{0,\ldots,15\}$, $b \in \{0,\ldots,15\}$ und
	$\mathrm{ctrl} \in \{0,\ldots,15\}$ gilt:
	\[
	\mathrm{ALU}_\mathrm{Verilog}(a, b, \mathrm{ctrl}) =
	\mathrm{ALU}_\mathrm{VHDL}(a, b, \mathrm{ctrl}).
	\]
\end{theorem}

\begin{proof}
	Direkte Konsequenz der Lemmata~\ref{lem:sub_equiv}–\ref{lem:div_equiv} und
	der drei nachfolgenden Äquivalenz-\newline lemmata für AND, OR und den Default-Fall.
	Die \texttt{case}-Struktur deckt alle 16 möglichen Steuercodes ab; für die
	definierten Codes (0–5) wurde Äquivalenz einzeln bewiesen, für die restlichen
	10 Codes liefern beide Designs $\texttt{0x00}$.
\end{proof}

\subsection{Schritt 3 – Display-Multiplexer}
\label{sssec:display_mux}

Der Display-Multiplexer ist der einzige \emph{sequentielle} Teilblock.  Er
enthält ein 17-Bit-Register \texttt{refresh\_counter} und ein 2-Bit-Register
\texttt{digit\_select}.

\paragraph{Verilog (sequentiell):}
\begin{lstlisting}[style=hdlstyle, caption={Multiplexer in Verilog}]
always @(posedge CLK) begin
if (refresh_counter == REFRESH_COUNT - 1) begin
refresh_counter <= 0;
digit_select    <= digit_select + 1;
end else begin
refresh_counter <= refresh_counter + 1;
end
end
\end{lstlisting}

\paragraph{VHDL (sequentiell):}
\begin{lstlisting}[style=hdlstyle, caption={Multiplexer in VHDL}]
multiplex_process: process(CLK)
begin
if rising_edge(CLK) then
if refresh_counter = REFRESH_COUNT - 1 then
refresh_counter <= 0;
digit_select <= std_logic_vector(
unsigned(digit_select) + 1);
else
refresh_counter <= refresh_counter + 1;
end if;
end if;
end process;
\end{lstlisting}

\paragraph{Invariante und Zustandsäquivalenz.}

Beide Register (\texttt{refresh\_counter} und \texttt{digit\_\newline select}) folgen
exakt derselben Zustandsübergangsrelation:

\begin{definition}[Zustandsübergangsrelation $\delta$]
	Sei $s = (\mathit{rc}, \mathit{ds})$ und $(\mathit{rc}, \mathit{ds})$ aus der Menge $\{0,\ldots,\text{REFRESH\_COUNT}-1\} \times \{0,1,2,3\}$ der Zustand.
	Die Ãœbergangsrelation lautet:
	\[
	\delta(s) =
	\begin{cases}
		(0,\; (\mathit{ds} + 1) \bmod 4) & \text{falls } \mathit{rc} = \text{REFRESH\_COUNT} - 1, \\
		(\mathit{rc} + 1,\; \mathit{ds}) & \text{sonst.}
	\end{cases}
	\]
\end{definition}

\begin{lemma}[Äquivalenz der Zustandsübergangsrelation]
	\label{lem:delta_equiv}
	Die Zustandsübergangsrelation $\delta$ ist in der Verilog- und in der
	VHDL-Fassung identisch.
\end{lemma}

\begin{proof}
	In der VHDL-Fassung erfolgt die Inkrementierung durch explizite Typumwandlung
	\texttt{unsigned(digit\_select) + 1}.  Da \texttt{digit\_select} ein
	2-Bit-Vektor ist, entspricht $(\mathit{ds} + 1) \bmod 4$ dem natürlichen
	Ãœberlauf eines 2-Bit-Registers, der in Verilog implizit durch
	\texttt{digit\_select <= digit\_select + 1} ausgelöst wird.  Die Bedingung
	\texttt{refresh\_counter == REFRESH\_COUNT - 1} (Verilog) und
	\texttt{refresh\_counter = REFRESH\_COUNT - 1} (VHDL) sind semantisch
	identische Gleichheitsvergleiche.  Der Initialisierungszustand beider Register
	ist in beiden Fassungen $(\mathit{rc},\,\mathit{ds}) = (0,\,0)$.
\end{proof}

Die Anode-Ausgabe \texttt{AN} und die Digit-Selektion sind kombinatorisch an
\texttt{digit\_select} gebunden:

\begin{table}[H]
	\centering
	\begin{tabular}{|c|c|c|c|}
		\hline
		\textbf{\texttt{digit\_select}} & \textbf{\texttt{AN}} & \textbf{\texttt{current\_digit}} & \textbf{Anzeige} \\
		\hline
		\texttt{00} & \texttt{1110} & \texttt{digit\_0} $= b_\mathrm{in}$  & B (Einerstelle) \\
		\texttt{01} & \texttt{1101} & \texttt{digit\_1} $= a_\mathrm{in}$  & A \\
		\texttt{10} & \texttt{1011} & \texttt{digit\_2} $=$ Ergebnis[3:0]  & Ergebnis (low) \\
		\texttt{11} & \texttt{0111} & \texttt{digit\_3} $=$ Ergebnis[7:4]  & Ergebnis (high) \\
		\hline
	\end{tabular}
	\caption{Digit-Auswahl und Anoden-Aktivierung (identisch in beiden Fassungen)}
\end{table}

Diese Tabelle ist wörtlich identisch in beiden Implementierungen umgesetzt.

\subsection{Schritt 4 – 7-Segment-Decoder}
\label{sssec:seg_decoder}

Der 7-Segment-Decoder ist eine rein kombinatorische 4-nach-7-Bit-Lookup-Tabelle.

\begin{definition}[Decoder-Funktion $\mathcal{D}$]
	\[
	\mathcal{D} : \{0,\ldots,15\} \to \{0,1\}^7, \quad
	\mathcal{D}(d) = \text{7-Segment-Bitmuster für Hexziffer } d.
	\]
\end{definition}

\begin{table}[H]
	\centering
	\small
	\begin{tabular}{|c|c|c|c|c|c|}
		\hline
		\textbf{Ziffer} & \textbf{SEG (gfedcba)} & \textbf{Verilog} & \textbf{VHDL} & \textbf{Ãœbereinstimmung}\\
		\hline
		\texttt{0} & \texttt{1000000} & \texttt{7'b1000000} & \texttt{"1000000"} & \checkmark \\
		\texttt{1} & \texttt{1111001} & \texttt{7'b1111001} & \texttt{"1111001"} & \checkmark \\
		\texttt{2} & \texttt{0100100} & \texttt{7'b0100100} & \texttt{"0100100"} & \checkmark \\
		\texttt{3} & \texttt{0110000} & \texttt{7'b0110000} & \texttt{"0110000"} & \checkmark \\
		\texttt{4} & \texttt{0011001} & \texttt{7'b0011001} & \texttt{"0011001"} & \checkmark \\
		\texttt{5} & \texttt{0010010} & \texttt{7'b0010010} & \texttt{"0010010"} & \checkmark \\
		\texttt{6} & \texttt{0000010} & \texttt{7'b0000010} & \texttt{"0000010"} & \checkmark \\
		\texttt{7} & \texttt{1111000} & \texttt{7'b1111000} & \texttt{"1111000"} & \checkmark \\
		\texttt{8} & \texttt{0000000} & \texttt{7'b0000000} & \texttt{"0000000"} & \checkmark \\
		\texttt{9} & \texttt{0010000} & \texttt{7'b0010000} & \texttt{"0010000"} & \checkmark \\
		\texttt{A} & \texttt{0001000} & \texttt{7'b0001000} & \texttt{"0001000"} & \checkmark \\
		\texttt{b} & \texttt{0000011} & \texttt{7'b0000011} & \texttt{"0000011"} & \checkmark \\
		\texttt{C} & \texttt{1000110} & \texttt{7'b1000110} & \texttt{"1000110"} & \checkmark \\
		\texttt{d} & \texttt{0100001} & \texttt{7'b0100001} & \texttt{"0100001"} & \checkmark \\
		\texttt{E} & \texttt{0000110} & \texttt{7'b0000110} & \texttt{"0000110"} & \checkmark \\
		\texttt{F} & \texttt{0001110} & \texttt{7'b0001110} & \texttt{"0001110"} & \checkmark \\
		\hline
	\end{tabular}
	\caption{Vollständiger Bit-für-Bit-Vergleich der 7-Segment-Decoder-Tabelle.
		Alle 16 Einträge stimmen überein.}
	\label{tab:seg_decoder}
\end{table}

\begin{lemma}[Äquivalenz des 7-Segment-Decoders]
	\label{lem:decoder_equiv}
	Es gilt $\mathcal{D}_\mathrm{Verilog}(d) = \mathcal{D}_\mathrm{VHDL}(d)$
	für alle $d \in \{0,\ldots,15\}$.
\end{lemma}

\begin{proof}
	Vollständige Aufzählung; Tabelle~\ref{tab:seg_decoder} zeigt bitweise
	Übereinstimmung aller 16 Einträge.
\end{proof}

\section{Strukturelle Äquivalenzprüfung mit dem Subgraph Algorithmus}
\label{subsec:strukturell_vv}

Wir wenden den Subgraph Algorithmus \cite{subgraph2026} auf die vollständigen
Hardware-Abhängigkeits\-graphen beider Designs an.

\subsection{Konstruktion der Hardware-Abhängigkeitsgraphen}

\begin{definition}[Hardware-Abhängigkeitsgraph $H(D)$, vgl.\ Definition~\ref{def:hdg}]
	Der Hard\-ware-Abhängigkeitsgraph $H(D) = (V, E)$ eines Designs $D$ hat als
	Knotenmenge $V$ alle Signale und als Kantenmenge $E$ alle
	Abhängigkeitsbeziehungen: $(u, v) \in E$ gdw.\ Signal $v$ unmittelbar von
	Signal $u$ abhängt.
\end{definition}

Die Knotenmenge für beide Fassungen ist:

\begin{table}[H]
	\centering
	\begin{tabular}{|c|l|l|}
		\hline
		\textbf{Index} & \textbf{Signal} & \textbf{Typ} \\
		\hline
		0 & \texttt{SWT[3:0]}  ($b_\mathrm{in}$)   & Primäreingang \\
		1 & \texttt{SWT[7:4]}  ($a_\mathrm{in}$)   & Primäreingang \\
		2 & \texttt{SWT[15:12]} ($\mathrm{ctrl\_in}$) & Primäreingang \\
		3 & \texttt{a\_unsigned}                      & intern, kombinatorisch \\
		4 & \texttt{b\_unsigned}                      & intern, kombinatorisch \\
		5 & \texttt{alu\_result}                      & intern, kombinatorisch \\
		6 & \texttt{refresh\_counter}                 & Register (sequentiell) \\
		7 & \texttt{digit\_select}                    & Register (sequentiell) \\
		8 & \texttt{current\_digit}                   & intern, kombinatorisch \\
		9 & \texttt{AN}                               & Primärausgang \\
		10 & \texttt{SEG}                             & Primärausgang \\
		\hline
	\end{tabular}
	\caption{Knotenmenge $V$, $|V| = 11$, identisch in beiden Designs.}
	\label{tab:knoten_vv}
\end{table}

Die Kantenmenge ergibt sich aus den Datenflussabhängigkeiten:

\begin{align*}
	E = \{
	&(1, 3),\; (0, 4),\;                       && a_\mathrm{in} \to a_\mathrm{unsigned},\; b_\mathrm{in} \to b_\mathrm{unsigned}\\
	&(3, 5),\; (4, 5),\; (2, 5),\;             && \text{ALU-Eingaben} \to \mathrm{alu\_result}\\
	&(6, 7),\;                                 && \mathrm{refresh\_counter} \to \mathrm{digit\_select} \\
	&(7, 8),\; (7, 9),\;                       && \mathrm{digit\_select} \to \mathrm{current\_digit},\; \mathrm{AN}\\
	&(0, 8),\; (1, 8),\; (5, 8),\;             && \mathrm{digit\_0..3} \to \mathrm{current\_digit}\\
	&(8, 10)\;                                 && \mathrm{current\_digit} \to \mathrm{SEG}
	\}
\end{align*}

\subsection{Adjazenzmatrix}

Für $n = 11$ Knoten (Indizes 0--10):

\[
A_\mathrm{Verilog} = A_\mathrm{VHDL} =
\begin{pmatrix}
	%    0  1  2  3  4  5  6  7  8  9 10
	0& 0& 0& 0& 1& 0& 0& 0& 1& 0& 0 \\ % 0 b_in
	0& 0& 0& 1& 0& 0& 0& 0& 1& 0& 0 \\ % 1 a_in
	0& 0& 0& 0& 0& 1& 0& 0& 0& 0& 0 \\ % 2 ctrl_in
	0& 0& 0& 0& 0& 1& 0& 0& 0& 0& 0 \\ % 3 a_unsigned
	0& 0& 0& 0& 0& 1& 0& 0& 0& 0& 0 \\ % 4 b_unsigned
	0& 0& 0& 0& 0& 0& 0& 0& 1& 0& 0 \\ % 5 alu_result
	0& 0& 0& 0& 0& 0& 0& 1& 0& 0& 0 \\ % 6 refresh_counter
	0& 0& 0& 0& 0& 0& 0& 0& 1& 1& 0 \\ % 7 digit_select
	0& 0& 0& 0& 0& 0& 0& 0& 0& 0& 1 \\ % 8 current_digit
	0& 0& 0& 0& 0& 0& 0& 0& 0& 0& 0 \\ % 9 AN
	0& 0& 0& 0& 0& 0& 0& 0& 0& 0& 0 \\ %10 SEG
\end{pmatrix}
\]

Da $A_\mathrm{Verilog} = A_\mathrm{VHDL}$, sind die Abhängigkeitsgraphen
strukturell identisch.

\subsection{Signaturberechnung}

Mit $n = 11$ und der Signaturfunktion aus Definition~\ref{def:signatur}
(vgl.\ Kapitel~\ref{sec:subgraph_equiv}) ergibt sich für jeden Knoten $j$ die
Spalten-Signatur $\hat{\sigma}_j = \sum_{i : A_{ij}=1} 2^i \pmod{2^{11}}$:

\begin{table}[H]
	\centering
	\begin{tabular}{|c|l|r|l|}
		\hline
		\textbf{Knoten $j$} & \textbf{Eingehende Kanten} & $\hat{\sigma}_j$ & \textbf{Berechnung} \\
		\hline
		0  & –                                   &   0 & Eingang \\
		1  & –                                   &   0 & Eingang \\
		2  & –                                   &   0 & Eingang \\
		3  & \{1\}                               &   2 & $2^1$ \\
		4  & \{0\}                               &   1 & $2^0$ \\
		5  & \{2, 3, 4\}                         &  28 & $2^2 + 2^3 + 2^4$ \\
		6  & –                                   &   0 & (Rückkopplung aus Takt) \\
		7  & \{6\}                               &  64 & $2^6$ \\
		8  & \{0, 1, 5, 7\}                      & 163 & $2^0+2^1+2^5+2^7$ \\
		9  & \{7\}                               & 128 & $2^7$ (d. Anode, aus digit\_select) \\
		10 & \{8\}                               & 256 & $2^8$ \\
		\hline
	\end{tabular}
	\caption{Spalten-Signaturen $\hat{\sigma}_j$ für $H(D)$ mit $n = 11$.}
	\label{tab:signaturen_vv}
\end{table}

\subsection{Äquivalenzentscheid}

\begin{theorem}[Strukturelle Äquivalenz Verilog $\leftrightarrow$ VHDL]
	\label{thm:vv_strukturell}
	Die Verilog-Fassung $D_V$ und die VHDL-Fassung $D_H$ von
	\texttt{alu\_7\_segment} sind strukturell äquivalent:
	\[
	H(D_V) \cong_s H(D_H).
	\]
\end{theorem}

\begin{proof}
	Da $A_\mathrm{Verilog} = A_\mathrm{VHDL}$, gilt
	$\hat{\boldsymbol{\sigma}}^V = \hat{\boldsymbol{\sigma}}^H$ (Tabelle~\ref{tab:signaturen_vv}).
	Mit $k^* = 0$ gilt
	$\hat{\boldsymbol{\sigma}}^V = \mathrm{rot}_0(\hat{\boldsymbol{\sigma}}^H)$.
	Das LCS-Kriterium liefert
	$\mathrm{lcs}(\hat{\boldsymbol{\sigma}}^V, \hat{\boldsymbol{\sigma}}^H) = 11 \geq 2$,
	womit nach Definition~\ref{def:subgraph_beziehung} und
	Satz~\ref{thm:strukturelle_aequivalenz} strukturelle Äquivalenz folgt.
	Nach Satz~\ref{thm:struktur_impliziert_funktion} impliziert strukturelle
	Äquivalenz funktionale Äquivalenz.
\end{proof}

\section{Formale Verifikation mit Yosys}
\label{subsec:yosys_direkt}

\subsection{Besonderheit: VHDL in Yosys}

Yosys unterstützt VHDL nicht nativ.  Für den direkten Vergleich sind zwei
Wege möglich:

\begin{enumerate}
	\item \textbf{GHDL-Plugin}: Das GHDL-Yosys-Plugin \cite{ghdl2024} erlaubt
	das Laden von VHDL-Dateien direkt in Yosys.
	\item \textbf{Vorab-Konvertierung}: Das VHDL-Design wird einmalig mit
	\texttt{ghdl --synth} in eine Verilog-Netzliste überführt, die dann
	in den Standard-Yosys-Workflow einfließt.
\end{enumerate}

Wir verwenden Variante~2 als reproduzierbaren Standardweg:

\begin{lstlisting}[style=bashstyle, caption={VHDL zu Verilog via GHDL-Synthese}]
	# VHDL synthetisieren und als Verilog-Netzliste exportieren
	ghdl --synth --std=08 alu_7_segment.vhd -e alu_7_segment \
	| yosys -p "read_verilog /dev/stdin; \
	write_verilog alu_7_segment_vhdl_net.v"
\end{lstlisting}

\subsection{Vollständiges Äquivalenz-Script}

\begin{lstlisting}[style=bashstyle, caption={Yosys-Äquivalenzprüfung: Verilog vs.\ VHDL-Netzliste}]
	#!/bin/bash
	# Direkte Aequivalenzpruefung: Verilog vs. VHDL
	# Voraussetzung: alu_7_segment.v und alu_7_segment_vhdl_net.v
	
	yosys << 'EOF'
	# --- Verilog-Fassung laden (Design Under Test) ---
	read_verilog alu_7_segment.v
	hierarchy -check -top alu_7_segment
	proc; opt; flatten
	rename alu_7_segment dut_verilog
	design -stash verilog_stash
	
	# --- VHDL-Netzliste laden (Referenz) ---
	read_verilog alu_7_segment_vhdl_net.v
	hierarchy -check -top alu_7_segment
	proc; opt; flatten
	rename alu_7_segment ref_vhdl
	design -stash vhdl_stash
	
	# --- Designs zusammenfuehren ---
	design -copy-from verilog_stash -as dut_verilog dut_verilog
	design -copy-from vhdl_stash    -as ref_vhdl    ref_vhdl
	
	# --- Aequivalenzmodul erzeugen ---
	equiv_make -multiclock dut_verilog ref_vhdl equiv_top
	
	hierarchy -check -top equiv_top
	proc; opt
	
	# --- SAT-basierte kombinatorische Pruefung ---
	equiv_simple -undef -seq 5
	
	# --- Induktiver Beweis fuer sequentielle Teile ---
	equiv_induct -undef -seq 10
	
	# --- Ergebnis ---
	equiv_status
	EOF
\end{lstlisting}

\subsection{Interpretation des Yosys-Outputs}

Ein erfolgreicher Äquivalenznachweis liefert folgenden Abschnitt im Log:

\begin{lstlisting}[style=bashstyle]
	Executing EQUIV_STATUS pass.
	Found 0 unproven $equiv cells.
	Found 0 proven $equiv cells (with cex).
	Equivalence successfully proven!
\end{lstlisting}

Die Meldung \textit{„Equivalence successfully proven"} bedeutet formal:

\begin{definition}[Yosys-Äquivalenzzertifikat]
	Yosys kodiert die Negation der Äquivalenz als SAT-Formel
	\[
	\Phi = \exists \vec{x} : \mathrm{dut\_verilog}(\vec{x}) \neq \mathrm{ref\_vhdl}(\vec{x}).
	\]
	Ist $\Phi$ unerfüllbar (UNSAT), gibt Yosys \textit{„Equivalence successfully
		proven"} aus.  UNSAT von $\Phi$ ist äquivalent zu
	$\forall \vec{x} : \mathrm{dut\_verilog}(\vec{x}) = \mathrm{ref\_vhdl}(\vec{x})$,
	d.\,h.\ vollständiger funktionaler Äquivalenz.
\end{definition}

\subsection{Komplexitätsbetrachtung}

Der Eingabe-Raum der Schaltung besteht aus:
\begin{itemize}
	\item 16~Schalter: $2^{16} = 65\,536$ Kombinations-Eingaben,
	\item 1 Takt-Signal (sequentiell, modelliert als Zustandsmaschine).
\end{itemize}

Eine vollständige Simulation aller Eingaben wäre für kombinatorische Teile
zwar machbar (Brute-Force über $65\,536$ Vektoren), für das sequentielle
Multiplexer-Register mit Zustandstiefe
$\text{REFRESH\_COUNT} = 100\,000$ jedoch unpraktisch.  Die
SAT-basierte Methode von Yosys löst dieses Problem durch symbolische
Ausführung: statt alle Zustände zu enumerieren, wird ein Beweis über alle
möglichen Traces geführt.

\section{Zusammenfassender Äquivalenzbeweis}
\label{subsec:hauptsatz}

\begin{theorem}[Vollständige Äquivalenz der CITEC-ALU-Designs]
	\label{thm:hauptsatz_vv}
	Die Verilog-Implemen\-tierung \texttt{alu\_7\_segment.v} und die
	VHDL-Implementierung \texttt{alu\_7\_segment.vhd} des CITEC/Bielefeld
	University FPGA-Labors sind vollständig äquivalent: Für alle
	Eingabesequenzen $(\vec{s}_0, \vec{s}_1, \ldots) \in (\{0,1\}^{16})^\omega$
	und identische Anfangszustände gelten zu jedem Taktzyklus $t$:
	\[
	\mathrm{SEG}_V(t) = \mathrm{SEG}_H(t)
	\quad\text{und}\quad
	\mathrm{AN}_V(t) = \mathrm{AN}_H(t).
	\]
\end{theorem}

\begin{proof}
	Wir führen den Beweis kompositional nach Satz~\ref{thm:kompositional}:
	
	\begin{enumerate}
		\item \textbf{Schnittstellenäquivalenz}: Lemma~\ref{lem:schnittstellen_equiv}
		zeigt, dass Ports und Bitbreiten überein\-stimmen.
		\item \textbf{Input-Mapping}: Lemma~\ref{lem:input_mapping} zeigt, dass
		die Bit-Selektion aus \texttt{SWT} in beiden Fassungen identisch ist.
		\item \textbf{ALU-Äquivalenz}: Satz~\ref{thm:alu_equiv} beweist die
		funktionale Äquivalenz für alle sechs Operationen und alle
		$16 \times 16 = 256$ Operandenkombinationen.
		\item \textbf{Sequentielle Äquivalenz}: Lemma~\ref{lem:delta_equiv} zeigt,
		dass die Zustandsübergangsrelation des Display-Multiplexers identisch
		ist; mit identischem Anfangszustand $(0,0)$ sind alle Zustände
		zu jedem Zeitpunkt $t$ gleich.
		\item \textbf{Decoder-Äquivalenz}: Lemma~\ref{lem:decoder_equiv} beweist
		die Übereinstimmung aller 16 Einträge der Lookup-Tabelle.
		\item \textbf{Strukturelle Äquivalenz}: Satz~\ref{thm:vv_strukturell}
		bestätigt die Übereinstimmung der Hardware-Abhängigkeitsgraphen.
	\end{enumerate}
	
	Da alle Teilblöcke äquivalent sind und die Verbindungsstruktur
	(Abbildung~\ref{fig:dekomposition}) in beiden Designs identisch ist, folgt
	die vollständige Äquivalenz nach Satz~\ref{thm:kompositional}.
\end{proof}

\section{Warum benötigt VHDL mehr Anweisungen?}

Als Nebenbefund der Äquivalenzprüfung lässt sich quantifizieren, wo die
\emph{syntaktische Mehrarbeit} in VHDL gegenüber Verilog entsteht, obwohl das
Syntheseergebnis identisch ist:

\begin{table}[H]
	\centering
	\small
	\begin{tabular}{|l|c|c|l|}
		\hline
		\textbf{Ursache} & \textbf{Verilog} & \textbf{VHDL} & \textbf{Begründung} \\
		\hline
		Port-Deklaration & 4 Zeilen & 6 Zeilen & Explizite \texttt{entity}/\texttt{architecture}-Trennung \\
		Typimporte       & 0 Zeilen & 3 Zeilen & \texttt{use IEEE.numeric\_std.all} etc.\ \\
		Interne Variablen & 0 & 6 & VHDL-Prozessvariablen für Typsicherheit \\
		Typumwandlungen  & 0 & 4 & \texttt{unsigned()}, \texttt{std\_logic\_vector()} \\
		Prozess-Boilerplate & 0 & 2 pro Prozess & \texttt{begin}/\texttt{end process name} \\
		Case-Syntax      & $\approx 1$ Zeile & $\approx 2$ Zeilen & \texttt{when X =>} vs.\ \texttt{default:} \\
		\hline
		\textbf{Gesamt (ca.)} & $\approx 120$ & $\approx 170$ & $+42\,\%$ \\
		\hline
	\end{tabular}
	\caption{Zeilenzahl-Vergleich: Syntaktischer Mehraufwand in VHDL ohne semantischen Unterschied.}
\end{table}

Die VHDL-Spezifikation (IEEE\,1076) fordert diese Explizitheit
\emph{bewusst}, um Mehrdeutigkeiten bei der Typeninferenz auszuschließen.
Die formale Äquivalenzprüfung bestätigt, dass dieser Aufwand rein syntaktischer
Natur ist: das Syntheseergebnis – und damit die realisierte Hardware – ist
für beide Designs identisch.

% =====================================================================
\chapter{Simulations-basierte Testautomatisierung und Code-Ãœber\-deckung}
\label{sec:pyuvm_aoi}
% =====================================================================

\section{Motivation und Einordnung}

Die in den Kapiteln~\ref{sec:subgraph_equiv} und~\ref{sec:verilog_vhdl_direkt}
vorgestellte formale Verifikation beweist die Korrektheit eines Designs für
\emph{alle} Eingaben in einem einzigen, mathematisch abgeschlossenen Schritt.
Diese Stärke ist zugleich ihre konzeptionelle Beschränkung: Formale
Äquivalenzprüfung beantwortet statische Fragen (sind zwei Designs äquivalent?)
und setzt voraus, dass das Referenzmodell korrekt spezifiziert ist. Was sie
\emph{nicht} unmittelbar beantwortet, ist:

\begin{itemize}
    \item Welche \emph{Eingabevektoren} wurden in der Simulation tatsächlich
          angelegt, und welche Codepfade blieben unbeobachtet?
    \item Deckt die Teststimulus-Menge alle funktionalen Verhaltensklassen ab
          (\emph{funktionale Ãœberdeckung})?
    \item Wie verhält sich das Design unter randomisierten, realitätsnahen
          Eingabesequenzen im Zeitverlauf?
    \item Lassen sich Fehler systematisch injizieren, und erkennt das
          Verifikationssystem diese zuverlässig?
\end{itemize}

Diese Fragen sind Gegenstand der \emph{simulations-basierten Testautomatisierung}
— einer eigenständigen, in der ASIC- und FPGA-Industrie weit verbreiteten
Verifikationsmethodik. Sie basiert auf dem UVM-Standard (Universal Verification
Methodology, IEEE~1800.2~\cite{uvm2020}), der eine standardisierte
Klassenstruktur für wiederverwendbare, skalierbare Testbench-Architekturen
bereitstellt.

\begin{definition}[Simulations-basierte Verifikation]
Simulations-basierte Verifikation bezeichnet die Methodik, ein
Hardware-Design (DUT, \emph{Device Under Test}) unter kontrollierten,
oft zufällig erzeugten Stimuli durch Simulation auszuführen und die
Korrektheit der Ausgaben anhand eines unabhängigen Referenzmodells (\emph{Golden
Model}) zu prüfen. Die Güte einer simulations-basierten Verifikation wird durch
\emph{Ãœberdeckungsmetriken} quantifiziert.
\end{definition}

In diesem Kapitel wird die Anwendung von \textbf{pyUVM} \cite{pyuvm2024} — einer
Python-Implemen\-tierung des UVM-Standards auf Basis von cocotb \cite{cocotb2024}
— auf den AOI-2-2-Schaltkreis dargestellt. Der Schaltkreis liegt sowohl in
Verilog als auch in VHDL vor; beide Sprachvarianten werden durch dieselbe
Testbench simuliert. Das Kapitel beschreibt vollständig die
Testbench-Architektur, alle Sequenz- und Ãœberdeckungsstrategien, den
Simulations-Workflow sowie die Werkzeuge zur Coverage-Analyse und
Laufzeit-Profiling.

\section{Der AOI-2-2-Schaltkreis als Demonstrationsobjekt}
\label{subsec:aoi_schaltkreis}

\subsection{Logische Spezifikation}

Das AND-OR-INVERT-Gate (AOI) der Konfiguration $2{-}2$ — kurz AOI-2-2 —
implementiert die Boolesche Funktion

\begin{equation}
    Y = \overline{(A \cdot B) + (C \cdot D)},
    \label{eq:aoi22}
\end{equation}

die das logische NICHT des logischen ODER zweier UND-Terme ist.
Das AOI-2-2-Gate ist ein Standardzelle, die in fast jedem CMOS-Prozess
als einzelnes Gatter mit geringer Fläche und hoher Schaltgeschwindigkeit
realisiert werden kann.

\begin{definition}[AOI-2-2-Funktion]
Seien $A, B, C, D \in \{0, 1\}$. Die AOI-2-2-Funktion lautet
$f_{\mathrm{AOI}}(A, B, C, D) = \neg((A \wedge B) \vee (C \wedge D))$.
\end{definition}

Die vollständige Wahrheitstabelle mit allen $2^4 = 16$ Eingabebelegungen ist:

\begin{table}[H]
\centering
\small
\begin{tabular}{|c|c|c|c||c|c|c||c|}
\hline
\textbf{A} & \textbf{B} & \textbf{C} & \textbf{D} & $A \cdot B$ & $C \cdot D$ & $(A \cdot B)+(C \cdot D)$ & $Y$ \\
\hline
0&0&0&0&0&0&0&1\\
0&0&0&1&0&0&0&1\\
0&0&1&0&0&0&0&1\\
0&0&1&1&0&1&1&\textbf{0}\\
0&1&0&0&0&0&0&1\\
0&1&0&1&0&0&0&1\\
0&1&1&0&0&0&0&1\\
0&1&1&1&0&1&1&\textbf{0}\\
1&0&0&0&0&0&0&1\\
1&0&0&1&0&0&0&1\\
1&0&1&0&0&0&0&1\\
1&0&1&1&0&1&1&\textbf{0}\\
1&1&0&0&1&0&1&\textbf{0}\\
1&1&0&1&1&0&1&\textbf{0}\\
1&1&1&0&1&0&1&\textbf{0}\\
1&1&1&1&1&1&1&\textbf{0}\\
\hline
\end{tabular}
\caption{Vollständige Wahrheitstabelle der AOI-2-2-Funktion. $Y=0$ (fett) tritt
bei 6 von 16 Eingabebelegungen auf; $Y=1$ bei 10 von 16.}
\label{tab:aoi_wahrheit}
\end{table}

Die Funktion liefert $Y = 0$ genau dann, wenn mindestens einer der beiden
AND-Terme $(A \cdot B)$ oder $(C \cdot D)$ den Wert~1 annimmt, d.\,h.\ wenn
beide Eingänge des ersten oder des zweiten Paares gleichzeitig High sind.

Im Demonstrationsdesign werden die vier logischen Eingänge auf die vier
niederwertigen Schiebeschalter des Basys-3-FPGA-Boards abgebildet
(\texttt{SWT[3:0]}), und das Ergebnis~$Y$ wird auf dem ersten Digit des
4-stelligen 7-Segment-Displays als \texttt{'0'} oder \texttt{'1'} angezeigt.

\subsection{Verilog-Implementation}
\label{sssec:aoi_verilog}

Die Verilog-Implementierung \texttt{aoi\_2\_2.v} beschreibt den Schaltkreis
als kombinatorisches Modul mit zwei Ausgangsports für die Siebensegmentanzeige:

\begin{lstlisting}[style=hdlstyle, caption={Verilog-Implementation des AOI-2-2-Schaltkreises (\texttt{aoi\_2\_2.v})}]
module aoi_2_2 (
    input  wire [3:0] SWT,   // SWT[3:0] fuer a,b,c,d
    output wire [6:0] SEG,   // 7-Segment-Segmente
    output wire [3:0] AN     // Anoden (Digit-Aktivierung)
);
    wire a, b, c, d, y;

    assign a = SWT[0];
    assign b = SWT[1];
    assign c = SWT[2];
    assign d = SWT[3];

    // AOI-Logik: Y = NOT((A AND B) OR (C AND D))
    assign y = ~((a & b) | (c & d));

    // Nur Digit 0 aktiv (active-low)
    assign AN = 4'b1110;

    // 7-Segment-Anzeige: '0' -> 1000000, '1' -> 1111001
    assign SEG = (y == 1'b0) ? 7'b1000000 : 7'b1111001;
endmodule
\end{lstlisting}

\paragraph{Merkwürdigkeiten und Designentscheidungen.}
Die gesamte Logik ist rein kombinatorisch — alle \texttt{assign}-Anweisungen
werden sofort und zeitparallel ausgewertet. Die Port-Hierarchie folgt dem
Basys-3-Board-Mapping: \texttt{SWT[0]}~$\widehat{=}$~A, \texttt{SWT[1]}~$\widehat{=}$~B,
\texttt{SWT[2]}~$\widehat{=}$~C, \texttt{SWT[3]}~$\widehat{=}$~D.

Der 7-Segment-Decoder ist als ternärer Ausdruck kodiert:
\begin{itemize}
    \item \texttt{7'b1000000} ($=0x40$) kodiert die Ziffer \texttt{'0'}: Segmente $a$–$f$
          an, Segment~$g$ aus.
    \item \texttt{7'b1111001} ($=0x79$) kodiert die Ziffer \texttt{'1'}: nur
          Segmente~$b$ und~$c$ an.
\end{itemize}

\subsection{VHDL-Implementation}
\label{sssec:aoi_vhdl}

Die VHDL-Implementierung \texttt{aoi\_2\_2.vhd} realisiert dieselbe Funktion
in stark typisierter IEEE-1076-Syntax. Kennzeichnend ist die explizite
Trennung zwischen Entity (Schnittstelle) und Architecture (Verhalten):

\begin{lstlisting}[style=hdlstyle, caption={VHDL-Implementation des AOI-2-2-Schaltkreises (\texttt{aoi\_2\_2.vhd})}]
library ieee;
use ieee.std_logic_1164.all;

entity aoi_2_2 is
  port(
    SWT : in  std_logic_vector(3 downto 0);
    SEG : out std_logic_vector(6 downto 0);
    AN  : out std_logic_vector(3 downto 0)
  );
end aoi_2_2;

architecture rtl of aoi_2_2 is
  signal a, b, c, d : std_logic;
  signal y          : std_logic;
begin
  a <= SWT(0);  b <= SWT(1);
  c <= SWT(2);  d <= SWT(3);

  -- AOI-Logik
  y <= not((a and b) or (c and d));

  AN <= "1110";  -- active-low, Digit 0

  -- 7-Segment-Decoder als process
  seg_display: process(y)
  begin
    case y is
      when '0'    => SEG <= "1000000";  -- Ziffer '0'
      when '1'    => SEG <= "1111001";  -- Ziffer '1'
      when others => SEG <= "1111111";  -- alle aus
    end case;
  end process seg_display;
end rtl;
\end{lstlisting}

\paragraph{Vergleich Verilog vs.\ VHDL.}
Beide Implementierungen realisieren exakt dieselbe Funktion (Gleichung~\ref{eq:aoi22}).
Die wesentlichen syntaktischen Unterschiede sind:
\begin{itemize}
    \item Verilog verwendet den ternären Operator \texttt{?:} für den Decoder,
          VHDL einen expliziten \texttt{process}-Block mit \texttt{case}-Anweisung.
    \item VHDL deklariert interne Signale explizit in der \texttt{architecture}-
          Deklarationszone; Verilog deklariert sie inline mit \texttt{wire}.
    \item Beide Fassungen setzen \texttt{AN = "1110"} konstant; beide kodieren
          die 7-Segment-Muster bitidentisch.
\end{itemize}
Die formale Äquivalenz beider Fassungen folgt aus der Analyse in
Kapitel~\ref{sec:verilog_vhdl_direkt}.

\section{Universal Verification Methodology in Python}
\label{subsec:pyuvm_framework}

\subsection{UVM -- Grundkonzepte und Klassen\-hierarchie}

Die Universal Verification Methodology (UVM) ist ein IEEE-standardisiertes
Rahmenwerk (IEEE~1800.2~\cite{uvm2020}) für die strukturierte Testbench-
Entwicklung in der Hardware-Verifikation. UVM definiert eine Hierarchie von
wiederverwendbaren Basisklassen und einen standardisierten Lebenszyklus
(Phasen), in dessen Rahmen Testkomponenten koordiniert operieren.

\begin{definition}[UVM-Phasenhierarchie]
Eine UVM-Simulation durchläuft die folgenden Pflichtphasen sequentiell:
\begin{enumerate}
    \item \textbf{build\_phase}: Instantiierung und Konfiguration von Komponenten.
    \item \textbf{connect\_phase}: Verbindung von Ports und Exporten.
    \item \textbf{end\_of\_elaboration\_phase}: Finalisierung der Designhierarchie.
    \item \textbf{start\_of\_simulation\_phase}: Initialisierungen vor dem Lauf.
    \item \textbf{run\_phase}: Eigentliche Simulation (zeitbehaftet, parallel).
    \item \textbf{extract\_phase, check\_phase, report\_phase}: Auswertung, Prüfung, Berichterstellung.
\end{enumerate}
\end{definition}

\textbf{pyUVM} \cite{pyuvm2024} implementiert die UVM-Klassen\-hierarchie in
reinem Python und nutzt \textbf{cocotb} als Ausführungs\-umgebung. Cocotb ist ein
Python-Framework für die Simulation digitaler Hardware; es verbindet den
Python-Interpreter mit dem Simulator (Icarus, Verilator, GHDL, ModelSim) über
eine gemeinsame Schnittstelle. pyUVM-Testbenches schreiben sich als
gewöhnliche Python-Module, was folgende Vorteile bietet:

\begin{itemize}
    \item \textbf{Spracheinheit}: Design (Amaranth) und Testbench (pyUVM) teilen denselben Python-Ökosystem.
    \item \textbf{Native Debugging}: Python-Debugger und Profiler sind direkt anwendbar.
    \item \textbf{Bibliotheksintegration}: NumPy, SciPy, Matplotlib und andere
          Bibliotheken können direkt in Testbenches genutzt werden.
    \item \textbf{UVM-Kompatibilität}: pyUVM folgt dem UVM-1800.2-Standard; mit UVM
          vertraute Engineers können sofort produktiv arbeiten.
\end{itemize}

\subsection{Kernklassen in der vorliegenden Testbench}

Abbildung~\ref{fig:uvm_architektur} zeigt die Schicht-Architektur der
implementierten Testbench. Die Klassen lassen sich in vier Schichten gliedern:

\begin{enumerate}
    \item \textbf{Stimulus-Schicht}: \texttt{AoiSeqItem}, \texttt{RandomSeq},
          \texttt{ExhaustiveSeq}, \texttt{CornerCaseSeq}, \texttt{Test\newline AllSeq}
    \item \textbf{Treiber-Schicht}: \texttt{Driver} (mit Analyse-Port)
    \item \textbf{Prüf-Schicht}: \texttt{Scoreboard}, \texttt{Coverage}
    \item \textbf{Integrations-Schicht}: \texttt{AoiEnv}, \texttt{AoiTest},
          \texttt{AoiExhaustiveTest}, \texttt{AoiCornerTest}, \texttt{AoiTestErrors}
\end{enumerate}

\begin{figure}[H]
\centering
\begin{tikzpicture}[
    block/.style={draw, rectangle, minimum width=3cm, minimum height=0.8cm,
                  text centered, font=\small},
    layer/.style={draw, dashed, rectangle, minimum width=10cm},
    node distance=1.2cm
]
    \node[block] (seqitem) {\texttt{AoiSeqItem}};
    \node[block, right of=seqitem, xshift=3.5cm] (seqs) {Sequenzen (4x)};
    \node[block, below of=seqitem, xshift=1.7cm] (seqr) {\texttt{uvm\_sequencer}};
    \node[block, below of=seqr] (driver) {\texttt{Driver}};
    \node[block, right of=driver, xshift=3.5cm] (dut) {DUT (Verilog/VHDL)};
    \node[block, below of=driver, xshift=-1.7cm] (scoreboard) {\texttt{Scoreboard}};
    \node[block, right of=scoreboard, xshift=3.5cm] (coverage) {\texttt{Coverage}};

    \draw[->] (seqs) -- node[above]{\small Items} (seqitem);
    \draw[->] (seqitem) -- (seqr);
    \draw[->] (seqr) -- node[right]{\small get\_next\_item} (driver);
    \draw[->] (driver) -- node[above]{\small SWT} (dut);
    \draw[->] (dut) -- node[below]{\small SEG} (driver);
    \draw[->] (driver) -- node[right]{\small ap.write} (scoreboard);
    \draw[->] (driver) -- node[right]{\small ap.write} (coverage);
\end{tikzpicture}
\caption{pyUVM-Testbench-Architektur für den AOI-2-2-Schaltkreis.}
\label{fig:uvm_architektur}
\end{figure}

\section{Testbench-Architektur im Detail}
\label{subsec:tb_architektur}

\subsection{Sequenz-Item: \texttt{AoiSeqItem}}

Ein \emph{Sequenz-Item} ist die atomare Dateneinheit, die zwischen Sequenzen
und dem Treiber ausgetauscht wird. Es kapselt genau einen Transaktionsvektor:

\begin{lstlisting}[style=pythonstyle, caption={\texttt{AoiSeqItem} -- Transaktions-Datenklasse}]
class AoiSeqItem(uvm_sequence_item):
    def __init__(self, name, a=0, b=0, c=0, d=0):
        super().__init__(name)
        self.a = a; self.b = b; self.c = c; self.d = d
        self.result = None   # wird vom Treiber befuellt

    def randomize(self):
        self.a = random.randint(0, 1)
        self.b = random.randint(0, 1)
        self.c = random.randint(0, 1)
        self.d = random.randint(0, 1)

    def __str__(self):
        return (f"{self.get_name()}: A={self.a} B={self.b} "
                f"C={self.c} D={self.d} -> Y={self.result}")
\end{lstlisting}

Das Feld \texttt{result} ist nach der Instantiierung \texttt{None} und wird
vom Treiber nach der DUT-Simulation mit dem beobachteten Ausgabewert belegt.
Damit dient dasselbe Item-Objekt sowohl als Stimulus als auch als Ergebnisträger.

\subsection{Sequenzen}
\label{sssec:sequenzen}

Sequenzen erzeugen Item-Ströme und übergeben diese über den Sequenzierer an den
Treiber. Es wurden vier Sequenzklassen implementiert:

\paragraph{\texttt{ExhaustiveSeq} -- Erschöpfende Abdeckung.}
Iteriert systematisch über alle 16 Eingabebelegungen $\{(a,b,c,d) \mid a,b,c,d \in \{0,1\}\}$
durch vier verschachtelte Schleifen. Diese Sequenz garantiert $100\,\%$
funktionale Eingangsüberdeckung.

\paragraph{\texttt{RandomSeq} -- Zufall\-stimulation.}
Erzeugt 20 Items mit zufällig randomisierten Eingaben. Dient zur Simulation
eines realitätsnahen, nicht a-priori bekannten Stimulus-Stroms und zur
Erkennung von Fehlern, die bei systematischen Tests möglicherweise übersprungen
würden.

\paragraph{\texttt{CornerCaseSeq} -- Eckfall-Tests.}
Testet acht manuell ausgewählte, semantisch bedeutsame Eingabewerte:
alle Nullen (Minimalfall, $Y=1$), alle Einsen (beide AND-Terme aktiv, $Y=0$),
nur der erste AND-Term aktiv, nur der zweite, und alternierende Muster.
Eckfälle decken häufig Implementierungsfehler auf, die zufällige Tests
statistisch selten treffen.

\paragraph{\texttt{TestAllSeq} -- Kombinierte Sequenz.}
Führt nacheinander \texttt{ExhaustiveSeq} und \texttt{Random\newline Seq} aus.
Damit wird in einem einzigen Testlauf sowohl vollständige Eingangsüberdeckung
als auch realistisches Randomisierungsverhalten erreicht.

\begin{lstlisting}[style=pythonstyle, caption={\texttt{ExhaustiveSeq} und \texttt{CornerCaseSeq}}]
class ExhaustiveSeq(uvm_sequence):
    async def body(self):
        for a in range(2):
          for b in range(2):
            for c in range(2):
              for d in range(2):
                item = AoiSeqItem("ex", a, b, c, d)
                await self.start_item(item)
                await self.finish_item(item)

class CornerCaseSeq(uvm_sequence):
    async def body(self):
        test_cases = [
            (0,0,0,0), (1,1,1,1),  # Extremwerte
            (1,1,0,0), (0,0,1,1),  # je ein AND-Term aktiv
            (1,0,0,0), (0,0,0,1),  # Einzeleingaben
            (1,0,1,0), (0,1,0,1),  # alternierende Muster
        ]
        for a,b,c,d in test_cases:
            item = AoiSeqItem("cc", a, b, c, d)
            await self.start_item(item)
            await self.finish_item(item)
\end{lstlisting}

\subsection{Treiber: \texttt{Driver}}
\label{sssec:driver}

Der Treiber ist das Bindeglied zwischen der pyUVM-Testbench und dem
DUT-Interface auf Hardware-Ebene. Er implementiert
\texttt{uvm\_driver} und führt folgende Schritte pro Item aus:

\begin{enumerate}
    \item Zusammensetzung des 4-Bit Schalterwertes:
          \texttt{SWT = (d << 3) | (c << 2) | (b << 1) | a}
    \item Anlegen des Wertes an \texttt{dut.SWT}
    \item Warten auf Einschwingen der kombinatorischen Logik: 2\,ns (\texttt{Timer(2, "ns")})
    \item Lesen und Dekodieren des Ausgabewertes von \texttt{dut.SEG}:
          \texttt{0x40}~$\widehat{=}$~$Y=0$, \texttt{0x79}~$\widehat{=}$~$Y=1$
    \item Befüllen von \texttt{item.result}
    \item Broadcast des vollständigen Items über den Analyse-Port \texttt{ap}
\end{enumerate}

\begin{lstlisting}[style=pythonstyle, caption={Treiber-Logik: DUT-Ansteuerung und Ergebnis-Erfassung}]
class Driver(uvm_driver):
    def build_phase(self):
        self.ap = uvm_analysis_port("ap", self)
        self.dut = cocotb.top

    async def run_phase(self):
        while True:
            item = await self.seq_item_port.get_next_item()
            swt = (item.d << 3) | (item.c << 2) | \
                  (item.b << 1) | item.a
            self.dut.SWT.value = swt
            await Timer(2, unit="ns")
            seg = int(self.dut.SEG.value)
            item.result = 0 if seg == 0x40 else 1
            cmd = (item.a, item.b, item.c, item.d, item.result)
            self.ap.write(cmd)
            self.seq_item_port.item_done()
\end{lstlisting}

Die 2\,ns Wartezeit ist ausreichend, da der AOI-2-2-Schaltkreis rein
kombinatorisch ist — es gibt keine Takt-Flanken-Abhängigkeit. Für
sequentielle Designs wäre hier ein \texttt{RisingEdge}-Trigger notwendig.

\subsection{Ãœberdeckungs-Komponente: \texttt{Coverage}}
\label{sssec:coverage_comp}

Die \texttt{Coverage}-Komponente ist ein \texttt{uvm\_subscriber} — eine
Spezialform von \texttt{uvm\_component}, die direkt auf den Analyse-Port des
Treibers abonniert ist. Sie empfängt jede Transaktion und trägt die beobachtete
Eingabebelegung $(a, b, c, d)$ in eine Python-\texttt{set}-Datenstruktur ein:

\begin{lstlisting}[style=pythonstyle, caption={\texttt{Coverage}: Ãœberdeckungs-Tracking und Berichterstattung}]
class Coverage(uvm_subscriber):
    def end_of_elaboration_phase(self):
        self.cvg = set()   # beobachtete (a,b,c,d)-Tupel

    def write(self, cmd):
        (a, b, c, d, _) = cmd
        self.cvg.add((a, b, c, d))

    def report_phase(self):
        total = 16            # 2^4 moegliche Eingaben
        covered = len(self.cvg)
        pct = covered / total * 100
        self.logger.info(
            f"Ueberdeckung: {covered}/{total} ({pct:.1f}%)")
        if covered < total:
            missing = [(a,b,c,d)
                for a in range(2) for b in range(2)
                for c in range(2) for d in range(2)
                if (a,b,c,d) not in self.cvg]
            self.logger.error(f"Fehlend: {missing}")
            assert False    # Test schlaegt fehl
        else:
            self.logger.info("Alle Eingabekombinationen abgedeckt!")
\end{lstlisting}

\begin{definition}[Funktionale Eingangsüberdeckung]
\label{def:func_coverage}
Die \emph{funktionale Eingangsüber\-deckung} des AOI-2-2-Schaltkreises ist
\[
    C_{\mathrm{fcov}} = \frac{|\{(a,b,c,d) : \text{beobachtet}\}|}{2^4} \in [0, 1].
\]
Eine Coverage von $C_{\mathrm{fcov}} = 1$ bedeutet, dass alle 16 Eingabevektoren
mindestens einmal angelegt wurden; der Testerfolg ist dann garantiert.
\end{definition}

\begin{theorem}[Ãœberdeckungsgarantie der ExhaustiveSeq]
\label{thm:coverage_exhaustive}
$\mathrm{ExhaustiveSeq}$ erreicht stets $C_{\mathrm{fcov}} = 1$.
\end{theorem}

\begin{proof}
Die vierfach verschachtelte Schleife iteriert genau einmal über jeden Wert
$(a, b, c, d)$ und $(a, b, c, d) \in \{0,1\}^4$; da $|\{0,1\}^4| = 16$, wird jede der
16 möglichen Eingabebelegungen exakt einmal angelegt und durch die
\texttt{Coverage}-Komponente registriert.
\end{proof}

\subsection{Scoreboard}
\label{sssec:scoreboard}

Das \texttt{Scoreboard} empfängt über eine FIFO-Queue alle Transaktionen vom
Treiber-Analyse-Port, berechnet für jede den erwarteten Ausgabewert mittels
der Python-Referenzfunktion \texttt{aoi\_prediction()} und vergleicht ihn mit
dem beobachteten DUT-Ausgabewert:

\begin{lstlisting}[style=pythonstyle, caption={Scoreboard: Soll-Ist-Vergleich in der \texttt{check\_phase}}]
class Scoreboard(uvm_component):
    def build_phase(self):
        self.cmd_fifo = uvm_tlm_analysis_fifo("cmd_fifo", self)
        self.cmd_get_port = uvm_get_port("cmd_get_port", self)
        self.cmd_export = self.cmd_fifo.analysis_export

    def connect_phase(self):
        self.cmd_get_port.connect(self.cmd_fifo.get_export)

    def check_phase(self):
        passed = True
        while self.cmd_get_port.can_get():
            _, cmd = self.cmd_get_port.try_get()
            a, b, c, d, actual = cmd
            expected = aoi_prediction(a, b, c, d, self.errors)
            if expected != actual:
                self.logger.error(
                    f"FAIL: A={a} B={b} C={c} D={d} "
                    f"-> erwartet {expected}, erhalten {actual}")
                passed = False
            else:
                self.logger.info(f"OK: {a},{b},{c},{d} -> {actual}")
        assert passed
\end{lstlisting}

Das Scoreboard empfängt sowohl \texttt{Scoreboard.cmd\_export} als auch
\texttt{Coverage.analy\newline sis\_export} über denselben Analyse-Port des Treibers —
ein klassisches UVM-\emph{Fan-Out}-Muster.

\subsection{Testumgebung: \texttt{AoiEnv}}
\label{sssec:aoi_env}

Die Umgebungsklasse \texttt{AoiEnv} aggregiert alle Verifikationskomponenten
und stellt die Verbindungen zwischen ihnen her:

\begin{lstlisting}[style=pythonstyle, caption={\texttt{AoiEnv}: Aufbau und Vernetzung}]
class AoiEnv(uvm_env):
    def build_phase(self):
        self.seqr = uvm_sequencer("seqr", self)
        ConfigDB().set(None, "*", "SEQR", self.seqr)
        self.driver    = Driver.create("driver", self)
        self.coverage  = Coverage("coverage", self)
        self.scoreboard = Scoreboard("scoreboard", self)

    def connect_phase(self):
        self.driver.seq_item_port.connect(
            self.seqr.seq_item_export)
        self.driver.ap.connect(
            self.scoreboard.cmd_export)
        self.driver.ap.connect(
            self.coverage.analysis_export)
\end{lstlisting}

Die \texttt{ConfigDB().set(None, "*", "SEQR", self.seqr)}-Anweisung in der
\texttt{build\_phase} registriert den Sequenzierer global in der
UVM-Konfigurationsdatenbank, sodass Sequenzen ihn über
\texttt{ConfigDB().get(None, "", "SEQR")} abrufen können — ein Standard-UVM-Idiom für
komponentenübergreifende Referenzen.

\subsection{Test-Klassen}
\label{sssec:test_klassen}

Vier Test-Klassen decken unterschiedliche Verifikationsszenarien ab:

\begin{table}[H]
\centering
\begin{tabular}{|l|l|p{4.5cm}|}
\hline
\textbf{Klasse} & \textbf{Sequenz} & \textbf{Zweck} \\
\hline
\texttt{AoiTest} & \texttt{TestAllSeq} & Vollständig: erschöpfend + zufällig; $C_{\mathrm{fcov}}=1$ erzwungen \\
\texttt{AoiExhaustiveTest} & \texttt{ExhaustiveSeq} & Nur erschöpfend; 16 Vektoren \\
\texttt{AoiCornerTest} & \texttt{CornerCaseSeq} & 8 Eckfälle; Coverage-Fehler deaktiviert \\
\texttt{AoiTestErrors} & \texttt{TestAllSeq} & Fehlerinjektion: erwartet Testfehler \\
\hline
\end{tabular}
\caption{Ãœbersicht der vier Test-Klassen und ihrer Verifikationsziele.}
\label{tab:test_klassen}
\end{table}

\texttt{AoiTestErrors} ist mit dem Dekorator \texttt{@pyuvm.test(expect\_fail=True)}
annotiert, wodurch ein \emph{Bestehen des Tests} dort registriert wird,
wenn — und nur wenn — das Scoreboard einen Fehler meldet. Dieser anti-zyklische
Test verifiziert das Verifikationssystem selbst: Er stellt sicher, dass die
Testinfrastruktur tatsächlich in der Lage ist, Fehler zu erkennen.

\begin{lstlisting}[style=pythonstyle, caption={Fehlerinjektions-Test \texttt{AoiTestErrors}}]
@pyuvm.test(expect_fail=True)
class AoiTestErrors(AoiTest):
    """Stellt sicher, dass der Fehler erkannt wird"""
    def start_of_simulation_phase(self):
        ConfigDB().set(None, "*", "CREATE_ERRORS", True)
\end{lstlisting}

\section{Bus Functional Model: \texttt{AoiBfm}}
\label{subsec:bfm}

Der Bus Functional Model (BFM) in \texttt{aoi\_2\_2\_utils.py} repräsentiert
eine alternative, tiefer abstrahierte Schicht der DUT-Kommunikation. Während
der Treiber aus Abschnitt~\ref{sssec:driver} direkt \texttt{cocotb.top}-Signale
beschreibt, kapselt das BFM diese Pegel-Zugänge in asynchrone Warteschlangen
und dedizierte Coroutinen:

\begin{lstlisting}[style=pythonstyle, caption={BFM-Struktur mit Schreib- und Ãœberwachungs-Koroutinen}]
class AoiBfm(metaclass=Singleton):
    def __init__(self):
        self.dut = cocotb.top
        self.input_queue    = Queue(maxsize=1)
        self.cmd_mon_queue  = Queue(maxsize=0)
        self.result_mon_queue = Queue(maxsize=0)

    async def driver_bfm(self):
        """Entnimmt Items aus input_queue, legt SWT an"""
        self.dut.SWT.value = 0
        while True:
            try:
                a,b,c,d = self.input_queue.get_nowait()
                self.dut.SWT.value = (d<<3)|(c<<2)|(b<<1)|a
                await Timer(2, units="ns")
            except QueueEmpty:
                await Timer(1, units="ns")

    async def cmd_mon_bfm(self):
        """Beobachtet SWT-Aenderungen, leitet Tupel weiter"""
        prev_swt = -1
        while True:
            await Timer(1, units="ns")
            cur = int(self.dut.SWT.value)
            if cur != prev_swt:
                a=cur&1; b=(cur>>1)&1; c=(cur>>2)&1; d=(cur>>3)&1
                self.cmd_mon_queue.put_nowait((a,b,c,d))
                prev_swt = cur

    async def result_mon_bfm(self):
        """Dekodiert SEG-Ausgang nach jeder SWT-Aenderung"""
        ...

    def start_bfm(self):
        cocotb.start_soon(self.driver_bfm())
        cocotb.start_soon(self.cmd_mon_bfm())
        cocotb.start_soon(self.result_mon_bfm())
\end{lstlisting}

Das BFM verwendet das \emph{Singleton}-Muster
(\texttt{metaclass=Singleton}): Es existiert zur Laufzeit genau eine
\texttt{AoiBfm}-Instanz pro Simulation, auf die alle Komponenten zugreifen.
Dies verhindert Races zwischen konkurrierenden Schreibzugriffen auf
\texttt{dut.SWT}.

Die drei Coroutinen laufen \emph{concurrent} via \texttt{cocotb.start\_soon()},
was dem Prinzip der parallelen Hardware-Signalverarbeitung entspricht: Treiber-,
Kommando-Monitor- und Ergebnis-Monitor-Logik sind voneinander entkoppelt.

\section{Prädiktionsmodell und Fehlerinjektion}
\label{subsec:praediktion}

Die Funktion \texttt{aoi\_prediction()} in \texttt{aoi\_2\_2\_utils.py} ist
das Python-Referenzmodell (Golden Model) für den AOI-2-2-Schaltkreis:

\begin{lstlisting}[style=pythonstyle, caption={Golden Model: \texttt{aoi\_prediction()}}]
def aoi_prediction(a, b, c, d, error=False):
    """
    Python-Modell der AOI-2-2-Funktion.
    Y = NOT((A AND B) OR (C AND D))
    """
    result = int(not ((a & b) | (c & d)))
    if error:
        result = 1 - result   # Bit-Invertierung fuer Fehlerinjektion
    return result
\end{lstlisting}

Diese einzeilige Implementation bildet Gleichung~\ref{eq:aoi22} unmittelbar
ab. Das optionale \texttt{error=True}-Flag invertiert das Ergebnis, was dem
Scoreboard ermöglicht, absichtlich falsche Vorhersagen zu machen und damit
das Reaktionsverhalten des Testinfrastruktur zu validieren.

\begin{theorem}[Korrektheit des Golden Models]
\label{thm:golden_model}
Für alle $(a,b,c,d) \in \{0,1\}^4$ gilt:
\begin{align}
	\mathrm{aoi\_prediction}(a,b,c,d,\mathrm{False}) = f_{\mathrm{AOI}}(a,b,c,d)
\end{align}
\end{theorem}

\begin{proof}
Vollständige Aufzählung: Tabelle~\ref{tab:aoi_wahrheit} zeigt alle 16
Eingabebelegungen. Die Python-Auswertung von
\texttt{int(not ((a \& b) | (c \& d)))} ist für jeden Wert direkt
über\-prüfbar und stimmt mit der Wahrheitstabelle überein.
\end{proof}

\section{Test-Strategien und Verifikationsabdeckung}
\label{subsec:test_strategien}

\begin{definition}[Vollständige Eingangsüberdeckung]
Ein Testlauf erreicht \emph{vollständige Eingangsüberdeckung} für den
AOI-2-2-Schaltkreis, wenn alle $2^4 = 16$ Eingabebelegungen mindestens
einmal simuliert wurden, d.\,h.\
$C_{\mathrm{fcov}} = 1$.
\end{definition}

Die vier Sequenzen decken unterschiedliche Teilbereiche des Stimulus-Raums ab:

\begin{table}[H]
\centering
\begin{tabular}{|l|r|c|p{4.5cm}|}
\hline
\textbf{Sequenz} & \textbf{Items} & $C_{\mathrm{fcov}}$ & \textbf{Stärke} \\
\hline
\texttt{ExhaustiveSeq} & 16 & $=1$ (garantiert) & Vollständige systematische Abdeckung \\
\texttt{RandomSeq} & 20 & $<1$ (erwartet) & Realitätsnaher Stimulus; statistisch $\approx\!91\%$ \\
\texttt{CornerCaseSeq} & 8 & $0{,}5$ & Gezielte Eckfälle; Design-Intent-Test \\
\texttt{TestAllSeq} & $\geq\!36$ & $=1$ & Kombiniert exhaustiv + zufällig \\
\hline
\end{tabular}
\caption{Gegenüberstellung der Sequenzen nach Stimulus-Anzahl und Überdeckungsgrad.}
\label{tab:sequenz_vergleich}
\end{table}

\paragraph{Statistischer Überdeckungsnachweis für \texttt{RandomSeq}.}
\label{par:random_coverage}
Die Wahrscheinlichkeit, dass nach $n$ zufälligen, gleichverteilten Items
alle $k = 16$ Belegungen mindestens einmal aufgetreten sind, entspricht
dem Coupon-Collector-Problem:

\begin{equation}
    P(\text{alle } k \text{ nach } n \text{ Items}) = \frac{k!}{k^n} \cdot S(n, k),
\end{equation}

wobei $S(n,k)$ die Stirling-Zahl zweiter Art ist. Für $k = 16$ und $n = 20$ ist
diese Wahrscheinlichkeit mit ca.\ $14{,}3\,\%$ relativ gering — weshalb
\texttt{RandomSeq} allein \emph{nicht} als hinreichende Verifikationsmethode
gilt, sondern stets in Kombination mit \texttt{ExhaustiveSeq} (via
\texttt{TestAllSeq}) eingesetzt wird.

\section{Make-basierter Simulations-Workflow}
\label{subsec:make_workflow}

Das \texttt{Makefile} implementiert einen dualen Simulations-Workflow, der beide
HDL-Varianten mit je einem spezialisierten Open-Source-Simulator unterstützt:

\begin{table}[H]
\centering
\begin{tabular}{|l|l|l|}
\hline
\textbf{HDL} & \textbf{Simulator} & \textbf{Make-Befehl} \\
\hline
Verilog & Verilator & \texttt{make coverage SIM=verilator} \\
VHDL    & GHDL      & \texttt{make sim TOPLEVEL\_LANG=vhdl SIM=ghdl} \\
\hline
\end{tabular}
\caption{Dualer Simulations-Workflow: Verilog mit Verilator, VHDL mit GHDL.}
\label{tab:sim_workflow}
\end{table}

\subsection{Verilog-Simulation mit Verilator}

Verilator kompiliert Verilog-Code in C++-Code und erzeugt hochperformante
Simulations-Executables. Im vorliegenden Makefile werden bei Auswahl
\texttt{SIM=verilator} zusätzlich Zeilend-, Toggle- und Branch-Überdeckung
aktiviert:

\begin{lstlisting}[style=bashstyle, caption={Verilator-spezifische Makefile-Konfiguration}]
ifeq ($(SIM), verilator)
    # Unterstuetzung fuer #delay-Anweisungen
    COMPILE_ARGS += --timing
    # Zeilenueberdeckung, Toggle- und Branch-Coverage
    COMPILE_ARGS += --coverage --coverage-line --coverage-toggle
    EXTRA_ARGS   += --coverage
endif
\end{lstlisting}

Nach der Simulation erzeugt Verilator eine \texttt{coverage.dat}-Datei.
Diese wird durch \texttt{verilator\_coverage} in eine annotierte
HTML-Ausgabe transformiert:

\begin{lstlisting}[style=bashstyle, caption={Ãœberdeckungsbericht mit Verilator}]
# Simulation mit Ueberdeckungs-Messung
make clean
make coverage SIM=verilator

# Annotierter HTML-Bericht
make coverage-report

# Universeller Python-Bericht
python3 generate_universal_coverage.py --sim verilator
\end{lstlisting}

\subsection{VHDL-Simulation mit GHDL}

GHDL ist ein quelloffener VHDL-Simulator mit IEEE-1076-2008-Unterstützung.
Er wird über den \texttt{SIM=ghdl}-Parameter aufgerufen:

\begin{lstlisting}[style=bashstyle, caption={VHDL-Simulation mit GHDL}]
# Simulation der VHDL-Implementierung
make clean
make sim TOPLEVEL_LANG=vhdl SIM=ghdl

# Universeller Ueberdeckungsbericht
python3 generate_universal_coverage.py --sim ghdl
\end{lstlisting}

\begin{bemerkung}
GHDL-Coverage-Unterstützung erfordert einen Build mit GCC- oder LLVM-Backend
(\texttt{-{}-coverage}). Im Standard-\texttt{mcode}-Backend steht
Überdeckungsmessung nicht zur Verfügung. Das Makefile enthält die
entsprechenden GHDL-Optionen kommentiert, um die Aktivierung bei verfügbarem
Backend zu erleichtern.
\end{bemerkung}

\subsection{Quelldatei-Auflösung im Makefile}

Das Makefile unterstützt zwei Suchpfade für die HDL-Quelldateien — den
Standard-Unterordner \texttt{hdl/verilog/} bzw.\ \texttt{hdl/vhdl/} sowie
ein optionales lokales Verzeichnis — und löst diese zur Compile-Zeit auf:

\begin{lstlisting}[style=bashstyle, caption={Robuste Quelldatei-Suche im Makefile}]
ifeq ($(TOPLEVEL_LANG),verilog)
  ifneq ("$(wildcard $(CWD)/hdl/verilog/aoi_2_2.v)","")
    VERILOG_SOURCES = $(CWD)/hdl/verilog/aoi_2_2.v
  else ifneq ("$(wildcard $(CWD)/aoi_2_2.v)","")
    VERILOG_SOURCES = $(CWD)/aoi_2_2.v
  else
    $(error "Datei aoi_2_2.v nicht gefunden")
  endif
endif
\end{lstlisting}

\section{Ãœberdeckungsbericht}
\label{subsec:coverage_report}

Das Skript \texttt{generate\_universal\_coverage.py} erzeugt einen
einheitlichen HTML-Ãœber\-deckungsbericht, der sowohl Verilator- als auch
GHDL-Simulationen unterstützt. Es implementiert eine abstrakte Parser-Hierarchie:

\begin{lstlisting}[style=pythonstyle, caption={Abstrakte Parser-Hierarchie für universelle Coverage-Berichte}]
class CoverageParser:        # abstrakte Basisklasse
    def parse(self): raise NotImplementedError

class VerilatorCoverageParser(CoverageParser):
    """Parst annotierte Verilator-Coverage-Dateien"""
    def parse(self):
        # Interpretiert Praefixe: ~ = partiell,
        # %000000 = nicht abgedeckt, sonst: Zaehler-Wert
        ...

class GHDLCoverageParser(CoverageParser):
    """Analysiert GHDL-Quelldatei heuristisch"""
    def parse(self):
        # Heuristik: Zeilen mit Zuweisungen/case/if
        # als abgedeckt markieren (vereinfacht)
        ...
\end{lstlisting}

Aus den Parsed-Daten wird ein HTML-Dokument mit:
\begin{itemize}
    \item Farbkodierter Zeilenüberdeckung: grün (abgedeckt), rot (nicht abgedeckt),
          gelb (partiell)
    \item Statistik-Header: Prozentsatz abgedeckter Zeilen, Simulator-Badge
    \item Simulator-spezifischer Fußzeile mit Erklärungen zur Überdeckungsmethodik
\end{itemize}

\begin{lstlisting}[style=bashstyle, caption={Aufruf des universellen Ãœberdeckungsberichts}]
# Fuer Verilator (nach 'make coverage-report'):
python3 generate_universal_coverage.py --sim verilator

# Fuer GHDL (nach 'make sim TOPLEVEL_LANG=vhdl SIM=ghdl'):
python3 generate_universal_coverage.py --sim ghdl
\end{lstlisting}

\section{Laufzeit-Profiling}
\label{subsec:profiling}

Das Skript \texttt{profile\_simulation.sh} kombiniert die Cocotb-Simulation
mit dem Python-Profiler \texttt{cProfile} und erzeugt eine
\texttt{.pstat}-Datei, die anschließend mit \texttt{analyze\_profile\newline .py}
ausgewertet wird:

\begin{lstlisting}[style=bashstyle, caption={Simulations-Profiling mit cProfile}]
#!/bin/bash
make clean
python3 -m cProfile -o simulation_profile.pstat \
    -m cocotb.runner \
    --toplevel aoi_2_2 \
    --module testbench
\end{lstlisting}

\begin{lstlisting}[style=pythonstyle, caption={\texttt{analyze\_profile.py}: Auswertung der Profil-Daten}]
def analyze_profile(profile_file="profile.pstat"):
    stats = pstats.Stats(profile_file)

    print("TOP 20 FUNKTIONEN NACH KUMULATIVER ZEIT:")
    stats.sort_stats("cumulative")
    stats.print_stats(20)

    print("TOP 20 NACH EIGENZEIT:")
    stats.sort_stats("time")
    stats.print_stats(20)

    print("TOP 20 NACH AUFRUFANZAHL:")
    stats.sort_stats("calls")
    stats.print_stats(20)
\end{lstlisting}

Das Profiling dient in erster Linie der Identifikation von Laufzeit-Engpässen
in der Python-Testbench, nicht im DUT selbst. Typische Hot Spots in cocotb-
basierten Testbenches sind: \texttt{Timer}-Triggers, Queue-Operationen und
die UVM-interne Phasen-Dispatch-Logik. Diese Informationen sind wertvoll für
die Optimierung großer Testsuiten mit Tausenden von Items.

\section{Simulations- vs.\ formal-basierte Verifikation}
\label{subsec:vergleich_sim_formal}

Die simulations-basierte und die formal-basierte Verifikation sind \emph{komplementär}:
Jede Methode hat spezifische Stärken und Schwächen, die die jeweils andere ergänzt.

\begin{table}[H]
\centering
\small
\begin{tabular}{|l|p{3.8cm}|p{3.8cm}|}
\hline
\textbf{Eigenschaft} & \textbf{Formal (Yosys/SAT)} & \textbf{Simulation (pyUVM)} \\
\hline
Vollständigkeit & Vollständig für alle Eingaben & Stichprobenartig \\
Überprüfte Eigenschaft & Äquivalenz zweier Designs & Korrektheit gegen Referenzmodell \\
Überdeckungsmetrik & n/a (vollständig per Konstruktion) & $C_{\mathrm{fcov}} \in [0,1]$ \\
Fehlertypen & Strukturelle und funktionale Fehler & Funktionale Fehler, Timing \\
Skalierung & Exponentiell (SAT, Worst-Case) & Linear in Testvektorzahl \\
Laufzeit für AOI-2-2 & $< 1\,$s & $< 5\,$s (alle Tests) \\
Fehlerinjektion & Durch Mutationsdesigns & Durch \texttt{CREATE\_ERRORS}-Flag \\
CI/CD-Integration & Hoch (deterministisch) & Hoch (determinierbar durch Seed) \\
Lernkurve & Steil (Yosys-Syntax) & Moderat (Python/UVM) \\
\hline
\end{tabular}
\caption{Systematischer Vergleich formaler und simulations-basierter Verifikation für den AOI-2-2-Schaltkreis.}
\label{tab:sim_formal_vergleich}
\end{table}

\paragraph{Optimale Kombination.}
Die in dieser Arbeit empfohlene Verifikationsstrategie lautet:

\begin{enumerate}
    \item \textbf{Simulations-basierte Erstverifikation}: pyUVM-Testbench
          deckt funktionale Fehler früh im Entwicklungsprozess auf; die
          Laufzeit ist gering und die Feedback-Schleife kurz.
    \item \textbf{Formale Äquivalenzprüfung}: Yosys beweist abschließend die
          Äquivalenz zwischen Verilog- und VHDL-Fassung für alle Eingaben,
          ohne Stichprobencharakter.
    \item \textbf{Ãœberdeckungsnachweis}: Die \texttt{Coverage}-Komponente stellt
          sicher, dass die Simulation alle Verhaltensklassen erfasst hat;
          Lücken werden als Test-Fehler behandelt.
\end{enumerate}

\begin{theorem}[Komplementarität von Simulation und formaler Verifikation]
\label{thm:komplementaritaet}
Eine simula\-tions-basierte Verifikation mit $C_{\mathrm{fcov}} = 1$
\emph{und} eine formale Äquivalenzprüfung (UNSAT) zusammen implizieren:
Das DUT-Design ist korrekt gegenüber dem Referenzmodell für alle Eingaben,
und das Referenzmodell wurde für alle möglichen Eingabeklassen evaluiert.
\end{theorem}

\begin{proof}
Die formale Äquivalenzprüfung (UNSAT) zeigt: Für alle $\vec{x}$, $\mathrm{DUT}(\vec{x}) = \mathrm{Ref}(\vec{x})$.
Die Simulation mit $C_{\mathrm{fcov}} = 1$ stellt sicher, dass alle 16
Eingabebelegungen mindestens einmal durch das Scoreboard geprüft wurden und
das Golden Model in jedem Fall korrekte Ausgaben lieferte
(Satz~\ref{thm:golden_model}). Beide Bedingungen zusammen garantieren
vollständige Korrektheit des DUT gegen eine validierte Spezifikation.
\end{proof}

\section{Zusammenfassung}
\label{subsec:pyuvm_zusammenfassung}

In diesem Kapitel wurde eine vollständige simulations-basierte Testbench für
den AOI-2-2-Schaltkreis entwickelt und auf dessen Verilog- und
VHDL-Implementierung angewendet. Die wesentlichen Ergebnisse sind:

\begin{enumerate}
    \item \textbf{Vollständige UVM-Architektur}: Sequenz-Items, vier Sequenzklassen,
          Treiber, Scoreboard, Coverage-Komponente und Umgebungsklasse wurden
          nach dem UVM-1800.2-Standard in pyUVM implementiert.
    \item \textbf{Duale HDL-Unterstützung}: Dieselbe Testbench simuliert sowohl
          die Verilog-Fassung (Verilator) als auch die VHDL-Fassung (GHDL),
          ohne Änderungen am Testbench-Code.
    \item \textbf{Ãœberdeckungsgarantie}: Die \texttt{ExhaustiveSeq} erreicht
          stets $C_{\mathrm{fcov}} = 1$ (Satz~\ref{thm:coverage_exhaustive});
          die \texttt{Coverage}-Komponente erzwingt diesen Nachweis als
          Test-Fehlerbedingung.
    \item \textbf{Golden Model und Fehlerinjektion}: Die
          \texttt{aoi\_prediction()}-Funktion stellt ein mathematisch korrektes
          Referenzmodell bereit (Satz~\ref{thm:golden_model}); das
          Fehlerinjektionsflag validiert die Reaktionsfähigkeit der
          Testinfrastruktur selbst.
    \item \textbf{Universelles Ãœberdeckungs-Reporting}: Das Skript
          \texttt{generate\_universal\_cover\newline age.py} erzeugt simulator-unabhängige
          HTML-Berichte für Verilator und GHDL.
    \item \textbf{Komplementarität}: Simulations-basierte und formal-basierte
          Verifikation ergänzen sich zur vollständigen Verifikationskette
          (Satz~\ref{thm:komplementaritaet}).
\end{enumerate}

\chapter{Schluss}

Diese Arbeit hat gezeigt:

\begin{enumerate}
    \item Amaranth kann VHDL-Designs funktional äquivalent reimplementieren.
    \item Die Äquivalenz lässt sich mit Open-Source-Tools (Yosys, SAT-Solver) formal beweisen.
    \item Der Workflow von Python-Code zu Hardware-Validierung ist praktikabel und vollstän\-dig reproduzierbar.
    \item Amaranth bietet Produktivitätsvorteile bei erhaltener Korrektheit: ca.~200 Zeilen Python gegenüber ca.~250 Zeilen VHDL, mit identischem Syntheseergebnis.
    \item Der Subgraph Algorithmus \cite{subgraph2026} stellt eine theoretisch fundierte,
          polynomial-zeitliche Methode zur strukturellen Äquivalenzprüfung von
          Hardware-Designs auf Graph-Ebene bereit (Kapitel~\ref{sec:subgraph_equiv}).
    \item Acht neu hergeleitete Sätze mit vollständigen Beweisen (Kapitel~\ref{sec:subgraph_equiv}) belegen
          Korrektheit, Laufzeit $O(n^3)$, Sensitivität gegenüber Einzelkanten-Mutationen,
          Transitivität, super-exponentiellen Speedup gegenüber naiver Permutationssuche
          sowie SETH-bedingte Optimalität des Verfahrens.
    \item Die SAT-basierte und die Subgraph\-Algorithmus-basierte Me\-tho\-de ergänzen sich
          komplementär: erstere prüft funktionale Äquivalenz auch nach Logikoptimierung,
          letztere prüft strukturelle Äquivalenz mit polynomieller Laufzeitgarantie.
    \item Kapitel~\ref{sec:verilog_vhdl_direkt} führt einen vollständigen, dreistufigen Äquivalenzbeweis
          zwischen einer Verilog- und einer VHDL-Implementierung der CITEC-ALU-Schaltung:
          (a)~manuelle Inspektion aller sechs Operationen und der sequentiellen Logik,
          (b)~strukturelle Prüfung mittels Subgraph Algorithmus auf dem
          $11 \times 11$-Abhängigkeitsgraphen sowie
          (c)~SAT-basierte Bestätigung durch Yosys.  Das Hauptresultat
          (Satz~\ref{thm:hauptsatz_vv}) beweist vollständige Ausgabeäquivalenz
          für alle Eingabesequenzen und identische Anfangszustände.
    \item Kapitel~\ref{sec:pyuvm_aoi} demonstriert eine vollständige simulations-basierte
          Testbench auf Basis von pyUVM \cite{pyuvm2024} und cocotb \cite{cocotb2024} für
          den AOI-2-2-Schaltkreis. Dieselbe Testbench simuliert sowohl die Verilog-Fassung
          (Verilator mit Zeilenüberdeckung) als auch die VHDL-Fassung (GHDL), erzwingt
          vollständige funktionale Eingangsüberdeckung ($C_{\mathrm{fcov}} = 1$) und
          validiert das Verifikationssystem selbst durch Fehlinjektion. Die Methode ergänzt
          die formale Verifikation um dynamische Ãœberdeckungs\-nachweise und Fehlerinjektionstests
          in einer einheitlichen, Python-nativen Verifikationskette.
\end{enumerate}

\section{Andere Python-basierte HDLs}

\paragraph{MyHDL}
MyHDL \cite{myhdl2024} ist ein älteres Python-basiertes HDL, das ähnliche Ziele wie Amaranth verfolgt. Im Vergleich bietet Amaranth:
\begin{itemize}
	\item Modernere Python-Syntax und -Features
	\item Bessere Integration mit Yosys
	\item Aktivere Community und Entwicklung
\end{itemize}

\paragraph{Migen}
Amaranth ist der Nachfolger von Migen. Wichtige Verbesserungen:
\begin{itemize}
	\item Klarere Trennung von kombinatorischer und synchroner Logik
	\item Verbesserte Fehlerbehandlung
	\item Standardisierte Build-System-Integration
\end{itemize}

\paragraph{Open-Source-Alternativen}
Yosys bietet vergleichbare Funktionalität für Open-Source-Workflows. Limitationen bestehen bei sehr großen Designs ($> 10^6$ Gates).

\section{Zukünftige Arbeiten}

Die vorliegende Arbeit hat eine grundlegende Methodik zur formalen Verifikation von
Amaranth-HDL-Designs etabliert und deren Praktikabilität anhand einer ALU mit
7-Segment-Ausgabe demonstriert. Sie legt damit ein Fundament, auf dem eine Vielzahl von
weiterführenden Forschungs- und Entwicklungsarbeiten aufgebaut werden kann. Die
folgenden Abschnitte skizzieren diese Perspektiven im Detail.

% -----------------------------------------------------------------------
\subsection{Verifikation komplexerer Hardware-Designs}

\paragraph{Vollständige Prozessorkerne}
Der natürliche nächste Schritt ist die Übertragung der Methodik auf vollständige
Prozessorkerne. RISC-V hat sich als de-facto-Standard für Open-Source-Prozessoren
etabliert; prominente Open-Source-Implementierungen wie PicoRV32 oder VexRiscv
existieren sowohl in Verilog/VHDL als auch in Teilen bereits als Amaranth-Portierungen.
Ein umfassender Äquivalenznachweis für einen 32-Bit-RISC-V-Kern würde die Skalierbarkeit
der vorgestellten Methodik direkt unter Beweis stellen. Dabei müssen insbesondere
mehrtaktige Operationen, Pipeline-Stufen und bedingte Sprünge korrekt modelliert werden,
was den Einsatz tiefer induktiver Beweisführung (etwa \texttt{equiv\_induct} mit
$k > 20$ Schritten) notwendig macht. Die Äquivalenzprüfung eines Prozessorkerns
stellt wegen der inhärenten Zustandsraumexplosion eine methodische Herausforderung dar,
die angepasste Abstraktions- und Decompositions-Strategien erfordert.

\paragraph{Bus-Protokolle und On-Chip-Interconnects}
Moderne SoC-Designs bestehen zu einem erheblichen Teil aus Bus-Protokoll-Logik.
Standardinterfaces wie AXI4, AXI4-Lite, Wishbone oder APB weisen präzise,
formal spezifizierbare Protokollgarantien auf (korrekte Handshake-Sequenzen,
keine Deadlocks, kein Datenverlust). Zukünftige Arbeiten könnten Amaranth-Implementierungen
dieser Protokolle formal gegen Referenzimplementierungen oder gegen formale Protokoll-
spezifikationen (z.~B. in PSL oder SVA) verifizieren. Dies wäre insbesondere für
Safety-kritische Anwendungen wertvoll, in denen Fehler in der Interconnect-Logik
schwerwiegende Systemausfälle verursachen können.

\paragraph{Digitale Signalverarbeitungs-Pipelines}
DSP-Pipelines (FIR-Filter, FFT, Korrelationsberechnungen) gelten als besonders
fehleranfällig, da numerische Eigenschaften wie Overflow-Verhalten, Rundungsfehler und
Sättigungsarithmetik schwer zu überblicken sind. Amaranth eignet sich durch seine
Python-Integration gut für die Beschreibung von DSP-Strukturen, da Filterkoeffizienten
und Topologien direkt aus NumPy-Berechnungen abgeleitet werden können. Die formale
Verifikation dieser Designs würde den Wertebereich aller Zwischen- und Ausgangssignale
sowie das Sättigungsverhalten beweisbar korrekt machen. Dies erfordert jedoch die
Einbindung von Bit-precise-Arithmetik in den SAT-Beweis, was den Einsatz von
SMT-Solvern (z.~B. Z3 oder Bitwuzla) nahelegt.

\paragraph{Speichercontroller und DRAM-Interfaces}
Die Ansteuerung von DRAM-Bausteinen (DDR3/DDR4/LPDDR) erfordert zeitlich präzise
Signalfolgen mit engen Timing-Fenstern. Eine Amaranth-basierte Implementierung eines
DDR-PHY-Controllers, verifiziert gegen eine JEDEC-konforme Referenz, wäre ein
bedeutender industrieller Beitrag. Gleichzeitig illustrierte ein solches Projekt die
Grenzen der rein logischen Äquivalenzprüfung: Timing-Properties können nicht allein durch
strukturelle Äquivalenz nachgewiesen werden und erfordern ergänzend statische
Timing-Analyse (STA) oder Timed-Model-Checking.

% -----------------------------------------------------------------------
\subsection{Erweiterung der formalen Verifikationstechniken}

\paragraph{Systemische Property-Verification mit SVA und PSL}
Neben der reinen Äquivalenz\-prüfung ist die property-basierte Verifikation ein
eigenständiges und mächtiges Paradigma. SystemVerilog Assertions (SVA) und die
Property Specification Language (PSL) erlauben es, Sicherheits- (\textit{safety}) und
Lebendigkeitseigenschaften (\textit{liveness}) direkt am Design zu annotieren. Beispiele
für relevante Properties der vorgestellten ALU wären:

\begin{itemize}
  \item \textbf{Overflow-Freiheit}: $\forall a, b \leq 15 : \text{ADD}(a, b) \leq 30$
        — keine impliziten Überläufe bei 4-Bit-Operanden.
  \item \textbf{Division-by-Zero-Sicherheit}: $b = 0 \Rightarrow \text{DIV}(a, b) = \texttt{0xFF}$
        — der Fehlerwert wird in jedem Cycle korrekt ausgegeben.
  \item \textbf{Display-Vollständigkeit}: Jede der vier Displaystellen wird in jedem
        Zyklus von $4 \times T_{\mathrm{refresh}}$ genau einmal aktiviert.
\end{itemize}

Zukünftige Arbeiten könnten PSL-Annotationen direkt in den Amaranth-Code integrieren
(ähnlich wie \texttt{assert} in simulationsbasierten Frameworks) und deren automatische
Ãœbersetzung in Yosys-Formal-Checks realisieren.

\paragraph{SMT-Solver-Integration}
Yosys unterstützt eine SMT-basierte Back-End-Verifikation über das Tool \texttt{sby}
(SymbiYosys). Während SAT-Solver für kombinatorische Äquivalenz\-prüfung gut geeignet sind,
erlauben SMT-Solver (Satisfiability Modulo Theories) die Einbeziehung arithmetischer
Theorien (Linear Integer Arithmetic, Bit-Vector-Logik). Dies ist besonders relevant für
DSP-Designs oder für den Nachweis arithmetischer Identitäten. Eine Direkt-Integration
des SymbiYosys-Workflows in die Amaranth-Build-Pipeline, analog zur beschriebenen
Yosys-Pipeline, wäre ein wertvoller Folgebeitrag.

\paragraph{Formale Verifikation mit temporaler Logik}
Model Checking mithilfe von temporaler Logik (LTL, CTL) erlaubt den Nachweis von
Eigenschaften über zeitliche Verläufe, die über Äquivalenzprüfung hinausgehen. So könnte
für das Display-Multiplexing formal gezeigt werden, dass keine Stelle dauerhaft
ausgeblendet bleibt (\textit{Liveness}) und keine zwei Stellen gleichzeitig aktiv sind
(\textit{Mutual Exclusion}). Open-Source-Tools wie NuSMV oder DIVINE könnten hier als
Ergänzung zu Yosys eingesetzt werden.

% -----------------------------------------------------------------------
\subsection{Automatisierung und CI/CD-Integration}

\paragraph{Kontinuierliche Verifikation in GitHub Actions}
Moderne Hardware-Projekte profitieren von automatisierten Qualitätssicherungspipelines
analog zur Praxis in der Softwareentwicklung. Die vorgestellte Verifikationsmethodik
eignet sich ideal für die Integration in CI/CD-Workflows: Bei jedem Commit könnten
automatisch die folgenden Stufen durchlaufen werden:

\begin{enumerate}
  \item Syntaxvalidierung des Amaranth-Python-Codes via \texttt{pylint}/\texttt{mypy}
  \item Verilog-Generierung via Amaranth-Backend
  \item Äquivalenzprüfung mit Yosys gegen die versionierte VHDL-Referenz
  \item Ressourcen-Synthese mit Yosys und automatische Erkennung von
        Ressourcen-Regres\-sionen ($>5\%$ LUT-Zunahme)
  \item Optionale FPGA-Synthese in Vivado (z.~B. täglich, nicht bei jedem Commit)
\end{enumerate}

GitHub Actions mit vorbereiteten Docker-Containern (z.~B.
\texttt{hdl/containers}) können Yosys, openFPGALoader und Python-Abhängigkeiten
reproduzierbar bereitstellen, ohne dass CI-Runner lokal installiert werden müssen.

\paragraph{Regressionstests und Testabdeckungsmetriken}
Analog zur Linienabdeckung in der Softwareverifikation existieren für formale Hardware-
Verifikation Metriken wie \textit{Mutation Coverage} (werden Fehlerinjektionen erkannt?)
und \textit{Toggle Coverage} (wechseln alle Signale in der Simulation ihren Wert?).
Zukünftige Arbeiten könnten automatisierte Mutationstests für Amaranth-Designs entwickeln:
Dabei werden systematisch synthetische Fehler (Bit-Flip in einer Konstante, vertauschte
Operanden) in das Design injiziert, und verifiziert, dass der Äquivalenz-Checker diese
erkennt. Dies qualifiziert gleichzeitig die Verifikationsstärke des Workflows selbst.

\paragraph{Versionierung von Verifikationsbeweisen}
Eine oft vernachlässigte Praxis ist die Versionierung nicht nur des Designs, sondern
auch der Verifikationsergebnisse. Das Log einer erfolgreich abgeschlossenen Yosys-
Äquivalenzprüfung sollte zusammen mit dem Design-Commit committet werden, sodass
nachvollziehbar ist, \textit{welche Version} zu welchem Zeitpunkt formal verifiziert
wurde. Hierfür kann ein einfaches Metadaten-Format (JSON mit Zeitstempel, Commit-Hash
und Yosys-Version) automatisch am Ende des CI-Skripts generiert und archiviert werden.

% -----------------------------------------------------------------------
\subsection{Vollständig quelloffene Toolchain}

\paragraph{F4PGA / Symbiflow für Xilinx-Targets}
Die in dieser Arbeit beschriebene Toolchain erfordert für den Synthese- und
Implementierungsschritt noch das proprietäre Vivado. Das Projekt F4PGA (ehemals
Symbiflow) \cite{f4pga2024} verfolgt das Ziel, für Xilinx Artix-7 und andere FPGAs
eine vollständig quelloffene Implementierungskette bereitzustellen. Zum
Entstehungszeitpunkt dieser Arbeit ist die Artix-7-Unterstützung in F4PGA noch nicht
produktionsreif; insbesondere fehlen offizielle Timing-Datenbanken. Sobald diese Lücken
geschlossen sind, ließe sich der gesamte Workflow von der Amaranth-Beschreibung bis
zur bitgenauen Programmierung des FPGA ohne proprietäre Software realisieren. Dies wäre
ein bedeutender Schritt für Lehre, Reproduzierbarkeit und Open-Science-Prinzipien.

\paragraph{GHDL als VHDL-Frontend für Yosys}
Mit GHDL \cite{ghdl2024} steht ein quelloffener VHDL-Simulator und -Synthesizer zur
Verfügung, der über ein Yosys-Plugin (\texttt{ghdl-yosys-plugin}) direkt VHDL-Designs
in Yosys laden kann — ohne den Umweg über eine KI-gestützte VHDL-zu-Verilog-Konvertierung.
Zukünftige Versionen des vorgestellten Workflows sollten GHDL als Standard-Frontend
verwenden: Das VHDL-Original wird direkt geladen, die manuelle Konvertierungsstufe
entfällt, und potenzielle Semantikdifferenzen durch den Übersetzungsschritt werden
ausgeschlossen. GHDL unterstützt VHDL-2008 und deckt damit auch fortgeschrittene
Sprachfeatures ab.

\paragraph{nextpnr für Lattice- und Intel-FPGAs}
Der Open-Source-Place-and-Route-Tool \texttt{nextpnr} unterstützt bereits vollständig die
Lattice-iCE40- und ECP5-Familien sowie (im Beta-Stadium) Intel Cyclone-10. In Kombination
mit \texttt{Project IceStorm} (iCE40) oder \texttt{Project Trellis} (ECP5) ergibt sich
eine stabile, vollständig quelloffene Implementierungskette. Die Übertragung der
vorgestellten Arbeiten auf kostengünstige Entwicklungsboards wie das iCEbreaker
(Lattice iCE40UP5k) hätte den Vorteil, dass keine proprietäre Software
erforderlich ist, und wäre damit besonders für Lehrveranstaltungen geeignet.

% -----------------------------------------------------------------------
\subsection{Erweiterte Testmethoden}

\paragraph{Property-based Testing mit Hypothesis}
Der beschriebene formale Äquivalenzbeweis ist vollständig und gilt für alle möglichen
Eingabekombinationen. Ergänzend hierzu kann \textit{Property-based Testing} auf
Simulationsebene wertvolle zusätzliche Erkenntnisse liefern, insbesondere in frühen
Designphasen, bevor die formale Verifikation aufgesetzt ist. Das Python-Framework
\texttt{Hypothesis} \cite{hypothesis2024} generiert automatisch Eingabesequenzen, die
Randfälle und Extremwerte gezielt abdecken, und kann mit dem Amaranth-Simulator
kombiniert werden:

\begin{lstlisting}[style=pythonstyle, caption={Hypothesis-basierter ALU-Test}]
from hypothesis import given, settings
from hypothesis import strategies as st
from amaranth.sim import Simulator

@given(
    a=st.integers(min_value=0, max_value=15),
    b=st.integers(min_value=1, max_value=15),
)
@settings(max_examples=1000)
def test_division_property(a, b):
    """Quotient darf nie groesser als Dividend sein"""
    result = simulate_alu(a, b, OP_DIV)
    assert result <= a
\end{lstlisting}

Property-based Tests decken Implementierungsfehler auf, die im formalen SAT-Modell
möglicherweise durch fehlerhafte Modellierung nicht sichtbar werden, und dienen somit
als orthogonale Verifikationsebene.

\paragraph{Simulation mit Cocotb und pyUVM}
\texttt{Cocotb} \cite{cocotb2024} ist ein Python-basiertes Testbench-Framework für
Hardware-Simulationen. Da Amaranth selbst Python-basiert ist, entsteht eine natürliche
Synergie: Testbenches, Stimulusgeneratoren und Assertion-Überprüfungen können in
derselben Python-Codebasis wie das Design selbst geschrieben werden. Die in
Kapitel~\ref{sec:pyuvm_aoi} vorgestellte pyUVM-Testbench für den AOI-2-2-Schaltkreis
legt dafür eine konkrete, vollständig reproduzierbare Basis: die duale
HDL-Unterstützung (Verilog/Verilator und VHDL/GHDL), die erzwungene
Ãœberdeckungsmetrik ($C_{\mathrm{fcov}} = 1$) und die Fehlerinjektions-Tests
können unmittelbar auf komplexere Designs übertragen werden.
Zukünftige Arbeiten könnten diese Architektur auf die ALU aus Kapitel~3 ausdehnen
und eine integrierte Verifikationspipeline entwickeln, die pyUVM-Simulation,
Yosys-Äquivalenzprüfung und Überdeckungsnachweis in einem einzigen CI-Schritt
verbindet.

\paragraph{Erweiterung der pyUVM-Testbench auf sequentielle Designs}
Der in Kapitel~\ref{sec:pyuvm_aoi} vorgestellte pyUVM-Ansatz ist explizit für einen
kombinator\-ischen 4-Eingangs-Schaltkreis konzipiert. Für sequentielle Designs —
etwa die ALU mit Takt-getaktetem Display-Multiplexer — sind folgende Erweiterungen
notwendig und unmittelbar umsetzbar: (a)~Ersetzen des Timer-basierten Einschwingens
durch \texttt{RisingEdge}-Trigger; (b)~Erweiterung der Coverage-Klasse auf
zustandsbehaftete Ãœberdeckungsmodelle (z.\,B.\ Cross-Coverage zwischen Operationscode
und Operanden); (c)~Einführung eines Referenz-Zustandsautomaten im Golden Model für
den Display-Multiplexer. Diese Erweiterungen würden pyUVM als vollwertige
Verifikationsplattform für alle in dieser Arbeit behandelten Designs etablieren.

\paragraph{Hardware-in-the-Loop-Tests}
Für zeitkritische Designs ist die simulationsbasierte Verifikation allein nicht
ausreichend. Hardware-in-the-Loop (HiL)-Tests betreiben das synthetisierte FPGA-Design
unter realen Bedingungen und lesen Ergebnisse über ein UART- oder JTAG-Interface
automatisiert aus. Ein Python-Testrahmen, der die Werte des Amaranth-Simulators mit
den auf dem FPGA gemessenen Ausgaben vergleicht, bildet eine vollständige, dreigliedrige
Verifikationskette: \textit{formale Äquivalenz} (Yosys), \textit{funktionale Simulation}
(Cocotb/Amaranth-Sim) und \textit{physikalische Validierung} (HiL auf Basys 3).

% -----------------------------------------------------------------------
\subsection{Hardware-Software-Co-Design}

\paragraph{Integration von Amaranth mit Embedded-Software-Stacks}
Ein zunehmendes Anwendungsfeld für FPGAs ist das Hardware-Software-Co-Design, bei dem ein
Soft-Core-Prozessor (z.~B. PicoRV32) zusammen mit applikationsspezifischen Hardware-
Acceleratoren auf dem FPGA instanziiert wird. Amaranth eignet sich hervorragend, um
solche Acceleratoren (z.~B. einen AES-Encoder, einen CRC-Calculator oder eine
Integer-Multiplikationseinheit) auf hohem Abstraktionsniveau zu beschreiben. Zukünftige
Arbeiten könnten zeigen, wie die formally-verified Amaranth-Hardware-Komponenten nahtlos
in einen SoC-Stack integriert werden, bei dem die Software-API und die Hardware-Logik
aus derselben Python-Codebasis abgeleitet werden.

\paragraph{Automatische Interface-Generierung}
Die Port-Definitionen von Amaranth-Modulen liegen als Python-Objekte vor. Daraus lässt
sich automatisch C-Header-Code generieren, der die Registeradressen und Bitefeldmasken
des Hardware-Interfaces für die Embedded-Software definiert. Dieses Prinzip (bekannt aus
Frameworks wie LiteX \cite{litex2024}) reduziert den Synchronisationsaufwand zwischen
Hardware- und Softwareentwicklung erheblich und verhindert Inkonsistenzen bei
Interface-Änderungen.

% -----------------------------------------------------------------------
\subsection{Skalierbarkeit und Grenzen der formalen Verifikation}

\paragraph{Zustandsraumexplosion bei tiefer Sequentialität}
Die induktive Äquivalenzprüfung mit \texttt{equiv\_induct} skaliert gut für
Designs mit geringer sequentieller Tiefe. Bei tiefen Pipelines (z.~B. einem
5-stufigen RISC-V-Pipeline mit Forwarding-Logik) kann die Menge der zu explorierenden
Zustandsübergänge die Kapazität des SAT-Solvers übersteigen. 
Kontermaßnahmen sind:
\begin{itemize}
    \item \textbf{Compositionelle Verifikation}: Das Design wird in unabhängige
          Teilkomponenten zerlegt, die separat verifiziert werden, und die Korrektheit
          der Komposition wird durch Annahme-Garantie-Kontrakte (Assume-Guarantee)
          formal abgesichert.
    \item \textbf{Abstraktions-Verfeinerung (CEGAR)}: Counterexample-Guided Abstraction
          Refinement beginnt mit einem abstrakten, vereinfachten Modell und verfeinert
          es iterativ, bis der Beweis gelingt oder ein echtes Gegenbeispiel gefunden ist.
    \item \textbf{Bounded Model Checking}: Statt eines vollständigen Induktionsbeweises
          werden nur die ersten $k$ Taktzyklen geprüft, was bei bekannten Designzyklen
          ausreicht.
\end{itemize}
Zukünftige Arbeiten sollten diese Techniken systematisch für wachsende Amaranth-Designs
evaluieren und Schwellenwerte (\textit{Complexity Horizons}) empirisch bestimmen.

\paragraph{Grenzen der Yosys-Äquivalenzprüfung bei komplexer Arithmetik}
Die vorliegende Arbeit hat gezeigt, dass Yosys für eine 4-Bit-ALU korrekte Beweise
liefert. Für Designs mit 64-Bit-Integer-Arithmetik, Floating-Point-Einheiten oder
parametrisierbaren Multiplizierer-Bäumen können dedizierte arithmetische SMT-Solver
(Bitwuzla, CVC5) bessere Beweisleistung erzielen. Eine vergleichende Studie zwischen
der Yosys-SAT-Methodik und SMT-basierten Ansätzen für arithmetikintensive Designs wäre
wissenschaftlich wertvoll.

% -----------------------------------------------------------------------
\subsection{Einsatz in der Hochschullehre}

\paragraph{Amaranth als Lehrmittel in Digitalelektronik-Kursen}
Die Python-Syntax von Amaranth senkt die Einstiegshürde für Studierende gegenüber
VHDL erheblich. Gleichzeitig erzwingt die strikte Trennung von kombinatorischer und
synchroner Logik (\texttt{module.d.comb} vs.\ \texttt{module.d.sync}) ein klares
Verständnis dieser fundamentalen Konzepte. Zukünftige Arbeiten könnten einen
didaktisch aufbereiteten Lehrpfad entwickeln, der:
\begin{enumerate}
    \item Mit einfachen kombinatorischen Schaltungen beginnt (Multiplexer, Decoder),
    \item Sequentielle Designmuster einführt (Zähler, Zustandsautomaten),
    \item Zur formalen Verifikation mit Yosys führt und
    \item Mit einer vollständigen FPGA-Implementierung abschließt.
\end{enumerate}

\paragraph{Reproduzierbarkeit von Forschungsergebnissen}
Die Reproducibility-Crisis in der Informatikforschung betrifft zunehmend auch
Hardware-Design-Arbeiten. Der in dieser Arbeit vorgestellte Workflow bietet durch den
Einsatz ausschließlich quelloffener Tools (Amaranth, Yosys, openFPGALoader) und
versionierter Abhängigkeiten eine hohe Reproduzierbarkeit. Zukünftige Arbeiten sollten
Docker-Container bereitstellen, die alle Abhängigkeiten in definierten Versionen
festlegen, sodass die Verifikationsergebnisse in Jahren noch reproduzierbar sind.

% -----------------------------------------------------------------------
\subsection{Industrielle Adoption und Standardisierung}

\paragraph{Integration in ASIC-Workflows}
Während diese Arbeit auf FPGA-basierte Designs fokussiert, ist die Methodik prinzipiell
auf ASIC-Designs übertragbar. Yosys unterstützt über kommerzielle Plugins
(Yosys-SMTBMC, Yosys-EQY) auch den ASIC-Kontext. Amaranth-Designs könnten, nach
Verifikation der Äquivalenz, in Standard-Cell-Syntheseflows (z.~B. mit OpenROAD) 
überführt werden. Dies wäre insbesondere für sicherheitskritische Anwendungen
(Automotive, Aerospace) relevant, in denen formale Verifikation gesetzlich
vorgeschrieben ist (DO-254, ISO 26262).

\paragraph{Beitrag zu offenen HDL-Standards}
Amaranth ist noch kein IEEE-standardisiertes Format. Eine Standardisierungsinitiative
ähnlich SystemVerilog (IEEE 1800) oder VHDL (IEEE 1076) würde die Toolchain-Interoperabilität
erhöhen und die industrielle Akzeptanz fördern. Die Forschungsgemeinschaft könnte durch
Publikationen vergleichender Studien und durch aktive Mitarbeit an der Amaranth-
Spezifikation diesen Prozess unterstützen.

\paragraph{Benchmarking gegen kommerzielle Äquivalenzprüfungs-Tools}
Synopsys Formality und Cadence Conformal sind Industrie-Standards mit hoher
Beweisleistung. Eine systematische Vergleichsstudie, die Yosys gegen diese kommerziellen
Tools auf standardisierten Benchmark-Suites (z.~B. ITC-99, ISCAS-89) antritt, würde
die Stärken und Schwächen des quelloffenen Ansatzes quantifizieren und eine fundierte
Empfehlung für verschiedene Design-Szenarien ermöglichen.

% -----------------------------------------------------------------------
\subsection{Subgraph Algorithmus-basierten Äquivalenzprüfung}

Das in Kapitel~\ref{sec:subgraph_equiv} entwickelte Verfahren eröffnet eine Vielzahl
weiterführender Forschungsrichtungen. Der folgende Abschnitt skizziert diese
systematisch und im Detail.

\paragraph{Erweiterung auf sequentielle Hardware-Designs}

Die in Kapitel~\ref{sec:subgraph_equiv} bewiesene Implikation (Satz~\ref{thm:struktur_impliziert_funktion})
gilt explizit für \emph{kombinatorische} Hardware-Designs, da nur bei diesen die
Datenabhängigkeit vollständig durch einen azyklischen Graphen beschrieben werden kann.
Sequentielle Designs — also solche mit Flip-Flops, Registern und Rückkopplungen — erzeugen
Zyklen im Abhängigkeitsgraphen, die eine direkte Anwendung des Subgraph Algorithmus
erschweren.

Ein zentraler Ansatz zur Erweiterung ist die \emph{Zeitentfaltung} (Time Unrolling):
Das sequentielle Design wird über $k$ Taktzyklen in ein äquivalentes kombinatorisches
Design transformiert, indem der Zustandsraum explizit in die Knotenstruktur eingebettet
wird. Der resultierende Entfaltungsgraph $\hat{H}(D, k)$ besitzt $k \cdot n$ Knoten (mit
$n$ Signalen im Original) und ist azyklisch, sofern keine kombinatorischen Schleifen
existieren. Der Subgraph Algorithmus kann dann auf $\hat{H}(D_1, k)$ und $\hat{H}(D_2, k)$
angewendet werden.

Die Laufzeit des erweiterten Verfahrens beträgt dann $O((kn)^3) = O(k^3 n^3)$, da die
Knotenanzahl sich $k$-fach erhöht. Für kleine Entfaltungstiefe $k$ — typisch $k \leq 10$
bei bekannter Pipeline-Tiefe — bleibt das Verfahren praktikabel. Zukünftige Arbeiten
sollten systematisch untersuchen, ab welcher Entfaltungstiefe $k$ die Äquivalenz
\emph{vollständig} erfasst wird, und diesen Wert für relevante Klassen sequentieller
Designs (Zähler, Zustandsautomaten, Pipeline-Register) empirisch bestimmen.

Darüber hinaus bietet die Erweiterung des Subgraph Algorithmus auf gewichtete Graphen
eine alternative Modellierungsmöglichkeit: Kanten erhalten Gewichte, die die zeitliche
Verzögerungs- oder Taktbindung repräsentieren. Eine solche Gewichtserweiterung würde es
erlauben, sowohl kombinatorische als auch getaktete Pfade im selben Graphmodell
darzustellen, ohne explizite Entfaltung durchführen zu müssen.

\paragraph{Hybride Verifikation: Kombination von Subgraph Algorithmus und SAT-Solver}

Das in Kapitel~\ref{subsec:vergleich} gezogene Fazit, dass SAT-basierte und
Subgraph Algorithmus-basierte Methoden komplementär sind, legt eine hybride
Verifikationspipeline nahe, die beide Stärken vereint. Das konkrete Vorgehen könnte
wie folgt ausstuiert werden:

\begin{enumerate}
    \item \textbf{Vorfilterung mit dem Subgraph Algorithmus (in $O(n^3)$):}
          Der Subgraph Algorithmus prüft zunächst, ob die Hardware-Abhängigkeitsgraphen
          strukturell äquivalent sind. Falls ja, ist die Äquivalenz nachgewiesen und der
          SAT-Solver muss nicht bemüht werden. Dies spart bei strukturell gleichen
          Designs den gesamten SAT-Aufruf ein.
    \item \textbf{SAT-basierte Tiefenprüfung bei Strukturdivergenz:}
          Falls der Subgraph Algorithmus Strukturunterschiede feststellt (d.\,h.\ die
          Designs sind nicht strukturell äquivalent), bedeutet das noch keine funktionale
          Nicht-Äquivalenz: Logikoptimierungen im Syntheseprozess können dieselbe Funktion
          mit anderem Graphen realisieren. In diesem Fall wird der SAT-Solver aktiviert,
          um die funktionale Äquivalenz zu prüfen. Der Subgraph Algorithmus hat dann
          wenigstens eine Schnellaussage über die Strukturdifferenz geliefert und
          möglicherweise den Suchraum für den SAT-Solver durch Identifikation
          identisch-strukturierter Teilgraphen (als fixe Punkte) eingeschränkt.
    \item \textbf{Rückkopplung: Gegenbeispiele des SAT-Solvers als Graphannotation:}
          Wenn der SAT-Solver ein Gegenbeispiel liefert (eine Eingabekombination, auf der
          beide Designs divergieren), kann diese Information genutzt werden, um den
          Abhängigkeitsgra\-phen zu annotieren: der divergente Pfad wird markiert,
          und der Subgraph Algorithmus kann in zukünftigen Versionen bevorzugt auf
          annotierten Teilgraphen angewendet werden.
\end{enumerate}

Eine solche hybride Pipeline würde insbesondere bei CI/CD-Anwendungen einen erheblichen
Laufzeitvorteil bieten: Die überwältigende Mehrheit der Check-ins ändert nichts an der
Abhängigkeitsstruktur (nur Kommentare, Parameternamen, Formatierung), sodass der
Subgraph Algorithmus allein suffizienten Nachweis der Äquivalenz in polynomieller Zeit
liefert, ohne den SAT-Solver zu starten.

\paragraph{Automatische Extraktion von Hardware-Abhängigkeitsgraphen aus Netzlisten}

Der in Kapitel~\ref{subsec:graph_repr} beschriebene Hardware-Abhängigkeitsgraph wurde
konzeptionell eingeführt. Die automatische Extraktion dieses Graphen aus bestehenden
Werkzeugketten ist ein eigenständiges Forschungsproblem. Konkret bieten sich zwei
Ansätze an:

\textbf{Extraktion aus Amaranth-Zwischenrepräsentation:} Das Amaranth-Framework
erzeugt intern eine RTLIL-Repräsentation (RTL Intermediate Language), die von Yosys
verarbeitet wird. Aus dieser Repräsentation lässt sich der Abhängigkeitsgraph direkt
ableiten: Jede RTLIL-Zelle (Gatter, Multiplexer, arithmetische Einheit) wird zum
Knoten, jede Verbindung zwischen Zellen zur Kante. Ein Python-Plugin für Yosys, das
diesen Graphen als Adjazenzliste oder -matrix exportiert, würde die nahtlose Integration
in die bestehende Verifikationspipeline erlauben.

\textbf{Extraktion aus Verilog-AST mit Pyverilog:} Für den VHDL-Zweig (über GHDL oder
Claude-konvertiertes Verilog) kann das Python-Framework \texttt{Pyverilog} den
abstrakten Syntaxbaum (AST) des Verilog-Codes analysieren und direkt in einen
Abhängigkeitsgraphen überführen. Dies hätte den Vorteil, unabhängig von der
RTLIL-Ebene zu arbeiten und damit auch höherstufige Abstraktionsphänomene
(Schleifenstrukturen, parameterisierte Blöcke) zu erfassen.

In beiden Fällen ist ein \emph{Graph-Normalisierungsschritt} notwendig: Knoten werden
in topologischer Reihenfolge nummeriert (Schritt~1 der Normalisierung), und redundante
Geistknoten (buffers, wires ohne Logik) werden kollabiert (Schritt~2). Erst nach dieser
Normalisierung sind die Graphen aus Amaranth und VHDL vergleichbar und können in den
Subgraph Algorithmus eingespeist werden. Die Entwicklung eines robusten
Normalisierungsalgorithmus, der mit den Quirks beider Werkzeugketten umgeht, ist
ein abgeschlossenes, praxisrelevantes Forschungsthema.

\paragraph{Theoretische Erweiterungen: Bedingte und unbedingte untere Schranken}

Satz~\ref{thm:optimalitaet_aequivalenz} (Kapitel~\ref{sec:subgraph_equiv}) beweist die
Optimalität des Verfahrens \emph{unter SETH}. SETH ist eine weithin akzeptierte, aber
unbedingt unbewiesene Komplexitätshypothese. Zukünftige theoretische Arbeiten könnten
in zwei Richtungen voranschreiten:

\textbf{Richtung 1 — Stärkere bedingte Schranke unter schwächerer Hypothese:}
SETH impliziert die \emph{Orthogonal Vectors Conjecture} (OVC), die wiederum untere
Schranken für LCS und Edit-Distanz impliziert. Möglicherweise lässt sich die
$\Omega(n^{3-\varepsilon})$-Schranke für das Subgraph-Äquivalenzprüfungsproblem
bereits aus der schwächeren NSETH (Non-uniform SETH) oder aus der
3SUM-Conjecture ableiten, was die Schranke auf eine breitere, robustere Basis stellen würde.

\textbf{Richtung 2 — Unbedingte $\Omega(n^{2+\delta})$-Schranke:}
Ein unbedingter Beweis einer superkubischen oder auch nur superquadratischen unteren
Schranke für das zyklische Subgraph-Problem würde — ohne jede Hypothese — die praktische
Relevanz des Algorithmus beziffern. Ein aussichtsreicher Ansatz wäre, die
Kommunikationskomplexität des Problems zu analysieren: In einem Zwei-Parteien-Protokoll,
bei dem Alice Matrix $A$ und Bob Matrix $B$ hält und beide entscheiden müssen, ob
$G \cong_s G'$, kann die Kommunikationskomplexität als untere Schranke für die
sequentielle Zeitkomplexität herangezogen werden.

\paragraph{Skalierung auf sehr große Hardware-Designs}

Der Subgraph Algorithmus hat eine Laufzeit von $\Theta(n^3)$, was für $n > 10\,000$
Knoten praktisch problematisch werden kann: Bei $n = 10\,000$ wären theoretisch
$10^{12}$ Operationen nötig. Für industrielle Designs mit Millionen von Signalen
sind approximative oder hierarchische Erweiterungen unumgänglich.

\textbf{Hierarchische Dekomposition:} Analog zur compositionellen Verifikation (siehe
oben) könnte das Hardware-Design in hierarchische Teilgraphen zerlegt werden.
Jede Hierarchieebene wird separat mit dem Subgraph Algorithmus verglichen; die
Korrektheit der Komposition wird durch Assume-Guarantee-Kontrakte abgesichert. Die
Gesamtlaufzeit reduziert sich von $O(n^3)$ auf $O(k \cdot (n/k)^3) = O(n^3 / k^2)$,
wenn das Design in $k$ gleichgroße Teilkomponenten zerlegt werden kann — ein quadratischer
Speedup durch Dekomposition.

\textbf{Approximative Subgraph-Prüfung mit Fehlertoleranz:} Für Designs, bei denen
geringfügige Implementierungsunterschiede (z.\,B.\ verschiedene Synthese-Optimierungen)
toleriert werden können, könnte ein approximativer Modus des Subgraph Algorithmus
eingeführt werden: Statt exakter Signaturübereinstimmung wird eine Hamming-Distanz
$d(\hat{\sigma}^A, \hat{\sigma}^B) \leq \tau$ toleriert. Dies erlaubt den Nachweis
von „annähernder Äquivalenz'' (Approximate Equivalence), die in der
Hardware-Validierung für Optimierungspässe nützlich sein kann.

\textbf{Parallelisierung der Rotationssuche:} Satz~\ref{thm:laufzeit_aequivalenz} zeigt,
dass die $n$ Rotationsvergleiche unabhängig voneinander sind (Lemma~4 in \cite{subgraph2026}).
Dies ist ideal für parallele Ausführung: Auf einem Mehrkern-Prozessor oder einer GPU
können alle $n$ LCS-Berechnungen gleichzeitig durchgeführt werden, was die
Laufzeit von $\Theta(n^3)$ auf $\Theta(n^2 \cdot \lceil n / p \rceil)$ mit $p$ Kernen
reduziert — bei $p = n$ Kernen also auf $\Theta(n^2)$, was der theoretischen unteren
Schranke der Eingabeleseschritte entspricht (Lemma~\ref{lem:injektivitaet}).

\paragraph{Erweiterung der Graphrepräsentation für heterogene Designaspekte}

Die aktuelle Definition des Hardware-Abhängigkeitsgraphen (Definition~\ref{def:hdg})
erfasst ausschließlich kombinatorische Datenabhängigkeiten. Für eine umfassendere
Äquivalenzprüfung sind folgende Erweiterungen relevant:

\textbf{Getypte Kanten (Typed Edges):} Verschiedene Arten von Abhängigkeiten — Datenpfad,
Kontrollpfad, Taktdomäne — werden durch verschiedene Kantentypen unterschieden. Der
Subgraph Algorithmus müsste entsprechend auf Multigraphen erweitert werden, wobei
die Signatur-Funktion für jeden Kantentyp eine separate Bit-Ebene verwendet:
\[
\sigma_j^{(\text{typ})}(M) = \sum_{i=0}^{n-1} M_{ij}^{(\text{typ})} \cdot 2^{i + n \cdot \text{typ}} + j \cdot 2^{n \cdot T},
\]
mit $T$ verschiedenen Kantentypen. Die Injektivität bleibt durch die disjunkte
Bit-Ebenen-Codierung erhalten.

\textbf{Knotengewichte (Node Weights):} Knoten erhalten Gewichte, die den Operationstyp
kodieren (Addition, Vergleich, Multiplexer usw.). Der Subgraph Algorithmus prüft dann
nicht nur die Verbindungsstruktur, sondern auch, ob entsprechend markierte Operationen
übereinstimmen. Dies erlaubt die Unterscheidung zwischen zwei Designs, die strukturell
äquivalent sind, aber verschiedene Operationstypen auf korrespondierenden Knoten
ausführen.

\textbf{Gerichtete vs.\ ungerichtete Subgraph-Prüfung:} Abhängigkeitsgraphen sind
naturgemäß gerichtet (Daten fließen von Eingabe zu Ausgabe). Der vorliegende Subgraph
Algorithmus operiert auf der Spaltenstruktur der Adjazenzmatrix (eingehende Kanten je
Knoten). Zukünftige Arbeiten könnten eine symmetrische Variante entwickeln, die sowohl
eingehende als auch ausgehende Kanten berücksichtigt und damit bidirektionale
Datenfluss-Abhängigkeiten (z.\,B.\ in Rückkopplungsschleifen) korrekt modelliert.

\paragraph{Visualisierung und interaktive Beweisexploration}

Die theoretischen Beweise in Kapitel~\ref{sec:subgraph_equiv} sind formal rigoros,
aber für Ingenieure ohne Informatik-Theorieausbildung schwer zugänglich. Zukünftige
Arbeiten könnten Werkzeuge entwickeln, die den Subgraph-Äquivalenzprüfungsprozess
interaktiv visualisieren:

\begin{itemize}
    \item \textbf{Graph-Diff-Visualisierung:} Darstellung der Hardware-Abhängigkeitsgraphen
          beider Designs nebeneinander, mit farblicher Hervorhebung übereinstimmender
          (grün), divergenter (rot) und redundanter (gelb) Knoten und Kanten.
          Rotationsverschiebungen werden als animierte Übergänge dargestellt.
    \item \textbf{LCS-Trace:} Anzeige der optimalen LCS-Lösung als hervorgehobene
          Teilsequenz innerhalb der Signatur-Arrays, sodass unmittelbar sichtbar wird,
          welche Signale in beiden Designs übereinstimmen.
    \item \textbf{Interaktiver Rotationsexplorer:} Schieberegler, mit dem der Ingenieur
          manuell die $n$ Rotationen durchläuft und für jede Rotation den aktuellen
          LCS-Wert und die Graphübereinstimmung verfolgt. Dies fördert das intuitive
          Verständnis des Algorithmus und erleichtert die Fehlerdiagnose bei
          Nicht-Äquivalenz.
\end{itemize}

Ein solches Werkzeug könnte als Plugin für bestehende Hardware-Design-Umgebungen
(z.\,B.\ als Yosys-Webinterface oder als Jupyter-Notebook-Widget) implementiert werden.

\paragraph{Anwendung des Subgraph Algorithmus auf höhere Abstraktionsebenen}

Bisher wurde der Subgraph Algorithmus auf der Ebene der RTL-Signals angewendet.
Die Methode lässt sich jedoch auf deutlich höhere Abstraktionsebenen übertragen:

\textbf{Funktionsblock-Ebene (IP-Core-Vergleich):} In modernen SoC-Designs werden
häufig vorgefertigte IP-Cores (z.\,B.\ AES-Beschleuniger, UART-Controller) aus
verschiedenen Quellen integriert. Der Subgraph Algorithmus könnte auf Ebene der
Blockdiagramme eingesetzt werden: Jeder IP-Core ist ein Knoten, jede Daten- oder
Steuerverbindung zwischen Cores ist eine Kante. Die strukturelle Äquivalenz zweier
SoC-Designs auf Block-Ebene lässt sich dann in $O(k^3)$ Zeit prüfen, mit $k$ der
(deutlich kleineren) Anzahl von IP-Cores. Dies ist besonders wertvoll für
Sicherheitsaudits, bei denen nachgewiesen werden soll, dass ein modifiziertes SoC-Design
keine unerwünschten neuen Verbindungen enthält.

\textbf{Algorithm-Level-Synthese (HLS):} High-Level-Synthesis-Werkzeuge (z.\,B.\ Vitis HLS,
Bambu) übersetzen C/C++-Algorithmen in RTL-Designs. Zwei verschiedene
HLS-Konfigurationen desselben Algorithmus (z.\,B.\ verschiedene Unrolling-Faktoren,
verschiedene Ressourcenbeschränkungen) erzeugen strukturell unterschiedliche, aber
funktional äquivalente RTL-Designs. Der Subgraph Algorithmus auf dem HLS-Datenflussgraphen
(DFG) — der naturgemäß dieselbe Struktur wie der Hardware-Abhängigkeitsgraph hat —
kann hier eingesetzt werden, um vor dem eigentlichen RTL-Vergleich bereits auf
algorithmischer Ebene Äquivalenz oder Divergenz festzustellen.

\textbf{Software-Äquivalenz via AST-Graphen:} Die ursprüngliche Motivation des Subgraph
Algorithmus \cite{subgraph2026} sind abstrakte Syntaxbäume (ASTs) von Programmen.
Ein unmittelbarer Übertrag der Hardware-Äquivalenzprüfungs-Methodik auf
Software-Äquivalenz liegt daher nahe: Zwei Implementierungen desselben Algorithmus
(z.\,B.\ in Python und in C) werden in ihre AST-Graphen überführt und mit dem Subgraph
Algorithmus verglichen. Dies würde die bestehende Anwendungsdomäne erheblich
erweitern und den Subgraph Algorithmus zu einem universellen, sprach- und
plattformübergreifenden Äquivalenzprüfungswerkzeug machen.

\paragraph{Formale Verankerung im Curry-Howard-Isomorphismus}

Die in Kapitel~\ref{sec:subgraph_equiv} geführten Beweise sind klassische
mathematische Beweise. Ein fortgeschrittenes, theoretisch bedeutsames Ziel wäre die
Formalisierung dieser Beweise in einem interaktiven Theorembeweiser wie Isabelle/HOL,
Lean~4 oder Coq. Dies hätte mehrere Vorteile:

\begin{itemize}
    \item \textbf{Maschinengeprüfte Korrektheit:} Alle acht Sätze aus
          Kapitel~\ref{sec:subgraph_equiv} — insbesondere Satz~\ref{thm:struktur_impliziert_funktion}
          (Implikation strukturell $\Rightarrow$ funktional) und
          Satz~\ref{thm:optimalitaet_aequivalenz} (SETH-Optimalität) — würden durch
          den Theorembeweiser maschinell verifiziert, ohne die Möglichkeit menschlicher
          Lücken im Induktionsbeweis.
    \item \textbf{Extraktion von verifizierten Algorithmen:} Aus dem Lean-4-Beweis
          des Korrektheitssatzes könnte mittels Curry-Howard direkt verifizierbarer
          Haskell- oder OCaml-Code extrahiert werden, der den Subgraph Algorithmus
          korrekt implementiert — eine in der formalen Methodik als \textit{verified
          extraction} bekannte Technik.
    \item \textbf{Wiederverwendbare Beweisbausteine:} Die Lemmas über zyklische Gruppen
          (Satz~\ref{thm:transitivitaet}), Signatur-Injektivität (Lemma~\ref{lem:injektivitaet})
          und LCS-Komplexität unter SETH (Lemma aus \cite{subgraph2026}) können als
          wiederverwendbare Proof-Bausteine in eine gemeinsame Bibliothek einfließen,
          die auch für verwandte Algorithmen genutzt werden kann.
\end{itemize}

\paragraph{Integration in sicherheitskritische Zertifizierungsrahmen}

Die Beweise in Kapitel~\ref{sec:subgraph_equiv} legen eine Grundlage für den Einsatz
des Subgraph Algorithmus in sicherheitskritischen Domänen, in denen formale
Verifikation gesetzlich vorgeschrieben ist. Konkret bieten sich folgende Anwendungsfelder an:

\textbf{DO-254 (Avionics Hardware):} Der Standard DO-254 (Design Assurance Guidance for
Airborne Electronic Hardware) schreibt für Design-Assurance-Level A (höchste Sicherheitsstufe)
explizite Äquivalenzprüfungen zwischen verschiedenen Entwicklungsstufen vor. Der
Subgraph Algorithmus könnte als zertifizierfähiges Werkzeug positioniert werden, sofern
seine Korrektheit formal (maschinell geprüft) nachgewiesen ist und eine
Tool-Qualifikation gemäß DO-330 durchgeführt wurde.

\textbf{ISO~26262 (Automotive):} Für ASIL-D-Sicherheitsfunktionen (höchste
Automobil-Sicherheitsstufe, z.\,B.\ in Bremssystemen oder autonomen Fahrfunktionen)
fordert ISO~26262 den Nachweis funktionaler Äquivalenz zwischen Spezifikation und
Implementation. Die Subgraph Algorithmus-basierte Methode könnte hier als
ergänzender, performanter Vorab-Check eingesetzt werden, der die Notwendigkeit
vollständiger SAT-basierter Verifikation auf strukturell divergente Fälle beschränkt.

\textbf{Common Criteria (IT-Sicherheit):} Im Rahmen der Common Criteria Evaluation
Methodology (CEM) wird für hochstufige Evaluierungen (EAL 6/7) ein formaler
Implementierungsnachweis gefordert. Hardware-Security-Module (HSMs), die kryptographische
Operationen in FPGA-Logik implementieren, könnten durch die Kombination aus
Subgraph Algorithmus-basierter Strukturprüfung und SAT-basierter Funktionalitätsprüfung
effizient zertifiziert werden.

\paragraph{Benchmarking des Subgraph Algorithmus auf Hardware-Äquivalenz-Benches}

Für eine wissenschaftlich belastbare Einordnung des Verfahrens ist ein systematisches
Benchmarking unabdingbar. Folgende Benchmark-Suites sind relevant:

\begin{itemize}
    \item \textbf{ITC-99 und ISCAS-89:} Standardisierte kombinatorische (ISCAS-85) und
          sequentielle (ISCAS-89) Schaltkreis-Benchmarks aus der Testgenerierungsforschung.
          Für jedes Benchmark-Design existieren bekannte Varianten (mit und ohne
          Optimierung), die als Äquivalenzprüfungs-Instanzen genutzt werden können.
    \item \textbf{OpenCores Hardware Designs:} Die Open-Source-Community OpenCores stellt
          hunderte von RTL-Designs bereit (UART, I2C, SPI, DMA-Controller, CPU-Kerne).
          Viele dieser Designs existieren in mehreren Varianten (verschiedene Autoren,
          Syntheseziele). Sie bieten einen realistischen Benchmark für die
          Subgraph Algorithmus-basierte Äquivalenzprüfung.
    \item \textbf{Synthetischer Mutations-Benchmark:} Für eine kontrollierte Evaluation
          des Sensitivitätsverhaltens (Satz~\ref{thm:sensitivitaet}) können Designs
          durch gezielte Einzel-Kanten-Mutationen modifiziert werden.
          Der Subgraph Algorithmus sollte in $100\%$ der Fälle die Mutation als
          Strukturdivergenz erkennen — eine Kennzahl, die direkt aus
          Satz~\ref{thm:sensitivitaet} und Korollar~\ref{kor:sensitivitaet} folgt.
\end{itemize}

Eine Benchmarkstudie würde darüber hinaus empirisch validieren, bei welcher Graphgröße
$n$ der theoretische $O(n^3)$-Vorteil gegenüber SAT-basierten Methoden (mit exponentiellem
Worst-Case) praktisch spürbar wird — ein Ergebnis mit unmittelbarer
Relevanz für die Systemdimensionierung in der Praxis.

\section{Schlusswort}

Amaranth demonstriert das Potential von Python-basierten HDLs als produktive Alternative
zu VHDL und Verilog. Die formale Verifizierbarkeit mit Open-Source-Tools macht Amaranth
zu einer ernstzunehmenden Option für professionelle Hardware-Entwicklung.

Das in Kapitel~\ref{sec:subgraph_equiv} neu entwickelte Verfahren der
Subgraph Algorithmus-basierten Äquivalenzprüfung ergänzt die etablierte SAT-basierte
Methodik um eine strukturelle, polynomial-zeitliche Perspektive. Mit acht rigoros
bewiesenen Sätzen — von der Korrektheit über die Laufzeitoptimalität bis hin zur
Transitivität und Sensitivität — liefert dieses Kapitel einen eigenständigen
theoretischen Beitrag, der über die konkrete ALU-Anwendung hinausgeht und einen
allgemeinen Rahmen für die graphbasierte Hardware-Äquivalenzprüfung bietet.

Die in Kapitel~\ref{sec:pyuvm_aoi} vorgestellte pyUVM-Testbench unterstreicht, dass
die Python-Nativität von Amaranth eine natürliche Erweiterung in die Domäne der
UVM-basierten Simulation erlaubt: Design und Verifikation teilen eine Sprache, eine
Entwicklungsumgebung und ein Ökosystem. Formale Verifikation (Yosys), strukturelle
Äquivalenzprüfung (Subgraph Algorithmus) und simulations-basierte Verifikation (pyUVM)
ergänzen sich zur vollständigen, mehrstufigen Qualitätssicherungskette für
Hardware-Designs in Verilog und VHDL.

Die im Ausblick skizzierten Erweiterungen — Sequentialität, hybride Pipelines,
automatische Graphextraktion, Skalierung auf industrielle Designs, formale Verankerung
in Theorembeweisern sowie Integration in Sicherheitszertifizierungsrahmen — zeigen,
dass das vorgestellte Verfahren ein fruchtbares Fundament für eine neue Forschungsrichtung
an der Schnittstelle von Algorithmentheorie und formaler Hardware-Verifikation bildet.

\begin{thebibliography}{99}

\bibitem{subgraph2026}
  Epp, S.,
  \textit{Subgraph Algorithmus},
  2026, Verfügbar bei \url{https://github.com/hjstephan86/science}.

\bibitem{amaranth2024}
  Amaranth HDL Project,
  \textit{Amaranth HDL Documentation},
  \url{https://amaranth-lang.org/},
  2024.

\bibitem{yosys2024}
  Claire Xenia Wolf,
  \textit{Yosys Open SYnthesis Suite},
  \url{http://www.clifford.at/yosys/},
  2024.

\bibitem{myhdl2024}
  MyHDL Project,
  \textit{MyHDL Documentation},
  \url{http://www.myhdl.org/},
  2024.

\bibitem{ieee1076}
  IEEE Standard,
  \textit{IEEE 1076-2019 - IEEE Standard VHDL Language Reference Manual},
  2019.

\bibitem{openfpgaloader}
  OpenFPGALoader Project,
  \textit{Universal utility for programming FPGA},
  \url{https://github.com/trabucayre/openFPGALoader},
  2024.

\bibitem{basys3}
  Digilent Inc.,
  \textit{Basys 3 FPGA Board Reference Manual},
  2024.

\bibitem{sat_equiv}
  Biere, A., Heule, M., van Maaren, H., Walsh, T.,
  \textit{Handbook of Satisfiability},
  IOS Press, 2009.

\bibitem{formal_verification}
  Clarke, E., Grumberg, O., Peled, D.,
  \textit{Model Checking},
  MIT Press, 1999.

\bibitem{f4pga2024}
  F4PGA Project,
  \textit{F4PGA -- Free and Open-Source FPGA Toolchain},
  \url{https://f4pga.org/},
  2024.

\bibitem{ghdl2024}
  Gingold, T.,
  \textit{GHDL -- Open Source VHDL 87/93/00/02/08 Simulator},
  \url{https://ghdl.github.io/ghdl/},
  2024.

\bibitem{hypothesis2024}
  MacIver, D. R. et al.,
  \textit{Hypothesis: Property-Based Testing for Python},
  \url{https://hypothesis.works/},
  2024.

\bibitem{cocotb2024}
  Cocotb Contributors,
  \textit{cocotb -- Coroutine-based Co-simulation Test Bench},
  \url{https://www.cocotb.org/},
  2024.

\bibitem{litex2024}
  Enjoy-Digital,
  \textit{LiteX -- A Python-Based SoC Builder},
  \url{https://github.com/enjoy-digital/litex},
  2024.

\bibitem{symbiflow2024}
  SymbiYosys Contributors,
  \textit{SymbiYosys -- Formal Hardware Verification Front-End},
  \url{https://symbiyosys.readthedocs.io/},
  2024.

\bibitem{pyuvm2024}
  Ray Salemi,
  \textit{pyUVM -- Python Universal Verification Methodology},
  \url{https://github.com/pyuvm/pyuvm},
  2024.

\bibitem{uvm2020}
  Accellera Systems Initiative,
  \textit{IEEE 1800.2-2020 -- IEEE Standard for Universal Verification Methodology
  Language Reference Manual},
  IEEE, 2020.

%% Weitere Quellen (nicht im Text zitiert) %%

\bibitem{verilog2005}
  IEEE Standard,
  \textit{IEEE 1364-2005 -- IEEE Standard for Verilog Hardware Description Language},
  IEEE, 2006.

\bibitem{systemverilog2017}
  IEEE Standard,
  \textit{IEEE 1800-2017 -- IEEE Standard for SystemVerilog -- Unified Hardware Design,
  Specification, and Verification Language},
  IEEE, 2017.

\bibitem{ovm2008}
  Cadence Design Systems; Mentor Graphics,
  \textit{Open Verification Methodology (OVM) Reference Manual, Version 2.0},
  2008.

\bibitem{vmm2006}
  Bergeron, J., Cerny, E., Hunter, A., Nightingale, A.,
  \textit{Verification Methodology Manual for SystemVerilog},
  Springer, 2006.

\bibitem{sv_uvm_praktikal}
  Vip, M.,
  \textit{SystemVerilog and UVM Verification: A Practical Guide},
  Independently published, 2021.

\bibitem{verilator2023}
  Snyder, W., Verilator Contributors,
  \textit{Verilator -- Fast Open-Source Verilog/SystemVerilog Simulator},
  \url{https://www.veripool.org/verilator/},
  2023.

\bibitem{iverilog2023}
  Williams, S.,
  \textit{Icarus Verilog -- A Free Verilog Simulator},
  \url{http://iverilog.icarus.com/},
  2023.

\bibitem{nvc2024}
  Newsome, N.,
  \textit{NVC -- VHDL Compiler and Simulator},
  \url{https://github.com/nickg/nvc},
  2024.

\bibitem{modelsim2022}
  Siemens EDA,
  \textit{ModelSim SE User's Manual},
  Siemens EDA, 2022.

\bibitem{questa2023}
  Siemens EDA,
  \textit{Questa Sim User's Manual},
  Siemens EDA, 2023.

\bibitem{sby2024}
  Wolf, C. X., SymbiYosys Contributors,
  \textit{SymbiYosys Formal Verification Flow},
  \url{https://symbiyosys.readthedocs.io/},
  2024.

\bibitem{boolector2019}
  Niemetz, A., Preiner, M., Wolf, C. X., Biere, A.,
  \textit{Boolector 3.0},
  Journal on Satisfiability, Boolean Modeling and Computation,
  Vol. 11, pp. 99--102, 2019.

\bibitem{z32024}
  Bjørner, N., de Moura, L., Lal, A.,
  \textit{Z3: A Production Theorem Prover},
  Microsoft Research, \url{https://github.com/Z3Prover/z3}, 2024.

\bibitem{bitwuzla2024}
  Niemetz, A., Preiner, M.,
  \textit{Bitwuzla: A Fast SMT Solver for Fixed-Size Bit-Vectors},
  \url{https://bitwuzla.github.io/},
  2024.

\bibitem{abc2023}
  Brayton, R., Mishchenko, A.,
  \textit{ABC: A System for Sequential Synthesis and Verification},
  UC Berkeley, \url{https://github.com/berkeley-abc/abc}, 2023.

\bibitem{cbmc2023}
  Kroening, D., Tautschnig, M.,
  \textit{CBMC: Bounded Model Checking for ANSI-C and C++ Programs},
  LNCS Vol. 8413, Springer, 2023.

\bibitem{spin2004}
  Holzmann, G. J.,
  \textit{The SPIN Model Checker: Primer and Reference Manual},
  Addison-Wesley, 2004.

\bibitem{nusmv2002}
  Cimatti, A., Clarke, E. M., Giunchiglia, E.,
  \textit{NuSMV 2: An OpenSource Tool for Symbolic Model Checking},
  Proc. CAV 2002, LNCS 2404, pp. 359--364, Springer, 2002.

\bibitem{cadence_jaspergold2023}
  Cadence Design Systems,
  \textit{JasperGold Formal Verification Platform},
  Cadence, 2023.

\bibitem{halbwachs93}
  Halbwachs, N.,
  \textit{Synchronous Programming of Reactive Systems},
  Kluwer Academic Publishers, 1993.

\bibitem{biere_bounded_2003}
  Biere, A., Clarke, E., Raimi, R., Zhu, Y.,
  \textit{Verifying Safety Properties of a PowerPC Microprocessor Using Symbolic Model
  Checking Without BDDs},
  Proc. CAV 1999, LNCS 1633, Springer, 1999.

\bibitem{bdd_bryant1986}
  Bryant, R. E.,
  \textit{Graph-Based Algorithms for Boolean Function Manipulation},
  IEEE Transactions on Computers, Vol. C-35, No. 8, pp. 677--691, 1986.

\bibitem{sat_dpll}
  Davis, M., Logemann, G., Loveland, D.,
  \textit{A Machine Program for Theorem Proving},
  Communications of the ACM, Vol. 5, No. 7, pp. 394--397, 1962.

\bibitem{cdcl_marques1999}
  Marques-Silva, J. P., Sakallah, K. A.,
  \textit{GRASP: A Search Algorithm for Propositional Satisfiability},
  IEEE Transactions on Computers, Vol. 48, No. 5, pp. 506--521, 1999.

\bibitem{minisat2003}
  Eén, N., Sörensson, N.,
  \textit{An Extensible SAT-Solver},
  Proc. SAT 2003, LNCS 2919, pp. 502--518, Springer, 2004.

\bibitem{kissat2020}
  Biere, A., Fleury, M., Heisinger, M.,
  \textit{CaDiCaL, Kissat, Paracooba, Plingeling and Treengeling Entering the SAT
  Competition 2020},
  Proc. SAT Competition 2020, 2020.

\bibitem{qbf_janota2016}
  Janota, M., Klieber, W., Marques-Silva, J., Clarke, E.,
  \textit{Solving QBF with Counterexample Guided Refinement},
  Artificial Intelligence, Vol. 234, pp. 1--25, Elsevier, 2016.

\bibitem{mips_formal2019}
  Wolf, C. X.,
  \textit{Verifying the RISC-V Compliance Suite with SymbiYosys},
  Acks Formal Verification Workshop, 2019.

\bibitem{riscv_isa2019}
  Waterman, A., Asanovic, K. (eds.),
  \textit{The RISC-V Instruction Set Manual, Volume I: Unprivileged ISA},
  RISC-V International, 2019.

\bibitem{xilinx_7series2023}
  AMD Xilinx,
  \textit{7 Series FPGAs Data Sheet: Overview (DS180)},
  AMD, 2023.

\bibitem{xilinx_ug901}
  AMD Xilinx,
  \textit{Vivado Design Suite User Guide: Synthesis (UG901)},
  AMD, 2023.

\bibitem{xilinx_ug904}
  AMD Xilinx,
  \textit{Vivado Design Suite User Guide: Implementation (UG904)},
  AMD, 2023.

\bibitem{intel_quartus2023}
  Intel Corporation,
  \textit{Quartus Prime Pro Edition User Guide: Getting Started},
  Intel, 2023.

\bibitem{lattice_ecp52023}
  Lattice Semiconductor,
  \textit{ECP5 and ECP5-5G Family Data Sheet},
  Lattice Semiconductor, 2023.

\bibitem{nextpnr2023}
  nextpnr Contributors,
  \textit{nextpnr -- A Portable FPGA Place and Route Tool},
  \url{https://github.com/YosysHQ/nextpnr},
  2023.

\bibitem{prjxray2023}
  SymbiFlow Contributors,
  \textit{Project X-Ray: Documenting the Xilinx 7-Series Bitstream Format},
  \url{https://github.com/SymbiFlow/prjxray},
  2023.

\bibitem{vunit2024}
  VUnit Contributors,
  \textit{VUnit -- A Unit Testing Framework for VHDL/SystemVerilog},
  \url{https://vunit.github.io/},
  2024.

\bibitem{uvvm2024}
  Bitvis AS,
  \textit{UVVM -- Universal VHDL Verification Methodology},
  \url{https://www.uvvm.org/},
  2024.

\bibitem{osvvm2024}
  SynthWorks Design,
  \textit{OSVVM -- Open Source VHDL Verification Methodology},
  \url{https://osvvm.org/},
  2024.

\bibitem{scv2005}
  Accellera,
  \textit{SystemC Verification Standard (SCV) 1.0},
  Accellera Systems Initiative, 2005.

\bibitem{ieee_sva}
  IEEE Standard,
  \textit{IEEE 1850-2010: Standard for Property Specification Language (PSL)},
  IEEE, 2010.

\bibitem{python_hdl2023}
  Hussain, B., Khan, F. A.,
  \textit{Python-Based Hardware Description Languages: A Comparative Study},
  ACM Computing Surveys, Vol. 55, No. 8, 2023.

\bibitem{chisel2012}
  Bachrach, J., Vo, H., Richards, B., Lee, Y., Waterman, A.,
  \textit{Chisel: Constructing Hardware in a Scala Embedded Language},
  Proc. DAC 2012, pp. 1215--1224, ACM, 2012.

\bibitem{clash2010}
  Baaij, C., Kooijman, M., Kuper, J., Boeijink, W., Gerards, M.,
  \textit{C\ensuremath{\lambda}aSH: Structural Descriptions of Synchronous Hardware using Haskell},
  Proc. DSD 2010, pp. 714--721, IEEE, 2010.

\bibitem{nMigen2019}
  m-labs,
  \textit{nMigen: A Python-Based Hardware Description Language (predecessor to Amaranth)},
  \url{https://github.com/m-labs/nmigen},
  2019.

\bibitem{migen2013}
  Bourdeauducq, S.,
  \textit{Migen: A Python Toolbox for Building Complex Digital Hardware},
  \url{https://github.com/m-labs/migen},
  2013.

\bibitem{spinalhdl2015}
  Papon, C.,
  \textit{SpinalHDL -- An HDL written in Scala},
  \url{https://github.com/SpinalHDL/SpinalHDL},
  2015.

\bibitem{llhd2020}
  Schuiki, F., Kurth, A., Grosser, T., Benini, L.,
  \textit{LLHD: A Multi-Level Intermediate Representation for Hardware Description Languages},
  Proc. PLDI 2020, pp. 258--271, ACM, 2020.

\bibitem{calyx2021}
  Nigam, R., Thomas, S., Li, Z., Sampson, A.,
  \textit{A Compiler Infrastructure for Accelerator Generators},
  Proc. ASPLOS 2021, ACM, 2021.

\bibitem{circt2021}
  CIRCT Contributors (LLVM Project),
  \textit{CIRCT: Circuit IR Compilers and Tools},
  \url{https://circt.llvm.org/},
  2021.

\bibitem{rtime_formal2017}
  Mukherjee, R., Kroening, D., Melham, T.,
  \textit{Hardware Verification Using Software Analyzers},
  Proc. ISVLSI 2015, pp. 7--12, IEEE, 2015.

\bibitem{coverage_metrics2006}
  Cheng, K.-T., Krishnakumar, A.,
  \textit{Automatic Functional Test Generation Using the Extended Finite State Machine Model},
  Proc. DAC 1993, ACM, 1993.

\bibitem{mutation_testing2017}
  Chowdhury, A., Wait, M., Patnaik, P.,
  \textit{Mutation Testing for Hardware Description Languages},
  IEEE Transactions on VLSI Systems, Vol. 25, No. 3, pp. 987--997, 2017.

\bibitem{constrained_random2011}
  Bergeron, J.,
  \textit{Writing Testbenches Using SystemVerilog},
  Springer, 2011.

\bibitem{functional_coverage2010}
  Piziali, A.,
  \textit{Functional Coverage for Design and Verification},
  Springer, 2010.

\bibitem{ieee_coverage}
  IEEE Standard,
  \textit{IEEE 2148-2020: Standard for Measuring Coverage of Functional Testing},
  IEEE, 2020.

\bibitem{vcd_format}
  IEEE Standard,
  \textit{IEEE 1364-2005 Annex: Value Change Dump Format specification},
  IEEE, 2005.

\bibitem{coq_proof2022}
  Bertot, Y., Castéran, P.,
  \textit{Interactive Theorem Proving and Program Development -- Coq'Art},
  Springer, 2004.

\bibitem{lean4_2021}
  Moura, L. de, Ullrich, S.,
  \textit{The Lean 4 Theorem Prover and Programming Language},
  Proc. CADE 2021, LNCS 12699, Springer, 2021.

\bibitem{alloy2006}
  Jackson, D.,
  \textit{Software Abstractions: Logic, Language, and Analysis},
  MIT Press, 2006.

\bibitem{tla2014}
  Lamport, L.,
  \textit{Specifying Systems: The TLA+ Language and Tools for Hardware and Software Engineers},
  Addison-Wesley, 2002.

\bibitem{temporal_logic_pnueli}
  Pnueli, A.,
  \textit{The Temporal Logic of Programs},
  Proc. 18th Annual Symposium on Foundations of Computer Science,
  pp. 46--57, IEEE, 1977.

\bibitem{ctl_clarke1981}
  Clarke, E. M., Emerson, E. A.,
  \textit{Design and Synthesis of Synchronization Skeletons Using Branching Time Temporal Logic},
  Logics of Programs Workshop, LNCS 131, pp. 52--71, Springer, 1981.

\bibitem{ltl_buchi1960}
  Büchi, J. R.,
  \textit{On a Decision Method in Restricted Second Order Arithmetic},
  Proc. International Congress on Logic, Method, and Philosophy of Science,
  Stanford University Press, pp. 1--11, 1962.

\bibitem{abstract_interpretation1977}
  Cousot, P., Cousot, R.,
  \textit{Abstract Interpretation: A Unified Lattice Model for Static Analysis of Programs
  by Construction or Approximation of Fixpoints},
  Proc. POPL 1977, pp. 238--252, ACM, 1977.

\bibitem{hybrid_verification2019}
  Henzinger, T. A., Jhala, R., Majumdar, R., Sutre, G.,
  \textit{Software Verification with BLAST},
  Proc. SPIN 2003, LNCS 2648, pp. 235--239, Springer, 2003.

\bibitem{counterexample_guided2000}
  Clarke, E. M., Grumberg, O., Jha, S., Lu, Y., Veith, H.,
  \textit{Counterexample-Guided Abstraction Refinement},
  Proc. CAV 2000, LNCS 1855, pp. 154--169, Springer, 2000.

\bibitem{pdr_ic3_2011}
  Bradley, A. R.,
  \textit{SAT-Based Model Checking Without Unrolling},
  Proc. VMCAI 2011, LNCS 6538, pp. 70--87, Springer, 2011.

\bibitem{interpolation_mcmillan2003}
  McMillan, K. L.,
  \textit{Interpolation and SAT-Based Model Checking},
  Proc. CAV 2003, LNCS 2725, pp. 1--13, Springer, 2003.

\bibitem{synthesis_survey2021}
  Nane, R. et al.,
  \textit{A Survey and Evaluation of FPGA High-Level Synthesis Tools},
  IEEE Transactions on CAD of Integrated Circuits and Systems, Vol. 35, No. 10,
  pp. 1591--1604, 2016.

\bibitem{hls_legup2013}
  Canis, A., Choi, J., Aldham, M., Zhang, V.,
  \textit{LegUp: An Open-Source High-Level Synthesis Tool for FPGA-Based Processor/
  Accelerator Systems},
  Proc. FPGA 2011, pp. 33--36, ACM, 2011.

\bibitem{vhls_xilinx2020}
  AMD Xilinx,
  \textit{Vitis High-Level Synthesis User Guide (UG1399)},
  AMD, 2023.

\end{thebibliography}

\end{document}

